% ============================================================
%  Chapter 06: Karl Marx — Historical Materialism, Class Struggle and Alienation
%  Inherits styling from preamble.tex
% ============================================================

\chapter{Karl Marx --- Historical Materialism, Class Struggle and Alienation}

\section{Historical Materialism and the Mode of Production}

\subsection{Why Historical Materialism Comes First}
Before the theory itself, two preliminary points were made in class --- one on how to take notes, and one on why this particular concept has to be mastered before anything else in Marx .

\paragraph{Note-making --- there is no single correct method}
\begin{itemize}
    \item Some students prefer making notes on a computer, because they can edit them later and add material they find relevant while reading newspapers; that is very convenient for them .
    \item Many others, despite knowing how to type, save, open and close files, prefer writing with a pen in a notebook . It is not that they cannot operate a computer; they simply find writing more suitable .
    \item Therefore no one method can be recommended --- whatever suits you, whatever is convenient to you .
    \begin{itemize}
        \item Generally it is preferable to make notes on a device so that you can edit or add to them later; but that does not mean you have to . It is completely according to your convenience .
    \end{itemize}
\end{itemize}

\paragraph{Why this concept is the foundation}
\begin{itemize}
    \item When Karl Marx --- or for that matter any scholar in your syllabus --- wrote a book, he never wrote it for a competitive examination . Scholars were free to write whatever they wanted, in whatever book .
    \item After their death, people arranged those books in chronological order, and then boards decided to lift portions of that material for an examination like this one .
    \item Consequently, elements of historical materialism appear inside the other three concepts as well --- in class struggle (third concept), in mode of production (second concept), and in alienation (fourth concept) .
    \item That is why historical materialism is very important: if you do not understand this, you will not be able to understand the remaining three .
\end{itemize}

\begin{CriticalPoint}
\begin{itemize}
    \item Historical materialism is not one topic among four --- it is the base on which mode of production, class struggle and alienation all rest . It must be finished first .
    \item Marx never wrote for an exam; the exam-friendly ordering of his work is a later, posthumous arrangement .
\end{itemize}
\end{CriticalPoint}

\paragraph{Engels on historical materialism}
\begin{itemize}
    \item Engels once said on historical materialism: just as Darwin developed the laws of organic development of nature, Marx developed the laws of the development of history .
    \item That is what historical materialism is about --- how history developed .
    \item The phrase ``historical materialism'' was not coined by Marx himself; whatever he wrote was in German .
    \item Through this theory Marx presents his own version of how history developed, and this version later came to be regarded as the materialist interpretation of history . Why ``materialist'' becomes clear only at the end of the theory .
\end{itemize}

\begin{QuickRevision}
\begin{itemize}
    \item Engels' formulation: Darwin $\rightarrow$ laws of organic development of nature; Marx $\rightarrow$ laws of development of history .
    \item ``Historical materialism'' is not Marx's own phrase; his writing was in German .
    \item His account of historical development is the materialist interpretation of history .
    \item Master this before mode of production, class struggle and alienation --- all three carry its elements .
\end{itemize}
\end{QuickRevision}

\subsection{The First Historical Act --- Production and Primitive Communism}
Marx begins the story at the very beginning of human society: what did the first human being actually do? 
\begin{enumerate}
    \item The first historical act in human society was production .
    \item Why production? Because survival is important . In order to survive, man has to produce --- food, cloth and shelter .
    \item Food, cloth and shelter are not readily available in nature . Man has to transform nature through labour into material goods .
    \item He could never do all of this by himself . He automatically had to cooperate with multiple people, and therefore it was inevitable to enter into social relations of production .
    \item The primary needs then stand satisfied .
\end{enumerate}
So although nature was transformed into material goods through labour, that labour was not individual labour --- it was collective labour, and entering into social relations of production was unavoidable for man .

\begin{KeyConcept}
\begin{itemize}
    \item First historical act = production, driven by the need to survive (food, cloth, shelter) .
    \item Nature does not supply these ready-made; labour transforms nature into material goods .
    \item Labour is collective, not individual --- hence social relations of production are inevitable, not chosen .
\end{itemize}
\end{KeyConcept}

\paragraph{The model this produces: primitive communism}
\begin{itemize}
    \item This model, according to Karl Marx, represented primitive communism --- a very simple, sweet picture of society .
    \item Primitive means old, ancient --- of that era .
    \item Communism here means the community gathered, produced and consumed --- everything happening in a communal fashion .
    \begin{itemize}
        \item ``Communal'' here does not mean religious conflict . It means community: gathered together, produced together, ate together, sat together around the campfire, danced together . Economic needs were fulfilled collectively .
    \end{itemize}
    \item Up to this stage Marx says there was no inequality .
    \begin{itemize}
        \item No inequality means no superior or inferior; no superior or inferior means no class; which simply means it was a classless society . There was no rich or poor --- everybody was equal .
    \end{itemize}
    \item The means of production at this stage were primitive tools --- stone, bow, arrow .
\end{itemize}

\begin{ExampleBox}
The teacher advised drawing the diagram of this stage according to your own convenience and available space --- the diagram is a personal memory aid, not a fixed format .
\end{ExampleBox}

\begin{QuickRevision}
\begin{itemize}
    \item Sequence: survival $\rightarrow$ production $\rightarrow$ collective labour $\rightarrow$ social relations of production $\rightarrow$ primary needs satisfied .
    \item This stage = primitive communism: communal ownership, communal production and consumption .
    \item No inequality $\rightarrow$ no superior/inferior $\rightarrow$ no class $\rightarrow$ classless society; no rich, no poor .
    \item Means of production at this stage: stone, bow, arrow .
\end{itemize}
\end{QuickRevision}

\subsection{Secondary Needs, New Technology and the Birth of Ownership}
\begin{itemize}
    \item Once the primary needs are met, secondary needs develop --- needs higher than primary: better food, better clothes, better shelter .
    \begin{itemize}
        \item A need is a need and cloth is cloth; but ``better'' clothes means branded clothes, ``better'' food means branded food, a better house means a house in a branded location . Needs increased .
    \end{itemize}
\end{itemize}

\paragraph{The lion and the iron cage}
\begin{itemize}
    \item Instead of a deer, they now want to hunt a lion . Hunting a lion is not easy, especially with the primitive tools these people had --- stone, bow, arrow .
    \item A lion is a huge, violent, wild animal; the primitive tools are not sufficient . Something else has to be developed to satisfy this new secondary need .
    \item So the means of production had to take an advanced form . The Stone Age is over; now comes the Iron Age --- iron tools .
    \item A few people had the idea that advanced tools were needed . Those people built an iron trap --- an iron cage --- inserted a bait, the lion got trapped, and they ate the lion .
\end{itemize}

\begin{KeyConcept}
\begin{itemize}
    \item The development of the means of production is nothing but technological advancement .
    \item Technological advancement then results in a change in the relations of production .
\end{itemize}
\end{KeyConcept}

\paragraph{How ownership was born}
\begin{itemize}
    \item In primitive communism the relations of production were symmetrical: no inequality, everybody equally involved in production, no superior, no inferior, no rich, no poor --- communal relations .
    \item With the advance in technology, only a few could develop the higher form of technology needed to satisfy the secondary needs . Those people became the owners .
    \item The next day, when others asked to borrow the iron cage to go hunting, that person said: ``This is mine.''  Even though all of them had collectively made it . This is how ownership came into existence .
    \begin{itemize}
        \item There were no laws then . Today you are the owner of your house and the proof is a paper . In those times there was no paper, no state, no law --- ownership was simply asserted: the idea is mine, so I am the owner .
    \end{itemize}
    \item Relations of production therefore moved from symmetrical to asymmetrical: some people are owners of the means of production; the rest, who are not owners, are dependent on them for the iron cage . Unless you have the iron cage you cannot trap the lion and cannot fulfil your increased desire .
\end{itemize}

\paragraph{The chain to remember}
\begin{itemize}
    \item Change in means of production $\rightarrow$ change in relations of production .
    \item Means of production reflect changes in technology .
    \item Therefore: change in technology ultimately led to change in the relations of production .
\end{itemize}

\begin{QuickRevision}
\begin{itemize}
    \item Primary needs satisfied $\rightarrow$ secondary needs develop (better food, cloth, shelter) .
    \item Old tools cannot meet new needs $\rightarrow$ iron tools, the iron cage $\rightarrow$ technological advancement .
    \item Only a few could develop the new technology; they claimed it --- this is the origin of ownership .
    \item Relations of production shift from symmetrical to asymmetrical: owners and dependants .
    \item Core line: change in technology $\rightarrow$ change in means of production $\rightarrow$ change in relations of production .
\end{itemize}
\end{QuickRevision}

\subsection{The Laws of Dialectics Operating in the Social World}
At this point the teacher stopped the narrative to show that Marx is not telling a story at random --- he is demonstrating how the laws of dialectics actually apply in the social world .

\paragraph{First law --- quantitative change leading to qualitative change}
\begin{itemize}
    \item Question put to the class: your needs were primary and now they are secondary . Is this quantitative or qualitative change?  Answer: it is quantitative .
    \item Why: till yesterday you were happy with one deer; now you want three deer, four deer . That is quantity increasing .
    \begin{itemize}
        \item Momos example: you go to the market and pay for one plate of momos; after finishing it you want one more plate --- that is quantitative . But when you have had momos every single day and get bored, and switch from momos to pizza --- that is qualitative .
        \item Similarly: from one deer to many deer is quantitative; from many deer to the lion is qualitative .
    \end{itemize}
    \item Quantitative changes accumulating and then a sudden qualitative change --- this is the first law . Marx is showing us how the laws are applied in the social world .
    \item Second application: change in technology --- from stone to iron --- is not quantitative, it is qualitative change .
    \item So there are two quantitative-to-qualitative transitions to keep in mind: primary to secondary need, and stone to iron technology .
\end{itemize}

\begin{CriticalPoint}
\begin{itemize}
    \item Do not treat the needs transition as qualitative at the first step --- the accumulation (one deer $\rightarrow$ many deer) is quantitative; only the jump (deer $\rightarrow$ lion, momos $\rightarrow$ pizza) is qualitative .
    \item The whole point of the exercise: Marx is applying the laws of dialectics to real social development, not merely narrating history .
\end{itemize}
\end{CriticalPoint}

\paragraph{Second law --- union of opposites, and conflict}
\begin{itemize}
    \item The result of quantitative change leading to qualitative change was the change in the relations of production . Relations that were earlier symmetrical are now hierarchical: one superior, one inferior .
    \item Yet both are important to each other . The owner of the cage needs those people --- today the cage may have to be repaired, tomorrow a new cage may have to be made . And those people are dependent on him because he possesses the cage; he rents it out to them .
    \item This is union --- the second law, first part: union of opposites .
    \item Conflict then arises: these people keep asking for the cage, and one fine day the owner refuses . So union is there and conflict is also there .
    \item In fact union implies conflict --- you are in union precisely because of two opposite needs .
    \begin{itemize}
        \item Marriage analogy: marriage is not just union, it is conflict as well . Two opposite needs are fulfilled by two opposite people . When two opposites attract, they also repel --- because in the first place they were opposites . Action and reaction .
    \end{itemize}
\end{itemize}

\begin{QuickRevision}
\begin{itemize}
    \item Primary $\rightarrow$ secondary needs = quantitative; deer $\rightarrow$ lion, momos $\rightarrow$ pizza = qualitative . First law in action .
    \item Stone $\rightarrow$ iron technology = qualitative change .
    \item Owner and dependant need each other $\rightarrow$ union of opposites (second law, first part) .
    \item Union implies conflict, because the union rests on two opposite needs --- marriage is the analogy given .
\end{itemize}
\end{QuickRevision}

\subsection{From Primitive Communism to Slavery}
\begin{itemize}
    \item With technological change the relations of production change; with the relations of production changing, the mode of production changes .
\end{itemize}

\paragraph{What ``mode of production'' means}
\begin{itemize}
    \item Mode of production simply implies how society organises itself and satisfies its needs .
    \item The mode of production in the first era was communism: the means of production were communally owned and the relations of production were symmetrical . It is the community that shares and owns; no class exists; no one person owns anything; everybody is equal .
\end{itemize}

\paragraph{The shift to slavery}
\begin{itemize}
    \item So much has changed that the mode of production must also change . From primitive communism we enter a new mode where one group is superior because it owns the means of production and the second group is inferior --- asymmetry .
    \item The new mode of production is slavery .
    \item Meaning: the people dependent on the iron cage have been enslaved . ``I will keep asking them to make cages for me whenever I want, because if they do not listen to me I will not give them my iron cage and they will die hungry.'' 
    \item They have become slaves; the owner has become master . No more community, no more classless society --- we now have classes: one master, all the others slaves .
\end{itemize}

\begin{ComparisonBox}
\begin{itemize}
    \item Primitive communism $\rightarrow$ communal ownership of MoP, symmetrical ROP, classless .
    \item Slavery $\rightarrow$ private ownership of MoP by the master, asymmetrical ROP, two classes: master and slaves .
\end{itemize}
\end{ComparisonBox}

\begin{QuickRevision}
\begin{itemize}
    \item Mode of production = the way society organises itself to satisfy its needs .
    \item Communism as a mode: communal ownership + symmetrical relations + no class .
    \item Slavery as a mode: master owns the means of production; the dependent group is enslaved .
    \item Chain: technology $\rightarrow$ means of production $\rightarrow$ relations of production $\rightarrow$ mode of production .
\end{itemize}
\end{QuickRevision}

\subsection{From Slavery to Feudalism}
\begin{itemize}
    \item In slavery there are two classes --- master and slaves . After some time, new needs develop again . The master has got bored of lion .
\end{itemize}

\paragraph{The discovery of precious stones and the birth of mining}
\begin{itemize}
    \item The slaves who went into the forest to hunt for the master came across something very shiny --- what we today call pearls, diamonds, jewellery .
    \item The slave did not know what to do with it and took it to the master . The master became very happy --- he did not know what it was, but he could figure out that it was very different from anything he had seen in his life .
    \item Visitors who came to the master admired it . The master then ordered the slaves to fetch more such stones --- otherwise no iron cage .
    \item They had found stones on the top of the soil; they reasoned that deeper in the soil there would be more . So they started mining the soil --- and that is how the mining industry came into existence .
    \begin{itemize}
        \item This is exactly what Marx is talking about: how things we take for granted actually got developed in the first place .
    \end{itemize}
\end{itemize}

\paragraph{New needs, new technology, new relations}
\begin{itemize}
    \item The master now wants precious stones, spices, and cloth made of silk rather than cloth made of leaf . His desires are changing; new needs are developing .
    \item With new needs, technology must again be developed --- mining the soil by hand might take five years to find anything, so a faster technology is required . New needs $\rightarrow$ new technology, as simple as that .
    \begin{itemize}
        \item The new technology consisted of basic tools . Compared with the technology we have attained they will always be called basic, but for that point of time they were quite advanced .
    \end{itemize}
    \item Most needs are now fulfilled by land: land gives cash crops, precious stones, precious textiles .
\end{itemize}

\begin{ExampleBox}
Indigo: a very precious colour, the royal colour --- royal blue . It is called ``royal'' because the very discovery of that blue happened in the age of masters and slaves, and the word royal is attached to it to this day .
\end{ExampleBox}

\begin{itemize}
    \item With the change in technology the relations of production change again . The masters are now not just masters --- they have become owners of the land . The land belongs to the master, and the slaves are now serfs .
    \item So: landlords are those people who at one point of time were the masters; serfs and peasants are those people who at one point of time were slaves . Only the names of the classes have changed --- the relation of production is very much asymmetrical still .
    \item From master--slave we have now entered feudal lord and serf, and with this comes the new mode of production: feudalism .
\end{itemize}

\paragraph{What each mode owned}
\begin{table}[h!]
\centering
\begin{tabularx}{\textwidth}{lYY}
\toprule
\rowcolor{AccentDark}
\thead{Mode of production} & \thead{Who owned what} & \thead{Classes} \\
\midrule
Primitive communism & No one person owned anything; the community was the owner & Classless \\
Ancient / slave society & The master owned all the tools --- and, for the first time, owned human beings & Master and slaves \\
Feudalism & Ownership of land; along with land the lord also held the serfs & Feudal lord and serf \\
\bottomrule
\end{tabularx}
\end{table}

\begin{CriticalPoint}
\begin{itemize}
    \item For the first time in human societies, under slavery man became the private property of another man . With the development of slavery, ownership of a human being came into existence .
    \item Feudalism means ownership of land --- and along with the land, the serfs .
\end{itemize}
\end{CriticalPoint}

\paragraph{The summary lines dictated in class}
\begin{enumerate}
    \item Technological advancement leads to changes in the means of production .
    \item Changes in the means of production lead to changes in the relations of production .
    \item Changes in both relations of production (ROP) and means of production (MoP) lead to change in the mode of production .
\end{enumerate}
Do not confuse ``means'' and ``mode'' of production . Mode of production was introduced here only because the word came up; it is explained fully in the next section .

\begin{QuickRevision}
\begin{itemize}
    \item New needs of the master (stones, spices, silk) $\rightarrow$ mining and new tools $\rightarrow$ land becomes the chief means of production .
    \item Masters become landlords; slaves become serfs; only the names of the classes change, the asymmetry persists .
    \item Slavery is the first society in which a human being becomes the private property of another human being .
    \item Three-line summary: technology $\rightarrow$ means of production $\rightarrow$ relations of production $\rightarrow$ (both together) mode of production .
\end{itemize}
\end{QuickRevision}

\subsection{Forms of Consciousness --- How Inequality Gets Justified}
The teacher flagged this as the most important thing in the whole topic .

\begin{itemize}
    \item The question nobody asks: when the master said ``this cage belongs to me, you cannot use it without my permission'' --- how did the slaves simply agree? 
    \begin{itemize}
        \item If a man tells his wife ``this is my property, you do not have a single piece of claim in it'', it is not a beautiful scenario --- there is conflict . So how did the master justify his superior status over the slaves? 
    \end{itemize}
    \item Here comes what Hegel called the spirit, the world spirit, consciousness .
    \item Marx says: ideas were now developed to justify the superiority of one class over another class . Masters started developing ideas to answer ``why should I obey you?'' 
    \item The masters consulted the priest --- or, more simply, consulted all the other masters --- and together they developed new ideas: let us justify his enslavement, his subordination, in the name of religion . If he does not obey us, he is committing sin; then he will not go to heaven .
    \item So: with asymmetrical relations developed new forms of consciousness, to justify class inequality .
    \begin{itemize}
        \item ``To justify class inequality'' simply means to justify why I should be the owner and you should listen to me, and why what I say is right and you are wrong . To justify that, you had to back yourself with something that could testify to your statement --- and that was religion, and law, and also the concept of family .
    \end{itemize}
\end{itemize}

\begin{KeyConcept}
\begin{itemize}
    \item Hegel: ideas come first . Marx: ideas are the result of material needs --- matter is primary .
    \item Ideas are developed to justify a claim over matter; religion, law and family arise as forms of consciousness that legitimise class inequality .
\end{itemize}
\end{KeyConcept}

\paragraph{Illustration 1 --- how the Bible was read to justify slavery}
\begin{itemize}
    \item The teacher first laid down the principle: facts do not speak for themselves --- you need someone to interpret the facts . The question is how the verses were \emph{interpreted} .
    \item The passage read out in class was the story of Noah's sons --- Shem, Ham and Japheth . Noah planted a vineyard, became drunk and lay uncovered in his tent . Ham saw his father's nakedness and told his brothers; Shem and Japheth walked backwards to cover him without looking . On waking, Noah cursed Ham's son Canaan to be a servant of servants to his brethren, blessed Shem, and said Canaan shall be servant to both Shem and Japheth . Noah lived 350 years after the flood .
    \item The reading offered in class: the one shown as drunk and uncovered was cast as the slave, and the curse is treated as divine sanction --- ``God has said it'' --- so the servitude stands justified .
    \item The critical questions the teacher then raised against that reading:
    \begin{itemize}
        \item What was so terrible about seeing Noah drunk? 
        \item Why curse Canaan rather than Ham? 
        \item How long was the servitude to last? 
        \item Surely Ham would have been the same colour as his brothers .
    \end{itemize}
    \item This is one illustration of how Christian slaveholders used scripture to justify slavery . The source referred to in class for the interested student was a BBC piece titled ``Attempts to justify slavery'' .
\end{itemize}

\paragraph{Illustration 2 --- the Laws of Manu}
\begin{itemize}
    \item The Manu Smriti extols both caste and class hierarchies, which in turn reinforce each other .
    \item It prescribes stringent punishments for transgressing the boundaries of birth-based professions, which if allowed will have terrible consequences .
    \item Chapter 8, shlokas 21--22 (as read in class): when a Shudra interprets the law for a king, his realm sinks like a cow in mud; the entire realm, stricken with famine and pestilence, quickly perishes when it is teeming with Shudras, overrun by infidels and devoid of twice-born people .
    \item Shloka 410: the king should make Shudras engage in the service of the twice-born .
    \item Shlokas 413--414: the Shudra was created solely to do slave labour for the Brahmin; even when released by his master, a Shudra is not freed from his slave status, for that is innate in him and nobody can remove it .
    \item The point being made: how religion --- Christianity in one case, Hinduism in the other --- and how law justified slavery .
\end{itemize}

\begin{CriticalPoint}
\begin{itemize}
    \item Slavery was legal . Today we call it a pathological social fact, but at one point of time it was a normal social fact .
    \item Things can be legal and unethical at the same time . Laws were developed, religion was developed, things were written down --- all to justify one person's status and class position .
\end{itemize}
\end{CriticalPoint}

\begin{QuickRevision}
\begin{itemize}
    \item Central question: why did the subordinated accept subordination? Answer: forms of consciousness .
    \item Hegel = ideas first; Marx = ideas are the result of material needs, matter is primary .
    \item Religion, law and family develop as justifications of class inequality .
    \item Two illustrations given: scriptural readings used by Christian slaveholders, and the Manu Smriti on caste-class hierarchy and Shudra servitude .
    \item Slavery was legal --- legality and ethics are two different things .
\end{itemize}
\end{QuickRevision}

\subsection{The Motor of History and the Superstructure}
\begin{itemize}
    \item The diagram drawn in class is not just a diagram --- it is the motor of history . Looking at societies from the top, Marx has analysed the evolution of societies, the evolution of material conditions leading to new relations of production and a new mode of production .
    \item At the surface, this is thesis, antithesis, synthesis --- a surface-level representation of the larger picture .
\end{itemize}

\paragraph{How the motor turns}
\begin{enumerate}
    \item With the development of higher needs, new technologies are developed to meet those needs .
    \item With new technologies, new relations of production develop .
    \item With new relations of production, a new mode of production comes into existence .
    \item Every mode of production has some internal economic contradictions --- exposed by higher needs or by the non-fulfilment of previous needs . Contradictions in the economy is the third law .
    \item Then human history develops towards a new mode of production; a new economic system and new social relations come into existence --- ancient to feudal, feudal to capitalism .
\end{enumerate}
An ever-moving motor does not stop . With this, Marx established the law of development of human history .

\begin{KeyConcept}
Overall conclusion of the theory: history is not the manifestation of spirit . History is the product of humanity and of the material needs of humanity . With ever-increasing needs, forms of consciousness also develop to justify the relations of production built around those needs .
\end{KeyConcept}

\begin{itemize}
    \item And therefore Karl Marx says: ``History of all society hitherto has been history of class struggle'' --- between masters and slaves, then lords and serfs, and then capitalism, where we are today .
    \item Whether we are still in capitalism or have progressed is something to be explored; but assume we are in so-called Marxian capitalism, as Marx described it during the Industrial Revolution . On that assumption there are only two classes: bourgeoisie and proletariat .
    \begin{itemize}
        \item Proletariats are the working class --- the teacher included himself and the students in it . The bourgeoisie are the owners of the coaching institutions and of firms; the people working in them are working-class people, dependent on the owners for survival .
    \end{itemize}
\end{itemize}

\paragraph{Why the conflict does not become revolution}
\begin{itemize}
    \item Just like any other mode of production, capitalism too will have conflict between these two classes . But what stops a worker from fighting his director? 
    \begin{itemize}
        \item Livelihood compulsion and job security --- of course, these things are there .
        \item But one more thing: family . Because I belong to a family I cannot take any decision independently . The family will say: he is your director, he has been feeding you, he gave you a job, and this is what you are doing to him?  These are the values we gave you, this is why we sent you to school and got you educated . So you go back and apologise --- and give a gift at Diwali .
    \end{itemize}
    \item Family is one system, one idea, which prevents revolution . It is a form of consciousness . You cannot see ``family'' --- you can see your mother, you can see the house you live in, but family itself is consciousness . And this consciousness does not let you think independently .
    \begin{itemize}
        \item Talk to working professionals who already have enough material needs met: they still work --- because of family, because of what a friend of the father would say, because the father is retiring, because people ask ``what are you doing these days, still studying?''  These are forms of consciousness that never let you think independently .
    \end{itemize}
    \item Media is also a form of consciousness, which will never let you think of bringing a revolution against your director . After a fight you open a video platform, watch a spiritual talk, listen to music that pleases your ears and pacifies your anger, or watch a movie --- and all the anger is over .
    \begin{itemize}
        \item This is what the critical theorists told us; that is why they are called neo-Marxists . This is the cultural industry . It works to subdue consciousness, to standardise the thinking process, and to create false desires so as to make you passive robots using false desires and false means . This prevents revolution .
    \end{itemize}
\end{itemize}

\paragraph{Infrastructure and superstructure}
\begin{itemize}
    \item Economy is the engine . Family, law, religion and nation are forms of consciousness .
    \item Because these forms of consciousness are situated on top of the economy --- which is the engine of society --- they constitute the superstructure; the economy is the infrastructure .
    \item The whole historical-materialism diagram was happening in the realm of the economy; the ideas were cropping up to justify what is there in the economy .
    \item Therefore the superstructure justifies the inequalities present in the infrastructure .
\end{itemize}

\begin{CriticalPoint}
Conclusion line dictated for writing: the theory of historical materialism reflects how societies develop class characters --- how societies develop social classes --- and how the relations of production between these classes are justified by the forms of consciousness, that is, by the superstructure .
\end{CriticalPoint}

\paragraph{Marx's place in the disciplines}
\begin{itemize}
    \item The line quoted in class on descending ``from heaven to earth'' --- that matter is primary and ideas are secondary --- was attributed in the lecture to Hegel [as stated in lecture] .
    \item Criticism of Marx comes later; it is not an easy task to criticise Karl Marx . You may love Marx or hate Marx, but you cannot ignore Marx .
    \item Today any academician --- economist, sociologist, psychologist, political scientist, philosopher --- is placed on a left wing / right wing divide, and that division came from Marxism .
    \item Your position in academics is the product of the opposition or the support that you subscribe to regarding Marxism . Either you are a Marxist or you are not --- which means your identity is somehow related to Marx . Even if you are not a Marxist, your identity is still a product of what people have agreed to believe is part of being a Marxist .
\end{itemize}

\begin{QuickRevision}
\begin{itemize}
    \item The motor of history: higher needs $\rightarrow$ new technology $\rightarrow$ new ROP $\rightarrow$ new MoP $\rightarrow$ internal contradictions $\rightarrow$ next mode .
    \item Every mode of production carries internal economic contradictions --- this is the third law .
    \item History is not the manifestation of spirit; it is the product of humanity and its material needs .
    \item ``History of all society hitherto has been history of class struggle'' .
    \item Family, media, law, religion and nation = forms of consciousness = superstructure; economy = infrastructure .
    \item The superstructure justifies the inequalities present in the infrastructure .
\end{itemize}
\end{QuickRevision}

\subsection{Doubts Taken in Class --- Terminology and the Origin of Ownership}

\paragraph{Can we call infrastructure ``base''?}
\begin{itemize}
    \item ``Base'' is simply an English word --- of course it is a base, everything stands on it . But ``infrastructure'' is the word used by Marx .
    \item If you are using the word superstructure, use infrastructure . Do not use ``economic base'' --- economic base is used in general language; it is of course right, but technically not sound .
\end{itemize}

\paragraph{Why could some people own the means of production and others could not?}
\begin{itemize}
    \item It is inevitable . Think of it this way: the day after tomorrow all of you will write the same test, but still some of you will be rank one, two, three and others will not .
    \item People will then ask to see the topper's copy, or take his or her help: ``how did you structure your answer, please help me'' . The moment you are taking his help, you have become subject to whatever he or she tells you --- and this is how hierarchy develops .
    \item This hierarchy is then legitimised in the form of the rank list .
    \item Talent also exists --- that is very much true .
    \begin{itemize}
        \item Linked forward: elite theory by Vilfredo Pareto --- some people are born elites . To be taken up separately .
    \end{itemize}
\end{itemize}

\begin{QuickRevision}
\begin{itemize}
    \item Use ``infrastructure'', not ``economic base'', when you are pairing it with ``superstructure'' .
    \item Differential ownership is treated as inevitable --- hierarchy emerges from differential performance and is then legitimised (the rank list is the classroom analogue) .
    \item Forward link: Pareto's elite theory --- some people are born elites .
\end{itemize}
\end{QuickRevision}

\subsection{Mode of Production --- Full Definition}

\paragraph{The motor analogy}
\begin{itemize}
    \item Everything is happening in the motor, and the name of the motor is the mode of production . That is the motor of history .
    \item The motor is rotating; history is evolving with the rotation of the motor; and the rotation of the motor involves conflict between classes . New modes of production develop; new classes come with new needs and new technology .
    \item The motor is situated in the base . Just as the motor of your car sits in a corner of the engine --- you never look at it, perhaps not in ten years --- you are concerned only with the steering wheel, brakes, clutches, gears, seat and music player .
    \item The mode of production, like any engine, consists of many parts . It is run by two wires: means of production (the first wire) and relations of production (the second wire) .
\end{itemize}

\paragraph{The three definitions to write down}

\paragraph{Mode of production}
\begin{itemize}
    \item Mode of production is defined as the way society organises itself to meet the needs of the population .
    \begin{itemize}
        \item Capitalism is a way society organises itself to meet needs . Socialism is a way society organises itself to meet needs . Feudalism is a way society organises itself to meet needs . Primitive communism was also a way society organised itself to meet needs .
        \item Question raised in passing: are they really \emph{our} needs, or has each need been manufactured as our need? 
    \end{itemize}
\end{itemize}

\paragraph{Means of production}
\begin{itemize}
    \item Means of production consist of technology, land, capital, machines and human labour (skilled, semi-skilled, unskilled --- so-called talented or untalented) . Land is also a form of capital .
    \item Meaning: the means of production are the tools required to fulfil the needs .
    \begin{itemize}
        \item Example given: if you want a branded pair of shoes, they will not fall from the sky . To manufacture them you must have skilled labourers, technological expertise, and the machine and assembly line .
    \end{itemize}
\end{itemize}

\paragraph{Relations of production}
\begin{itemize}
    \item Relations of production are the relationship between the classes involved in the production process .
    \begin{itemize}
        \item In the mode of production called communism the relations of production were symmetrical --- the relation between classes was not unequal, because there was no class at all; it was a classless society .
        \item Our needs increased; to gratify them technology increased; means of production reflect the change in technology; then relations of production changed; then classes came . So the relations of production change every time with the evolution of history .
    \end{itemize}
    \item Next line: the relationships between the classes are the product of their relationship to the means of production . Some classes are owners of the means of production and some are not --- and this is how the relationship between the classes is defined .
\end{itemize}

\begin{CriticalPoint}
This is only the first layer of knowledge --- new words, new terms . Do not get worried when you come across new words .
\end{CriticalPoint}

\paragraph{``Forces of production'' --- a translation artefact}
\begin{itemize}
    \item Marx was one, but Marxists are billions, because there is no one interpretation of Marx . Many people interpreted him according to their ideology, social location and education . When the author dies, interpretation begins .
    \item Force of production is a result of interpretation of the original writings . Marx in the original writings never described anything like ``force of production''; the originals are in German, and translators rendered the same term as ``means'' in some places and ``force'' in others .
    \item Means of production and forces of production are not different . It is only a difference of literature --- like the difference in pronunciation between American and British English .
    \begin{itemize}
        \item Some say force of production = technology + human labour . But human labour is part of the means of production as well, so there is no difference .
    \end{itemize}
\end{itemize}

\paragraph{A note on abbreviations}
\begin{itemize}
    \item When you write MOP in capitals, mean Mode of Production; write means of production in small letters --- so that later you can tell what you wrote .
\end{itemize}

\begin{QuickRevision}
\begin{itemize}
    \item Mode of production = the way society organises itself to meet the needs of the population .
    \item Means of production = technology, land, capital, machines, human labour --- the tools required to fulfil needs .
    \item Relations of production = relationship between the classes involved in production, defined by their relationship to the means of production .
    \item The mode of production is the motor of history, run by two wires: means of production and relations of production .
    \item Forces of production = means of production; the difference is only one of translation and interpretation .
\end{itemize}
\end{QuickRevision}

\subsection{The Cycle of Modes of Production}
\begin{itemize}
    \item It is the same engine completing one cycle after another . Only the cycles are being completed; the engine is the same .
    \item Cycle 1: the mode of production was primitive communism . One cycle completed --- new needs developed, new technology developed, relations of production changed with respect to the change in technology or means of production .
    \item Then comes the new mode of production: ancient society . Then feudalism . Then capitalism .
\end{itemize}

\begin{DataFact}
\begin{itemize}
    \item Marx died in the year 1883 .
    \item At that point there was no socialism of the kind Marx envisaged; he believed the world was living in the era of capitalism .
\end{itemize}
\end{DataFact}

\begin{itemize}
    \item Just like every previous mode of production collapsed because of inherent contradictions, capitalism will sow the seeds of its own destruction . Capitalism is destined --- doomed --- to collapse . And from the ruins, from the ashes of capitalism, will rise socialism, the new mode of production . This part is futuristic .
\end{itemize}

\paragraph{The crux line to write down}
\begin{itemize}
    \item All the modes of production developed as a result of inherent contradictions in the previous modes of production, which were the result of the contradiction between relations of production and means of production .
\end{itemize}

\paragraph{Where the line falls between history and prophecy}
\begin{itemize}
    \item Draw a line after capitalism . Till capitalism it was nothing but historical materialism . After capitalism it is prophecy . Marx becomes a prophet, telling the world that socialism is inevitable, communism is inevitable .
\end{itemize}

\paragraph{Dialectics inside the cycle}
\begin{itemize}
    \item In all the modes of production we had thesis, antithesis, then synthesis . The synthesis becomes a new thesis; antithesis develops again due to inherent contradictions in ancient society, and the new synthesis is feudalism .
    \item Therefore feudalism is the negation of negation --- the third law .
    \item Three forms of dialectics are going on simultaneously --- it is a very big, very complex engine:
    \begin{itemize}
        \item Dialectic between modes of production --- feudalism could not have developed had there been no dialectic between the two previous modes .
        \item Dialectic between relations of production and means of production within a mode of production .
        \item Dialectic between the classes --- for example between masters and slaves .
    \end{itemize}
\end{itemize}

\paragraph{Socialism and communism are not the same}
\begin{itemize}
    \item Socialism is not equal to communism --- these are two different things, and specifically to Karl Marx they are very different .
    \item Communism is the last stage --- the end of history . And it is not merely communism; it is pure communism . The first was primitive, the last is pure .
    \begin{itemize}
        \item ``History repeats itself, first as tragedy, second as farce'' --- communism comes again, this time as pure communism .
    \end{itemize}
\end{itemize}

\paragraph{The Asiatic mode of production}
\begin{itemize}
    \item There is one more mode of production which the teacher instructed students never to write along with the four: the Asiatic mode of production .
    \item Reasons: there is not enough literature available on it, and Marx himself never wrote extensively on it .
    \item The problem is with Indians and Asians: when Marx described the Asiatic mode of production he was writing about India and China; Indians got fascinated by his writings on India and started exploring it --- but there is nothing more there .
    \item Marx never visited India and did not have a real idea of India . He had an orientalist, European vision of India: to him India was all about villages --- self-sufficient villages .
    \begin{itemize}
        \item If needs are limited (Gandhi's ``limit your needs''), needs never grow; if needs never grow, technology never grows; if technology never grows, the means of production never change; if means never change, relations of production never change; and classes never emerge . We remained self-sufficient .
        \item It was colonialism --- the British --- that increased our needs, and then we jumped to capitalism .
    \end{itemize}
    \item This is to be dealt with in Paper 2 . The four that matter: primitive, ancient, feudal, capitalist --- plus socialism and then pure communism .
\end{itemize}

\begin{QuickRevision}
\begin{itemize}
    \item Same engine, successive cycles: primitive communism $\rightarrow$ ancient $\rightarrow$ feudal $\rightarrow$ capitalist $\rightarrow$ socialism $\rightarrow$ pure communism .
    \item Marx died in 1883, believing the era was capitalism; capitalism is doomed to collapse from its own inherent contradictions .
    \item Crux line: every mode of production arises from the inherent contradictions of the previous one --- contradictions between relations and means of production .
    \item Till capitalism = historical materialism; after capitalism = prophecy .
    \item Feudalism is the negation of negation (third law); three dialectics run at once --- between modes, between relations and means, and between classes .
    \item Socialism $\neq$ communism . Communism is the last stage, pure communism, the end of history .
    \item Asiatic mode of production: thin literature, Marx never wrote extensively, based on an orientalist view of India --- do not list it with the four .
\end{itemize}
\end{QuickRevision}

\subsection{Feudalism in Detail}

\paragraph{Definition dictated in class}
\begin{itemize}
    \item Feudalism is a mode of production characterised by land as the primary means of production, in relation to which we have two social classes: landlords and serfs (S-E-R-F-S) --- that is, landlords and peasants .
    \item Landlords are the owners of the land . Serfs do not own the land; serfs only work on the land .
\end{itemize}

\paragraph{The French Revolution estate system as illustration}
\begin{table}[h!]
\centering
\begin{tabularx}{\textwidth}{lYY}
\toprule
\rowcolor{AccentDark}
\thead{Estate} & \thead{Share of population} & \thead{Control over land} \\
\midrule
First Estate & Hardly about 2\% of the population & About 20\% of the land [as stated in lecture] \\
Second Estate & About 4--5\% of the population & Disproportionately large relative to numbers \\
Third Estate & More than 90\% of the population & Largely landless; the little land held was concentrated among the elites of the third estate \\
\bottomrule
\end{tabularx}
\end{table}

\begin{itemize}
    \item The three estates had disproportionate control over land relative to their population . The majority was in the third estate, and those people were largely landless --- mostly peasants or landless labourers with no control over land .
    \item The king was the owner of the land --- if not the real owner, at least the ceremonial, symbolic owner .
\end{itemize}

\paragraph{Why landlords appear}
\begin{enumerate}
    \item The king all by himself cannot manage so much land .
    \item So he asks the landlords --- the warrior class, the second-estate nobility --- to manage the land on his behalf . These are the lords, earls and dukes who now hold the king's land .
    \item These landlords themselves are not going to till the soil and harvest, so they employ serfs and peasants from the third estate to work on the land granted to them by the first estate .
    \item And this is how the relations of production evolve: means of production = land; classes = landlords and serfs; landlords exercise full power over serfs .
\end{enumerate}

\paragraph{Who gets the produce}
\begin{itemize}
    \item Serfs are in no way owners of the land --- they are only there to produce the crop from the land .
    \item After production, the maximum share of the crop is usurped / appropriated by the landlords and given to the king, because the land belongs to the king .
    \begin{itemize}
        \item A certain share is kept in the pocket of the landlord; a certain share is given to the serfs; and the larger share goes to the king .
    \end{itemize}
    \item The king is impressed by the services of the landlords and is very happy that they manage the land nicely . The landlords are also happy with the king, because the king reposes faith in them .
    \item What do the landlords give the king in return? Their services --- including fighting for the king, which is also one form of service .
    \item That means the mutual needs of each other are what maintain and establish the relations of production . This is analysing the state system through Marxism .
\end{itemize}

\begin{QuickRevision}
\begin{itemize}
    \item Feudalism = mode of production with land as the primary means of production; classes = landlords and serfs/peasants .
    \item French Revolution estates: $\sim$2\% First Estate, $\sim$4--5\% Second Estate, $>$90\% Third Estate, land control grossly disproportionate .
    \item The king is the ceremonial/symbolic owner; he delegates land management to the nobility, who employ serfs .
    \item Produce: largest share to the king via the landlord, a share retained by the landlord, a share to the serfs .
    \item Relations of production are held together by mutual need --- land for service, service including warfare .
\end{itemize}
\end{QuickRevision}

\subsection{Peasant Struggles Read as Class Struggle}
\begin{itemize}
    \item The teacher noted this applies to any peasant struggle in India --- there have been so many peasant movements .
    \item Examples named: Tebhaga Struggle, Indigo Revolt, Moplah Rebellion, Deccan Revolt --- all peasant movements .
\end{itemize}

\paragraph{Tebhaga Struggle}
\begin{itemize}
    \item Tebhaga means two-thirds .
    \item In Bengal, the farmers were sowing their crops --- whichever crop the landlords told them to produce .
    \begin{itemize}
        \item Just as in Champaran those people were producing indigo because they were instructed to do so by the British . You were not producing out of your own will; you were producing because someone had told you to produce that thing only .
    \end{itemize}
    \item This shows that the production process has become external to you --- you are not able to produce what you want .
    \begin{itemize}
        \item The same thing happened in slave society: slaves produced only what the masters told them to; slaves could not produce anything out of their own will; slaves were the private property of the masters .
    \end{itemize}
    \item In Tebhaga, out of the whole produce, two-thirds was being given to the landlords . The peasants said two-thirds is a lot; we will give you only one-third, because otherwise we would have no means available to satisfy our needs . And since two-thirds is a lot, we will protest against you . That is why it was called the Tebhaga Struggle .
\end{itemize}

\paragraph{Who led it}
\begin{itemize}
    \item The protest was led by the Kisan Sabha . There is no English version of the name because that is the name; it is a farmers' or peasants' organisation .
    \item The Kisan Sabha was a branch --- a pressure group --- of the Communist Party of India . And the communists are those people who always fight for the oppressed .
    \item The Communist Party of India was leading the Tebhaga struggle against the landlords . This shows nothing but class struggle .
\end{itemize}

\begin{CriticalPoint}
\begin{itemize}
    \item You would never normally read Tebhaga as class struggle between peasants and landlords --- but that is exactly what it is .
    \item Peasants refuse the two-thirds share, landlords do not agree, and struggle follows: class struggle between landlords and serfs .
\end{itemize}
\end{CriticalPoint}

\paragraph{Indigo Revolt, Champaran and the Moplah Rebellion}
\begin{itemize}
    \item The Indigo Revolt and the Champaran Satyagraha are two different events, but the Indigo Revolt too is nothing but class struggle .
    \item The Moplah Rebellion is also class struggle . The difference is only that of Hindu and Muslim --- otherwise these are classes: zamindars and peasants . The main struggle is economic .
    \begin{itemize}
        \item It was not just Hindu and Muslim --- of course that dimension is there . But religion as a superstructural institution legitimises inequality: being Muslim in a so-called Hindu society, the oppression was legitimised . The detailed mechanism is taken up later .
    \end{itemize}
\end{itemize}

\begin{QuickRevision}
\begin{itemize}
    \item Tebhaga = two-thirds; peasants demanded that the landlord's share be cut from two-thirds to one-third .
    \item Led by the Kisan Sabha, a pressure group / branch of the Communist Party of India .
    \item Production process external to the producer --- the same feature seen in slavery and in Champaran indigo .
    \item Indigo Revolt, Champaran, Deccan Revolt, Moplah Rebellion --- all readable as class struggle; in Moplah, religion is the superstructure legitimising an economic conflict .
\end{itemize}
\end{QuickRevision}

\subsection{The Collapse of Feudalism and the Bourgeois Revolution}
\begin{itemize}
    \item For any mode of production to collapse there must be a new form of needs, and corresponding to those new needs there must be new technology or new advancements .
\end{itemize}

\paragraph{What the new needs were --- the trade evidence}
\begin{itemize}
    \item From South Indian history --- Pallavas, Cholas, Pandyas --- there was a point of time when the Roman Empire was indebted to the Pallavas . Also the Cheras / Chera dynasty, Kerala --- ``Keralaputra'' .
    \item These were the people known as exporters of black gold --- pepper .
    \begin{itemize}
        \item The same pepper that we do not particularly value today, and some people even avoid, was a luxury item in those times . The only thing that could drain the wealth of Rome was that one black gold --- pepper .
    \end{itemize}
    \item Along with pepper: silk, textile and spices --- South Indian history is rich in these details .
    \item The Roman Empire had fallen victim to the debts it had accumulated because of the large imports of these precious items . So increased needs ultimately led to the fall of the empire .
\end{itemize}

\paragraph{The process in France}
\begin{itemize}
    \item The needs of the nobility and the kings increased; with the population increasing, their needs also increased .
    \item A lot of land was given out to serfs; many things were now being produced from land; the number of hours invested in production of land also increased .
    \item With growing needs, the repercussions had to be faced by the people working on the ground --- the serfs --- because whatever they produced was going to meet the desires and needs of the nobility, the kings and the upper classes .
\end{itemize}

\paragraph{The chain of consequences}
\begin{enumerate}
    \item Gradually towns and townships started getting developed .
    \item With ever-growing needs, technology changed; with technological change, production increased .
    \item With increased production, industries and factories got established .
    \item To seek more production of precious items, trade started --- and then international trade .
    \item When international trade happens, to seek more and more needs, colonialism happened . India was seen as the golden bird, and that wealth flowed towards the kings and queens who had these desires and needs .
    \item This whole chain led to the rise of industries .
\end{enumerate}

\paragraph{The birth of the merchant class}
\begin{itemize}
    \item The people in the third estate were the bourgeoisie --- the middle class . It was they who became the owners of the factories established to meet ever-increasing needs in the name of mass production .
    \item And that is how the merchant class came into existence . Otherwise we had only landlords .
    \item These merchants got into conflict with the landlords and the nobility .
    \item And that is why Marx would say it is a bourgeois revolution . We read it as the third estate overthrowing the first estate, but it is actually bourgeois replacing monarchs and nobility . What does it show? Class struggle .
    \begin{itemize}
        \item The teacher clarified that calling the French Revolution a bourgeois revolution is attributable not to Karl Marx per se but to Marxists .
    \end{itemize}
    \item When the merchants became the owners of factories they became the bourgeoisie --- not just a middle class any more; now the whole wealth was in their hands, because maximum production was happening in factories . The more you produce, the more you want, the more you become a millionaire .
    \item New forms of inequality came into existence: instead of serfs, we now had the working class . Most of those same serfs were now working for the bourgeoisie . The old serfs had become the new working class --- the proletariat .
\end{itemize}

\begin{CriticalPoint}
This is how you read history when you have put on the glasses of Marxism .
\end{CriticalPoint}

\paragraph{Doubts on this section}

\paragraph{Does ``inherent contradiction'' in the French Revolution mean the conflict between nobility and merchants?}
\begin{itemize}
    \item Yes . The first conflict was landlords versus serfs . Then, when needs increased, the merchants became owners of factories, and class struggle happened again --- this time between merchants and landlords, two powerful classes struggling with each other . And that is how aristocracy was replaced by industrial capitalist society .
\end{itemize}

\paragraph{Why did merchants and nobility struggle?}
\begin{itemize}
    \item Struggle for power . Struggle for resources . Because with the ownership of resources you gain power . Who holds the resources --- that is what struggle is all about .
    \begin{itemize}
        \item The same logic was illustrated with contemporary party politics: why are rival political parties in struggle? Power .
    \end{itemize}
\end{itemize}

\begin{QuickRevision}
\begin{itemize}
    \item Collapse requires new needs plus new technology .
    \item Trade evidence: pepper as ``black gold'', Roman indebtedness to South Indian dynasties, silk, spices, textiles .
    \item Chain: growing needs $\rightarrow$ towns $\rightarrow$ technology $\rightarrow$ production $\rightarrow$ factories $\rightarrow$ trade $\rightarrow$ international trade $\rightarrow$ colonialism $\rightarrow$ rise of industries .
    \item Third-estate bourgeoisie become factory owners $\rightarrow$ merchant class $\rightarrow$ conflict with landlords and nobility .
    \item French Revolution as bourgeois revolution: bourgeois replacing monarchs and nobility, not simply third estate over first .
    \item Serfs become the proletariat; aristocracy is replaced by industrial capitalist society .
    \item Root of the struggle: power and resources .
\end{itemize}
\end{QuickRevision}

\subsection{Capitalism, Socialism and Communism --- Working Meanings}

\paragraph{Capitalism}
\begin{itemize}
    \item Capitalism is an economic system which has organised itself to meet certain needs and to extract profit out of it .
    \item The main motive is profit . It is profit which is running this motor . Without profit this motor will collapse .
    \item In capitalism we have industries --- the industrial class, the bourgeoisie --- and the people working for the industrial class, the proletariat .
    \begin{itemize}
        \item Although in the real world we do not have these two binary opposites so cleanly; there are many complexities . But let us speak on Marx's behalf .
    \end{itemize}
\end{itemize}

\paragraph{Socialism}
\begin{itemize}
    \item Socialism is also a mode of production which has organised itself to meet the needs . But in socialism the factories will not be owned by one industrialist class, or two or three private enterprises --- the factories will be owned by the state .
    \item What we today call PSUs --- public sector undertakings: ONGC, BHEL, SAIL, GAIL; Navratna and Maharatna categories . These are factories owned by the state --- steel factories, lignite factories, and so on .
    \item There are some of the most important sectors where the state does not even allow private enterprise . That is also changing .
    \begin{itemize}
        \item As reported in newspapers, most of the things which were considered sacred and not open to private investment are now also being sold to private players . Against this, many people object, saying the ruling party is right-wing capitalist . They believe the state should be the guardian of the weaker sections; and that if state enterprises are sold to private players, and if FDI and FII are allowed in the most important enterprises, then it is a threat to sovereignty or security . [Details deferred in class.] 
    \end{itemize}
    \item So socialism is simply state ownership of resources or of the industries --- industries built to produce those resources, or to process them, or to distribute them .
\end{itemize}

\paragraph{From capitalism to socialism to communism}
\begin{enumerate}
    \item Marx believed capitalism will also fail, and from the ashes of capitalism, socialism will come . (How capitalism will fail is dealt with in the next class.) 
    \item In capitalism we have workers, and workers are oppressed --- so much exploitation . These are the workers who will head the state in socialism .
    \item Socialism will establish the dictatorship of the proletariat . The proletariat, who were the exploited class in capitalism, will now head or take charge of socialism and form a state .
    \item When they form the state, all the inequalities of capitalism will wither away --- because all resources will be distributed equally among everyone by the state, headed by the working class . No one will be richer, no one poorer .
    \item Once resources are equally distributed and produced with the help of cooperatives --- regional rural banks, village cooperatives, cooperative societies --- then comes communism, because the duty of the state is over . The state shall wither away .
    \item Communism will be stateless and classless, and production will be ``from each according to his ability, to each according to his needs'' . This is the happy picture of communism .
\end{enumerate}

\begin{KeyConcept}
Conclusion: socialism is an intermediary stage between capitalism and pure communism . Socialism is just a stage, not an end in itself . After socialism will come communism . Communism will be stateless; socialism has a state controlled by the working class carried over from capitalism .
\end{KeyConcept}

\paragraph{Why this is the materialist interpretation of history}
\begin{itemize}
    \item Everything begins from the economy: needs arising, then classes coming, then conflict happening, then evolution happening . This is materialist .
    \item Hegel by contrast talks of ideas --- monarchy, democracy, world spirit, some world-historical figure who will come and do something . But history does not change like that .
    \item Karl Marx goes into the details to tell us how these ideas are actually developed by the classes to justify their claim over limited resources, or to justify inequality in that mode of production . That is it .
\end{itemize}

\begin{QuickRevision}
\begin{itemize}
    \item Capitalism: economic system organised to meet needs and extract profit; profit is the sole driver of the motor .
    \item Socialism: state ownership of resources and industries --- the PSU model; dictatorship of the proletariat .
    \item Equal distribution $\rightarrow$ state's duty over $\rightarrow$ state withers away $\rightarrow$ communism: stateless, classless, ``from each according to his ability, to each according to his needs'' .
    \item Socialism is an intermediary stage, not an end in itself .
    \item Materialist interpretation: economy $\rightarrow$ needs $\rightarrow$ classes $\rightarrow$ conflict $\rightarrow$ evolution; against Hegel's primacy of ideas .
\end{itemize}
\end{QuickRevision}

\subsection{Remaining Doubts from the Class}

\paragraph{Relations of production in socialism}
\begin{itemize}
    \item Relations of production in socialism are between the state and the working class --- state and workers, with the state headed by workers . State as one class; all the workers as the second .
\end{itemize}

\paragraph{Meaning of consciousness according to Marx}
\begin{itemize}
    \item It is the being which determines the consciousness --- this is what Karl Marx says . From matter come the ideas .
    \item Therefore some people, reading Marx, would say it is false consciousness, and that true consciousness will come . How we measure something as true or false is a later detail .
\end{itemize}

\paragraph{Do the three laws function among the working class in socialism?}
\begin{itemize}
    \item The problem is that till the four stages, the motor was going on and history was being interpreted . After the four stages, Marx has become a prophet --- and when you speak like a prophet, you do not speak with logic .
    \item This is also one of the criticisms of Marx: whatever he said till capitalism, he does not repeat in socialism and communism; it becomes an altogether different story .
\end{itemize}

\paragraph{How does the third law of dialectics define the origin and evolution of patriarchy in the family?}
\begin{itemize}
    \item Engels talks about the origin of the family in private property . Marriage and family will be taught separately, and the origin of patriarchy and family will be discussed there .
    \item This is historical sociology . When you talk about the origin of patriarchy, the origin of private property, the origin of class inequality --- that is historical sociology, and it can be written in the History versus Sociology question .
    \begin{itemize}
        \item Marxism would not have been possible if there was no Industrial Revolution . Marxism was a response to the Industrial Revolution . Therefore sociology is rootless without history, but history is fruitless without sociology --- because reading the Industrial Revolution through the lens of Marx gives a flavour absent from other textbooks .
    \end{itemize}
\end{itemize}

\paragraph{Commodity fetishism --- brief introduction}
\begin{itemize}
    \item When Marx talks about fetishism he is talking about the fetishism of commodities .
    \begin{itemize}
        \item Example: people with a particular brand of phone --- a certain colour, a certain logo . This connects to conspicuous consumption, show-off, symbolic goods . You know that phone is not so different from a competitor's, but we have developed a fetish for that brand's products .
    \end{itemize}
    \item Fetishism is actually the name of a religion --- being in awe or reverence of something that seems very magical to you . So a brand seems magical to you; you will even buy its accessories at a high price, feeling there must be something inside it .
    \item Commodity fetishism is when you are not able to realise that the goods do not have any value per se --- and yet you fall in love with those goods .
\end{itemize}

\paragraph{Are needs also a form of consciousness?}
\begin{itemize}
    \item No . Needs are needs . Our needs lead to the rise of certain technological advancements which are going to fulfil those needs .
    \item But today our needs are designed . Our needs are no longer natural --- whatever you need today is a socially constructed need .
    \begin{itemize}
        \item Through television they tell you that if you do not have this, then you are left behind . Advertising uses nationalism --- nation as a form of consciousness --- to sell you commodities . You do not need that bike; but once it has been described as the heartbeat of the nation, you must buy it, otherwise you become anti-national . Needs are created .
    \end{itemize}
\end{itemize}

\paragraph{Terminology caution}
\begin{itemize}
    \item Do not say ``pure capitalism'' . The term is pure communism .
\end{itemize}

\begin{QuickRevision}
\begin{itemize}
    \item Socialism's ROP: state (headed by workers) and the workers .
    \item ``It is the being which determines the consciousness'' --- matter first, ideas after; hence false consciousness .
    \item After capitalism Marx speaks as a prophet, not with logic --- this is itself a standard criticism .
    \item Engels on the origin of the family in private property; the whole enquiry is historical sociology .
    \item Commodity fetishism: goods have no value per se, yet we revere them as magical --- fetishism is originally the name of a religion .
    \item Needs are not consciousness, but today's needs are socially constructed and manufactured, often through nationalism in advertising .
\end{itemize}
\end{QuickRevision}

\sectiondivider

\section{Understanding, Scientific Method and the Foundations of Class}

\subsection{Why the Paper Takes Time --- Diagnosis from the First Test}
\begin{itemize}
    \item More than 110 students wrote the first test . Those who did not attempt it were asked to explain why --- it is important to know the reason .
    \item Not attempting is not right, especially when it is your first test . Missing the first test shows a careless attitude --- that you do not give importance to these units or to tests .
    \begin{itemize}
        \item Context noted in class: the test fell on the date on which the prelims examination would have been conducted had it not been delayed .
    \end{itemize}
    \item Learning alone will not help . Writing is where you implement the policies you have learned . If implementation is not good, something has to be worked out at the implementation stage .
\end{itemize}

\paragraph{The main problem reported: time constraint}
\begin{itemize}
    \item Most of the time is gone in the thinking process . Of the one and a half hours, at least 20--30 minutes cumulatively went in thinking --- not in one go, but in fragments: five minutes of thinking, a few seconds of writing, then back into thinking .
    \item Spending 20--30 minutes thinking reflects one or two things:
    \begin{itemize}
        \item Lack of revision . Though revision is not quantifiable --- someone's four revisions may equal someone else's two .
        \item The larger problem: lack of understanding . It is not about revision only; it is about understanding what you are already revising .
    \end{itemize}
    \item When your understanding is crystal clear, revision will not take a lot of time . You can save a lot of revision time if your understanding is good .
\end{itemize}

\begin{QuickRevision}
\begin{itemize}
    \item 110+ students wrote the first test; time constraint was the common problem .
    \item 20--30 minutes lost to fragmented thinking during the paper .
    \item Two causes: lack of revision, and --- the bigger one --- lack of understanding .
    \item Good understanding shortens revision, not the other way round .
\end{itemize}
\end{QuickRevision}

\subsection{The Concrete Test of Understanding}
Let us not talk in vague terms like ``good understanding'' and ``lack of revision'' . How do we get to know, concretely, whether we have understood? 
\begin{itemize}
    \item Test 1: Suppose you have read positivism many times . Can you explain it to your father or mother --- to someone who does not know positivism? 
    \begin{itemize}
        \item ``Teaching'' here does not mean taking a one-hour lecture . If you can describe the concept in very simple language in two minutes, you have attained a good level of understanding .
    \end{itemize}
    \item Test 2 (the harder one): think of the poorest person who could not afford education . Imagine explaining Marxism to a rickshaw puller . If you have that level of understanding, revision is nothing .
    \item Then you revise only to remember keywords and facts --- you do not revise to understand, because understanding is already there .
\end{itemize}

\begin{ExampleBox}
Worked example --- explaining positivism simply . Applying the natural-science method to the social world is part of positivism . But what is positivism exactly? Simple version: positivism says the social world has certain laws; society is driven by laws . It is not just the universe out there that is governed by laws --- the social world is governed by laws too . The study of the laws of the social world is positivism .
\begin{itemize}
    \item The next question will be: give me an example --- which laws? How are our actions governed by laws? If you get stuck there, that is precisely what you need to work on .
\end{itemize}
\end{ExampleBox}

\paragraph{How to run this process on your own}
\begin{itemize}
    \item You do not require other people for this . Revise three or four times, then sit for five minutes and ask: what have I read exactly, and how would I explain this? 
    \item When you explain something to someone, you do not explain features --- you explain the essence, the meaning of the term, in simple language .
    \item Whatever you were struggling with --- the thoughts you could not answer --- you scribble on paper . That forms your own mind map, and it will help you forever . Refer the mind map first, then handouts or class recordings for revision .
\end{itemize}

\begin{QuickRevision}
\begin{itemize}
    \item Test of understanding: explain it in simple language in two minutes to someone with no background --- a parent, or a rickshaw puller .
    \item Explain the essence, not the features .
    \item Once understanding is secure, revision is only for keywords and facts .
    \item Scribble what you struggle with --- that becomes your own permanent mind map .
\end{itemize}
\end{QuickRevision}

\subsection{Positivism Illustrated --- Two Cases of the Nuclear Family}

\paragraph{Case 1 --- the actor at the centre}
\begin{itemize}
    \item ``I decided to stay in a nuclear family because I believe a nuclear family gives me freedom. In a joint family there is no freedom, because too many people are there who exercise control on you.'' 
    \begin{itemize}
        \item And they derive their authority to exercise control because of lineage --- because of the way the family has been running for generations: ``this is the way we have been conducting ourselves, you cannot do this or that.'' 
    \end{itemize}
    \item Who is at the centre of this whole statement --- the nuclear family, or me? It is me . I am deciding . And it gives the other person the impression that I can decide my own fate .
\end{itemize}

\paragraph{Case 2 --- the structure at the centre}
\begin{itemize}
    \item ``With the development of industrialisation, families got broken down from joint to nuclear.'' 
    \item Now the question: is it my choice, or is it structurally that nuclear families are coming into existence? Is my choice itself the result of industrialisation? 
    \item Who is bigger --- me exercising my decision, or the structural process of industrialisation which actually gave me the power to choose, because industrialisation requires small families? 
    \item Case 1 is all about me . Case 2 sounds like a law . ``With industrialisation, joint families will break down into nuclear'' --- that is a law . And that is why we always talk about nuclear families in industrial societies .
\end{itemize}

\begin{ExampleBox}
Chemistry analogy: the joint family is one chemical substance . React it with a catalyst named industrialisation . What do you get? Nuclear families + individualism + fragments of the joint family .
\end{ExampleBox}

\paragraph{Which case explains positivism?}
\begin{itemize}
    \item Case 2 . In the first case I am choosing whom I want to stay with . Case 2 tells me it is not you who is choosing --- structures and larger forces are governing your very act of choosing . The idea that you are choosing at all reflects that the larger process has succeeded in creating individualism .
\end{itemize}

\paragraph{The 1920s counterfactual}
\begin{itemize}
    \item Imagine a hundred years back, in the 1920s . Nobody really had the power to break away from a joint family, because there was nowhere to go: 
    \begin{itemize}
        \item No job opportunities in urban centres of the kind we have today .
        \item No scope for migration; and even where migration existed, it was not at a mass level .
        \item The means of transport and communication were not there .
        \item You would not have received information about employment in an urban centre, because there was no newspaper as such --- the printing press existed, but not as a mass phenomenon as it is today .
    \end{itemize}
    \item So there was no point in breaking away . Today, if I break away from the joint family, it is only because of the power of individualism .
    \item But why I want to form a nuclear family is a separate matter of enquiry . I have attached certain meanings to the nuclear family --- I believe it gives freedom . In that case it is I who am at the centre; another person may attach a different meaning to the nuclear family .
\end{itemize}

\begin{QuickRevision}
\begin{itemize}
    \item Case 1 = the individual at the centre; Case 2 = structure at the centre and sounds like a law .
    \item Case 2 is the positivist statement: larger forces govern even the act of choosing .
    \item Industrialisation as catalyst $\rightarrow$ nuclear families + individualism + fragments of the joint family .
    \item The 1920s had no urban jobs, mass migration, transport or mass press --- so no structural possibility of breaking away .
    \item The meaning an actor attaches to the nuclear family is a separate line of enquiry .
\end{itemize}
\end{QuickRevision}

\subsection{Durkheim --- Positivist in One Place, Not in Another}
\begin{itemize}
    \item Durkheim also gave us certain laws . On his study of suicide he states that suicide rates will increase with industrialisation --- that is a law-like statement, a clear cause and effect . Reading that theory, Durkheim seems a positivist .
    \item But when reading his theory of religion, he does not seem a positivist scholar --- because in religion he gives no laws . Although he tried to study religion in a scientific fashion, he did not develop laws out of it .
    \item To study something scientifically and to develop laws are two different things . Many scientists study everything in a scientific fashion; that does not mean they come up with a law every time .
    \begin{itemize}
        \item When you conducted experiments in the school chemistry or physics lab, were you developing laws? No --- you were conducting experiments, and those experiments were conducted in a scientific fashion .
    \end{itemize}
    \item Law development is the higher stage . Durkheim did conduct a scientific study on the elementary forms of religion --- that is not the same as producing a law .
\end{itemize}

\begin{CriticalPoint}
A scholar is not positivist or non-positivist wholesale --- the same scholar can be positivist in one work and not in another . Judge the work, not the name .
\end{CriticalPoint}

\begin{QuickRevision}
\begin{itemize}
    \item Durkheim on suicide $\rightarrow$ law-like, positivist . Durkheim on religion $\rightarrow$ scientific study, but no laws, so not positivist there .
    \item Studying scientifically $\neq$ producing laws; law-building is a higher stage .
\end{itemize}
\end{QuickRevision}

\subsection{Why ``Elementary Forms''? --- The Power of Basic Questions}
\begin{itemize}
    \item The questions that check your basic knowledge are the toughest ones . When the teacher said the paper would be the toughest, what he meant was that it would be very basic .
    \item Basic question: why did Durkheim study the elementary form of religion and not the modern form? 
    \item Answer: in order to understand something complex, you first have to understand the basic . Religion is a very complex phenomenon; to understand it we must understand the basic .
\end{itemize}

\begin{ExampleBox}
\begin{itemize}
    \item Why do scientists study the composition of asteroids? Ask the scientist and he will say: because I want to explain the origin of the universe --- how the universe came into existence, how the earth got formed . To study that, he studies a small rocky particle .
    \item That scientist is also studying the elementary form of the existence of the universe . The asteroid, that rocky composition, is that elementary form .
\end{itemize}
\end{ExampleBox}

\paragraph{Why this could not have been taught earlier}
\begin{itemize}
    \item If this had been said twenty days back, you would not have understood it . Things grow organically . This is the next layer .
    \item Knowledge is cumulative --- it adds on to something that already exists . You will understand today because a base is already established in your mind .
\end{itemize}

\paragraph{Why the test mattered}
\begin{itemize}
    \item Before a test, we learn, we understand, we memorise --- we do not think . Thinking is what takes 20--30 minutes on the day of the test, because you never thought before .
    \item When you write a test you start thinking on your own; you get into the roots of the complexities . And when you analyse the roots of those complex things, you realise: this was so easy . It is an organic process .
\end{itemize}

\begin{QuickRevision}
\begin{itemize}
    \item Basic questions are the toughest; a ``tough'' paper here means a basic paper .
    \item Durkheim studied the elementary form because the basic must be grasped before the complex --- the asteroid/origin-of-universe analogy .
    \item Knowledge is cumulative and organic; the test is what forces the thinking that learning and memorising never do .
\end{itemize}
\end{QuickRevision}

\subsection{Always Write Examples}
\begin{itemize}
    \item Whenever you study anything theoretical, always look out for examples .
    \item Unless you have examples in your head to justify what you have written in theory, your answer will not reflect the level of understanding expected from a candidate .
    \item How else would the other person know that you have understanding? Only when you write examples .
    \item If you do not write examples, you are not showing your true potential to the examiner . Always make it a point to write examples .
\end{itemize}

\begin{QuickRevision}
\begin{itemize}
    \item Theory without examples cannot demonstrate understanding to an examiner .
    \item Collect examples as you study each theoretical point --- this is an answer-writing requirement, not a decoration .
\end{itemize}
\end{QuickRevision}

\subsection{Where the Syllabus Sits --- BA, BA Honours or MA?}
\begin{itemize}
    \item The sociology syllabus is very much in alignment with the BA --- but not exactly the bachelor's level . It is more like an intermediary stage between BA and MA: it is BA honours .
    \item When you look at the syllabus you will see that the BA syllabus by itself, and the understanding you have at BA, is not sufficient .
    \item Over the last two years, questions have also been asked from areas taught at the MA level .
\end{itemize}

\begin{CriticalPoint}
\begin{itemize}
    \item Advice given: do not chase MA-level books for one stray question . Out of 20 questions, suppose one is from the MA level --- if you start hunting for books to answer that one, you may end up missing the larger picture, that is, the 19 questions you are supposed to do .
    \item You would only get the answer to that one question asked; you would still not be prepared to tackle such questions in future .
    \item We do not have to ace all the questions --- we attempt 19 out of 28, and we have choice in that respect [figures as stated in lecture] .
\end{itemize}
\end{CriticalPoint}

\begin{QuickRevision}
\begin{itemize}
    \item Syllabus level = BA honours: between BA and MA .
    \item Some recent questions come from MA-level areas .
    \item Do not sacrifice the 19 answerable questions chasing the one that is not .
\end{itemize}
\end{QuickRevision}

\subsection{Which Mode of Production Does India Have?}
The question was deliberately asked without having been taught, to force independent thinking: ``I do not really want to teach you everything --- I want you to be always in a thinking process.'' 

\begin{table}[h!]
\centering
\begin{tabularx}{\textwidth}{lYY}
\toprule
\rowcolor{AccentDark}
\thead{Mode of production} & \thead{Present in India?} & \thead{Where / how} \\
\midrule
Primitive communism & Yes & Tribal communities --- Gonds and other tribals \\
Feudal & Yes & UP and Bihar as a highly feudal belt; zamindari, land ownership and landlords; the landlord (``Chaudhari sahab'') mentality and feudal mindset in Haryana, UP and Bihar \\
Capitalist & Yes & The industrial and market economy \\
Socialist & Yes & State ownership and the welfare apparatus \\
\bottomrule
\end{tabularx}
\end{table}

\begin{itemize}
    \item So we have all the modes of production --- and it is not only us; many countries have all these modes .
    \item And that is why: no country is purely capitalist, no country is purely socialist . These are only models .
    \begin{itemize}
        \item A related aside: ethnocentric behaviour --- assertions of caste/community pride among groups such as Jats and Gujjars --- has become highly visible only over the last ten years or so . Where did it come from? To be explained when ethnicity and ethnocentrism are taught .
    \end{itemize}
\end{itemize}

\paragraph{Why ``model'' is the right word --- the light analogy}
\begin{itemize}
    \item When you read the composition of light --- particles, waves, photons --- in reality you would not see only photons, only waves or only particles . You would see a mix .
    \item The wave paradigm tells us light is waves; the photon paradigm tells us light consists of photons . Each paradigm is a theoretical perspective; when you look through it, that is all you see .
    \item So in science also we have models . Study the universe through the Einsteinian perspective and you explain it with one model; through the Newtonian perspective, with a different model .
    \item Science too uses models to explain reality --- and social scientists are no exception . They have Einstein, we have Marx; they have Newton, we have Max Weber .
    \begin{itemize}
        \item And these are people you will meet outside sociology as well --- in political science, in English and literature courses, in linguistics, in cultural theory, in anthropology, and in philosophy .
    \end{itemize}
\end{itemize}

\begin{QuickRevision}
\begin{itemize}
    \item India contains all modes of production simultaneously --- primitive communist, feudal, capitalist and socialist .
    \item No country is purely capitalist or purely socialist; modes of production are models, not descriptions .
    \item Light analogy: wave/particle/photon paradigms each show only what they are built to show .
    \item Science itself works with models; Marx and Weber are sociology's Einstein and Newton .
\end{itemize}
\end{QuickRevision}

\subsection{The Steps of the Scientific Method}
Illustrated end-to-end with a single running example: how many environment questions will appear in the prelims? 
\begin{enumerate}
    \item Research question . You cannot follow a method unless you have a question to get an answer to . Here: what percentage of questions in the prelims will be from environment? 
    \item Review of literature . To answer it, you first refer to the previous year questions --- say five years of them . That is your literature review .
    \item Hypothesis . After reviewing, you form an assumption: approximately 30 questions will be from environment .
    \item Data collection . You appear in the examination and see 35 questions from environment . The data is: 35 questions .
    \item Verification of the hypothesis against that data .
\end{enumerate}
\begin{itemize}
    \item And then you become a prophet predicting the future --- ``every year 35 questions will come, mark my words'' --- and one year only two questions from environment appear .
\end{itemize}

\begin{CriticalPoint}
The way you ask a question can itself influence the opinion of the people you are asking --- linked forward to the Hawthorne effect, to be read separately .
\end{CriticalPoint}

\begin{QuickRevision}
\begin{itemize}
    \item Five steps: research question $\rightarrow$ review of literature $\rightarrow$ hypothesis $\rightarrow$ data collection $\rightarrow$ verification .
    \item The framing of the question can shape the response --- Hawthorne effect flagged for later .
\end{itemize}
\end{QuickRevision}

\subsection{What ``Scientific'' Means --- Testability and Falsifiability}

\paragraph{Remove the preconceived notion}
\begin{itemize}
    \item Scientific does not mean something is absolutely true . If I say this is a scientific statement or a scientific theory, that does not mean the theory is absolutely correct .
    \item It only means that the theory or statement is testable . You can test it; it is not a random statement that cannot be tested .
    \item It is not scientific because it is absolutely true --- because even in science there is no absolute truth . We do not have absolute theories or laws to explain the functioning of the universe .
    \begin{itemize}
        \item Example: ``every year 35 questions will come from environment in the examination'' is a testable statement --- therefore scientific .
    \end{itemize}
\end{itemize}

\paragraph{Whose concept is this?}
\begin{itemize}
    \item When we say something is scientific or not scientific --- is that according to Popper, or is it a universal concept? It is a universal concept .
    \item What is specific to Karl Popper is the science--pseudoscience demarcation principle . Every natural scientist discusses how we determine the scientificity of something, and when an experiment is to be called scientific .
    \item Even the greatest scientists in history agree there is no absolute truth in the natural world; they are still exploring, and science is still maturing .
    \item Science does not claim to have a monopoly of truth . Some individual scientists may sound arrogant, criticising social science, religion or magic --- but science as a body of knowledge makes no such claim .
    \item Why? Because science rests on organised skepticism . Scientists accept criticism precisely because they know the knowledge they possess is not absolute; otherwise there would be no reason to accept criticism . They are humble enough to accept loopholes in their theories, because they are thinking of the contributions their theories make to the larger world rather than of themselves .
\end{itemize}

\paragraph{Testing a social-science statement}
\begin{itemize}
    \item Take: ``Higher unemployment is the source of a higher crime rate.''  Does this not sound like a law?  It is a cause--effect relation --- concomitant variation, exactly like ``higher toxic substance in the water leads to a higher toxicity level of the water'' .
    \item How would I test it? Go through the composition of all the criminals and see how many of them were unemployed .
    \begin{itemize}
        \item If, out of a hundred criminals, 90 or more are unemployed, the theory is valid .
        \item But if only 20 or 30 are unemployed and 70\% are well employed --- driving for ride-hailing services, working with delivery platforms, working in large consulting firms --- they are also criminals: white-collar criminals .
    \end{itemize}
\end{itemize}

\begin{KeyConcept}
\begin{itemize}
    \item When we hear the word ``crime'' we start thinking of poor people stealing money or attacking someone . But the biggest crime is money laundering, and the biggest source of black money and money laundering happens at the level of white-collar people .
    \item That crime is the consequence of employment, not unemployment --- you are employed in sectors where there is either no state regulation, or the regulation lacks stringent provisions .
\end{itemize}
\end{KeyConcept}

\paragraph{Generalisation}
\begin{itemize}
    \item Is this not generalisation from 100 samples? Yes --- and anything that seems to be a generalisation is fine as long as the generalisation is testable .
    \item You will find generalisation in natural science also, and they are able to test those generalisations .
    \begin{itemize}
        \item Einstein believed the great solar eclipse of 1919 would tell him whether light bends . There is no other reason for him to wait for that eclipse --- the idea was generalised in his mind . Nothing is wrong in that generalisation per se, as long as it is testable --- and it was .
    \end{itemize}
    \item But not every generalisation is testable . In the social world many things cannot be tested because of the way they are defined: how do you define crime? What is your criterion for categorising someone as employed or unemployed? The social world has this many subjectivities .
    \item Therefore, for something to be testable it must first be concretely defined .
\end{itemize}

\paragraph{Induction, deduction and the swans}
\begin{itemize}
    \item In inductive reasoning you go from the few to the general . Going to the general is not wrong; but if the generalisation is not testable, it is not scientific .
    \item ``All swans are white'' --- you observed a few white swans and developed the generalisation . This is easily testable, and that is why it is scientific . It is wrong --- that is a different thing --- but it is scientific because it is testable .
    \item Did we test it initially? No . Deductive reasoning is about testing that which was hypothesised . And this is what Karl Popper's version of science is all about .
\end{itemize}

\paragraph{Testing means trying to prove wrong}
\begin{itemize}
    \item Testing does not mean proving it right .
    \begin{itemize}
        \item If I mark an answer in prelims and then search only for statements that prove my answer right --- never considering that I may be wrong --- that is not a test, however much it looks like one .
    \end{itemize}
    \item Testable means falsifiable . Falsifiable means the statement, law or theory should have the capacity to be proven wrong . If you cannot prove something wrong, it is not scientific --- anything at all can be proven right in this world; there is not even one thing that cannot be proven right .
\end{itemize}

\paragraph{Two contrasting examples}
\begin{table}[h!]
\centering
\begin{tabularx}{\textwidth}{lYY}
\toprule
\rowcolor{AccentDark}
\thead{Statement} & \thead{Scientific?} & \thead{Why} \\
\midrule
``Yesterday Santa Claus came to my room and gave me this gift.'' & No & It cannot be tested. Here nonsense and untestability happen to coincide. Saying ``many people report Santa gifting them things'' is only verification, not a test in Popper's sense --- you cannot think of a condition that would falsify it. \\
``There is one planet between the Sun and Venus.'' & Yes & Not because you already know Mercury is there. It is scientific because it has the power to be proven wrong: go and look; if no planet is seen, the statement had the capacity to be falsified. If a planet is seen, it was tested right --- but the condition of being provable wrong existed in the first place. \\
\bottomrule
\end{tabularx}
\end{table}

\begin{QuickRevision}
\begin{itemize}
    \item Scientific $\neq$ true . Scientific = testable .
    \item Testable = falsifiable = has the capacity to be proven wrong (Popper) .
    \item Popper's own contribution is specifically the science--pseudoscience demarcation principle; ``scientific'' itself is a universal concept .
    \item Science rests on organised skepticism and claims no monopoly of truth .
    \item Induction goes few $\rightarrow$ general; deduction tests the hypothesis . ``All swans are white'' is scientific though false .
    \item Concrete definition is a precondition of testability --- hence the difficulty in the social world .
\end{itemize}
\end{QuickRevision}

\subsection{Applying It to a Previous Year Question}

\paragraph{The PYQ that was deferred}
\begin{itemize}
    \item There was one previous year question set aside for Unit 4: ``Is non-positivist methodology scientific?'' 
    \item You cannot attempt that question properly if you do not know what ``scientific'' means . If you know non-positivism but assume positivism = science, you will automatically conclude that non-positivism is not science .
    \item The question is actually very basic --- and once your basics are strong, you will attempt it very nicely .
\end{itemize}

\paragraph{How to think through the unemployment--crime statement in an answer}
\begin{itemize}
    \item Which theoretical strand would say something like this? Obviously you cannot expect a phenomenologist to speak like this --- phenomenology would never say there is one cause and effect in the social world . There are multiple causes behind multiple effects, and objectivity is an illusion .
    \item But here I can see a highly objective statement which looks like a cause--effect proposition, highly scientific, something that can be tested . That is a positivist statement .
    \item Positivism has two branches --- functionalist and Marxist . Do we need to go into those details? You have the option of not going into them . But if you do:
    \begin{itemize}
        \item You would not be talking about functionalism here [as stated in lecture] .
        \item Marxism as a philosophy would not support this statement either . To Karl Marx, the bourgeoisie are the criminals --- the owning class, the ruling class is the criminal, not the working class . If you have defined crime as crime committed by working-class people, it is not a Marxist statement; Marxism would never support it .
    \end{itemize}
    \item At the surface, though --- removing all subjectivities, not knowing how crime or unemployment have been defined --- the statement reads as one given by a positivist scholar who believes in cause and effect in the social world . That much you can say without going into the nitty-gritties of functionalist versus Marxist .
\end{itemize}

\begin{CriticalPoint}
A note on how to read the scholars: come out of the idea that you are reading a subject for 500 marks . Marx never wrote for sociology; he did not know that a subject would exist one day and that his philosophy would be labelled positivist .
\end{CriticalPoint}

\paragraph{The answer-writing method for such a question}
\begin{enumerate}
    \item Establish that the statement is testable, therefore scientific --- and that ``scientific'' here means falsifiable, not true .
    \item Think about possibilities beyond the given statement . You do not have to quote data or reports .
    \item Ask: if it is falsifiable, what could be the possible reasons for the statement being wrong? A leads to B --- but perhaps it is C that leads to B, or D, or E . Perhaps it is not higher unemployment but some other factor that is the source of a higher crime rate .
\end{enumerate}

\begin{QuickRevision}
\begin{itemize}
    \item PYQ: ``Is non-positivist methodology scientific?'' --- unanswerable without a clear meaning of ``scientific'' .
    \item A cause--effect, objective proposition = positivist; a phenomenologist would reject it outright .
    \item Marxism would also reject the unemployment--crime claim, because to Marx the owning class is the criminal class .
    \item Answer method: establish testability/falsifiability, then generate alternative explanations rather than citing data .
\end{itemize}
\end{QuickRevision}

\subsection{Class Struggle Begins --- What is Class?}
\begin{itemize}
    \item Mode of production is over; the next concept is class struggle . But before understanding class struggle we should understand class, and what class means in the language of Marx .
\end{itemize}

\paragraph{The historical sequence recalled}
\begin{table}[h!]
\centering
\begin{tabularx}{\textwidth}{lY}
\toprule
\rowcolor{AccentDark}
\thead{Society} & \thead{Classes} \\
\midrule
Primitive society & Classless \\
Ancient society & Master and slaves \\
Feudal society & Lords and serfs \\
Capitalist society & Bourgeoisie and proletariat \\
\bottomrule
\end{tabularx}
\end{table}

\begin{itemize}
    \item This evolution of history has its roots in material conditions --- and that is clearly the materialist interpretation of history, done in terms of the struggle between classes .
\end{itemize}

\paragraph{How did classes come into existence?}
\begin{enumerate}
    \item Secondary needs develop after the gratification of the primary needs .
    \item Whoever developed the technology to satisfy those new needs possessed ownership over that means of production .
    \item Once they possessed ownership over the means of production, it changed the relations of production .
    \item Otherwise the relations of production were very much communal --- harmonious, peaceful, everybody doing similar work, with similar food, clothing, lifestyle and shelter .
    \item The moment needs developed, class formation started happening . That class was the owner of the technology or means of production through which the new needs could be gratified .
\end{enumerate}
\begin{itemize}
    \item ``Hitherto'' means until now --- that is, up to capitalism . So: the history of all societies hitherto has been the history of class struggle .
\end{itemize}

\paragraph{The definition}
\begin{itemize}
    \item Remember this: Marx himself never defined class . But from the writings of Marx we can conclude a definition .
\end{itemize}

\begin{KeyConcept}
Class, to Marx, is a group of people having similar relations to the means of production .
\end{KeyConcept}

\paragraph{A note on the two words}
\begin{itemize}
    \item Bourgeoisie is a noun; bourgeois is the adjectival form . Bourgeoisie simply means the owner of the means of production --- the industrialist class, the so-called capitalist class .
\end{itemize}

\paragraph{Applying the definition}
\begin{itemize}
    \item Bourgeoisie as a class: a group of people having similar relations to the means of production --- and the similar relation is that all of them are owners of the means of production .
    \item Proletariat as a class: a group of people having similar relations to the means of production --- the similar relation being that they do not own the means of production, and all of them are selling their labour to the capitalist class .
\end{itemize}

\begin{ExampleBox}
\begin{itemize}
    \item Take a person employed to cook noodles at someone's shop, and the teacher employed to teach one subject . One store is run by a man who has employed someone to cook; another, a level up, is run by a man who has employed someone to teach sociology .
    \item Both employees are the same --- similarly placed in the social world . Neither is an owner, and the relation uniting them is that both are there to sell their labour .
\end{itemize}
\end{ExampleBox}

\begin{QuickRevision}
\begin{itemize}
    \item Classless $\rightarrow$ master/slave $\rightarrow$ lord/serf $\rightarrow$ bourgeois/proletariat; the evolution is rooted in material conditions .
    \item Classes arise when the developers of new technology take ownership of the new means of production .
    \item ``Hitherto'' = until now, i.e. up to capitalism .
    \item Marx never defined class; the working definition: a group of people having similar relations to the means of production .
    \item Bourgeoisie = all owners; proletariat = all non-owners who sell their labour . Occupation and status are irrelevant to the classification .
\end{itemize}
\end{QuickRevision}

\subsection{Class as Objective Experience and as Subjective Condition}
\begin{itemize}
    \item There is a concept of the objective experience and the subjective condition of class . These are two different concepts .
    \item Class as objective experience simply means that the workers have realised that they are exploited --- and they know how they are exploited .
    \item Class has a subjective condition when workers do not realise that they are exploited --- when workers are living in false consciousness .
\end{itemize}

\begin{ExampleBox}
\begin{itemize}
    \item Suppose today I am very happy: I am making enough money, I am able to do whatever I want, express myself, take time out to meet friends, enjoy every luxury I have seen on television . So I am happy .
    \item But to Karl Marx I am stupid --- a highly idiotic person --- because I am not able to realise that I have been exploited, and that I am exploited every day .
    \item Only Marx knows how exploitation takes place in such a sweet fashion that you would never realise you were exploited .
\end{itemize}
\end{ExampleBox}

Note for the reader: in the following class the same pair of terms is presented with the labels attached the other way round --- subjective dimension = class in itself, objective dimension = class for itself . Both formulations as delivered are preserved in these notes; see section~\ref{sec:class-in-for-itself} .

\begin{QuickRevision}
\begin{itemize}
    \item Objective experience of class = workers realise they are exploited and know how .
    \item Subjective condition of class = workers do not realise it; they live in false consciousness .
    \item Exploitation in capitalism is sweet-tasting --- that is precisely why it is not noticed .
\end{itemize}
\end{QuickRevision}

\subsection{Labour Power, Use Value and Exchange Value}
\begin{itemize}
    \item When you get a job, what are you doing exactly? You are selling your labour .
    \item You get a certain amount per month only because you have also given something . What have you given the company in exchange? Your labour --- and, to be more precise, your labour power .
    \item Labour power = the labour that is sold to the employer for a wage .
    \item If you are selling your labour power in exchange for a wage, then labour power is also a commodity . You sold your commodity and got \rupee30,000 for it --- exactly as you might have sold a car .
\end{itemize}

\paragraph{Every commodity has two values}
\begin{table}[h!]
\centering
\begin{tabularx}{\textwidth}{lYY}
\toprule
\rowcolor{AccentDark}
\thead{} & \thead{Use value} & \thead{Exchange value} \\
\midrule
What it is & The services rendered for yourself or for your own needs, not for anyone else & The same services sold to someone else for money \\
What it reflects & Fulfilment of human need per se & Commercial gain obtained by selling the commodity \\
Name of the labour & Labour & Labour power \\
\bottomrule
\end{tabularx}
\end{table}

\begin{ExampleBox}
\begin{itemize}
    \item I make two cups of noodles at home for myself and my brother . I opened the packet, boiled the water for three minutes, stood waiting, added the ingredients --- I gave roughly ten minutes of my life . That is labour, and it has use value: I used my labour for my own purpose and for my brother .
    \item My brother is so impressed that he opens a shop for me . I now make exactly the same noodles, but charge \rupee30 from every customer . The same labour, the moment it is performed for someone else, is charged for . The commodity now has exchange value .
\end{itemize}
\end{ExampleBox}

\begin{itemize}
    \item Terminology to fix: when you do labour only to fulfil your own needs, at home, without any commercial interest, that is regarded purely as labour . When you exchange the labour for money for commercial gain, that is known as labour power .
\end{itemize}

\begin{QuickRevision}
\begin{itemize}
    \item Labour power = labour sold to the employer for a wage; it is itself a commodity .
    \item Every commodity has use value and exchange value .
    \item Use value = fulfilment of human need; exchange value = commercial gain .
    \item Same act, different value: noodles at home vs noodles at the shop .
\end{itemize}
\end{QuickRevision}

\subsection{The Source of Profit}

\paragraph{The question Marx set himself}
\begin{itemize}
    \item What is the source of profit to the capitalist class? A very basic research question in his mind .
    \item Why did he want the answer? Because Marx read Adam Smith and David Ricardo --- the two most popular economists of those times --- and also Malthus and Say (Say's law) . These four were dealing with how the economy runs and what the source of profit to the capitalist class is .
    \item Their position: the price of a commodity --- how a pencil comes to cost \rupee10 --- is purely based on supply and demand, that is, on market forces .
    \item Marx's counter-question: if everything is purely decided by market forces, not by any individual, if it is completely in the hands of supply and demand, then how is it that despite this natural force, some people are able to appropriate profit and some are not?  How is it that some become rich and some become poor?  Is that also justified by the law of supply and demand --- or is there something hidden? 
    \item Marx believed something is hidden --- there are many issues no other economist had dealt with .
\end{itemize}

\begin{KeyConcept}
The source of profit is the exchange value of labour power . Remember this statement forever .
\begin{itemize}
    \item And the line behind it: labour has the power to create values more than the worth of itself --- labour has the potential to create values more than its own worth .
\end{itemize}
\end{KeyConcept}

\paragraph{Worked illustration}
\begin{itemize}
    \item You sold your commodity --- labour power --- to the company for \rupee30,000 per month . But the company extracted values out of your labour worth \rupee1 lakh .
    \item Take the content-developer case: salary \rupee30,000 . The content was sold on the internet, in bookshops, to distributors and libraries, and garnered value worth \rupee5 lakh .
    \begin{itemize}
        \item $\text{Profit} = \rupee5,00,000 - \rupee30,000 = \rupee4.7\text{ lakh}$ .
        \item The source of this profit is obviously the man who was given \rupee30,000 . If he was not there, you would not have had this much profit .
    \end{itemize}
    \item Therefore: if you have a lower exchange value, you will have higher returns . If you do not give enough wage or increment to the labourer, you will have a lot of profit in your hand .
\end{itemize}

\paragraph{What exploitation means}
\begin{itemize}
    \item Exploitation is not beating the workers . Remember this .
    \item Exploitation is the result of the difference between what the worker has got and what the worker has created . Created: \rupee5 lakh . Received: \rupee30,000 . Difference: \rupee4.7 lakh .
    \item And this explains the rate of exploitation: the bigger the difference, the higher the rate of exploitation . If the creation were worth \rupee10 lakh, the difference would be \rupee9.7 lakh --- showing that the person who got \rupee30,000 is the most exploited man in that company .
    \item The degree of exploitation is easily calculable: what you are getting versus what you are producing . What you offer the company is far more worthy than the wage you get in return; the two are not in proportion to each other .
\end{itemize}

\begin{CriticalPoint}
\begin{itemize}
    \item Marxian economics is generally not taught in universities --- specifically in highly right-wing universities --- for the reason that it may develop sentiments the state would not want among students . In some departments it is taught [as stated in lecture] .
    \item The details of Marxian economics are not important for us . What is important is the question Marx had in his mind .
\end{itemize}
\end{CriticalPoint}

\begin{QuickRevision}
\begin{itemize}
    \item Marx's research question: what is the source of profit to the capitalist class? 
    \item Classical economists (Smith, Ricardo, Malthus, Say) answered: supply and demand . Marx asked why, then, some appropriate profit and some do not .
    \item Answer: the source of profit is the exchange value of labour power .
    \item Because labour creates value greater than its own worth .
    \item Exploitation = $(\text{value created}) - (\text{value received})$; rate of exploitation rises with that gap .
\end{itemize}
\end{QuickRevision}

\subsection{The Working Day}
\begin{itemize}
    \item To go into the micro details, take the wage per hour: \rupee100 per hour, eight hours a day (9 a.m. to 5 p.m.) = \rupee800 .
    \item Marx says: what you wanted to earn for your own needs or your own sustenance was already met in the first four hours .
    \item What are you doing in the next four hours? You are working for the person who employed you, so that you could give him profit .
    \item The working day has been divided: some amount of work you do to meet your own needs; some amount of work you do to meet the needs of the employer .
    \item And the wage you get is only for the work that you did for your own needs . You are not going to get a wage for the work that you did to meet the needs of the employer .
\end{itemize}

\begin{ExampleBox}
\begin{itemize}
    \item The variant worked in the doubt session: \rupee2,000 per day, eight hours of work . Marx would say the \rupee2,000-worth of labour was already produced in the first five hours .
    \item In the extra three hours the capitalist extracts the revenue or profit . You have done the whole work in the first four or five hours; the extra hours are for the capitalist class, and you get no reward for them --- your reward is fixed at \rupee2,000 .
\end{itemize}
\end{ExampleBox}

\begin{QuickRevision}
\begin{itemize}
    \item Working day = necessary labour time + surplus labour time .
    \item First four (or five) hours: the worker earns his own subsistence . The remaining hours: the worker produces for the employer, unpaid .
    \item The wage covers only the first portion .
\end{itemize}
\end{QuickRevision}

\subsection{Value of Commodity = C + V + S}
\begin{itemize}
    \item Marx believes the value of a commodity is equal to $C + V + S$ .
    \item $S = \text{surplus value}$. $V = \text{variable capital}$. $C = \text{constant capital}$ .
    \begin{itemize}
        \item There are two types of commodity in play here: labour power is a commodity (what I have offered to the company in exchange for a wage), and a smartphone is a commodity . The worth of the phone and the worth of the labour power are both $C + V + S$ .
    \end{itemize}
\end{itemize}

\begin{table}[h!]
\centering
\begin{tabularx}{\textwidth}{lYY}
\toprule
\rowcolor{AccentDark}
\thead{Term} & \thead{Definition} & \thead{What it consists of} \\
\midrule
C --- Constant capital & The capital required to produce the commodity. ``Constant'' means it is not changing, not dynamic --- it is static. & Machines and tools --- the machine out of which shoes or phones are produced \\
V --- Variable capital & The labour power. Not machines and tools, but you --- the labourer who gave his or her labour power to the company in exchange for money. & The worker's labour power \\
S --- Surplus value & The higher chunk that the capitalist class appropriated --- the difference between the value of wage and the value of capital. & \rupee4.7 lakh in the running example; \rupee9.7 lakh in the larger one \\
\bottomrule
\end{tabularx}
\end{table}

\begin{itemize}
    \item What was the source of the \rupee5 lakh? It was that labourer only, who got \rupee30,000 . If that labourer was not there, there would not have been \rupee5 lakh or \rupee3 lakh of income; there would not have been any surplus .
    \item Human labour is the source of capital .
\end{itemize}

\paragraph{The most important line here}
\begin{KeyConcept}
Constant capital does not add any value to the commodity . Rather, it is the variable capital that adds all the value to the commodity . That is: machines do not add any value; labour power is the source of all values .
\end{KeyConcept}

\begin{itemize}
    \item Illustration: do you really think that the computer, the word processor, the paint program or the design software gave any value to my content?  The content was all mine . I am the source of profit --- not the software .
    \begin{itemize}
        \item That software is available to everyone . How you use it is what makes you different and makes your content different from the rest . So if that content garnered \rupee5 lakh, it is not because of the word processor or the design software --- it is because of the labour, the variable capital, and that alone .
    \end{itemize}
\end{itemize}

\begin{QuickRevision}
\begin{itemize}
    \item $\text{Value of commodity} = C + V + S$ .
    \item $C = \text{constant capital} = \text{machines and tools}$, static, already invested .
    \item $V = \text{variable capital} = \text{labour power}$, the worker .
    \item $S = \text{surplus value} = \text{the appropriated difference between value created and wage paid}$ .
    \item Constant capital adds no value; variable capital adds all the value . Human labour is the source of capital .
\end{itemize}
\end{QuickRevision}

\subsection{Capitalism as an Inverted World}
\begin{itemize}
    \item Exchange value of labour is the source of profit . The exchange value was very low --- \rupee30,000 --- while the worth of the services lent to the company was \rupee1 lakh . So the surplus is \rupee70,000 .
    \item Who is going to have that \rupee70,000 --- you, or the person who employed you?  The person who employed you keeps it in his pocket . But who created it? You did . You created \rupee70,000 worth of product and did not get \rupee70,000; the capitalist got it .
    \item Is capitalism not an inverted world? This is what Karl Marx says: capitalism is an inverted world . You do the work but do not get the amount; he does not do the work but gets all .
\end{itemize}

\paragraph{Homo Faber}
\begin{itemize}
    \item In primitive communism, in ancient society and in feudal society, what was recognised was Homo Faber --- man as a working, creative, expressive being .
    \item It is only through work that you express your creativity and your potential . But today, in capitalism, the one who works does not get the money and the one who does not work gets all the money .
    \item Capitalism runs purely and purely on profit . Nothing else is the driving force of this engine --- only profit, only surplus value .
\end{itemize}

\paragraph{Living labour and dead capital}
\begin{KeyConcept}
Marx's line: living labour creates dead capital, but the dead capital dominates over the living capital .
\end{KeyConcept}

\begin{itemize}
    \item What is dead capital? The content . It is dead in the sense that it is not a living species . The source of the content was you, who created it .
    \item But this dead capital has come to dominate over the living capital, because the living capital was worth only \rupee30,000 while the dead capital is worth \rupee5 lakh .
    \item Your creation has become over and above and external to you in this world . Your creation has assumed an identity above you --- you are worth \rupee30,000, your creation is worth \rupee5 lakh . This is the inverted world we live in .
    \begin{itemize}
        \item You already knew these things before joining the course, but somehow you did not think like this and did not go into these details . What seemed to be common sense --- how did it become common sense to us? That is what Marx is arguing .
    \end{itemize}
\end{itemize}

\begin{QuickRevision}
\begin{itemize}
    \item Surplus is created by the worker and pocketed by the capitalist --- capitalism is an inverted world .
    \item Homo Faber: man as a working, creative, expressive being --- recognised in earlier modes, denied in capitalism .
    \item Capitalism's only driving force is profit / surplus value .
    \item ``Living labour creates dead capital, but dead capital dominates over living capital'' .
    \item The creation assumes a reality above and external to its creator .
\end{itemize}
\end{QuickRevision}

\subsection{Doubts Taken in Class}

\paragraph{State the definition of surplus value clearly}
\begin{itemize}
    \item The value your work produces and what you get are two different things --- the potential you have versus what you get . The difference between these is the surplus value, appropriated by the capitalist class .
\end{itemize}

\paragraph{What do the extra working hours have to do with it?}
\begin{itemize}
    \item Assume \rupee2,000 per day for eight hours --- approximately \rupee50,000--\rupee60,000 per month allowing for holidays . \rupee2,000 is fixed income per day, but you must work eight hours .
    \item Marx says the \rupee2,000-worth of labour was already produced in the first four or five hours . In the extra three hours the capitalist extracts the revenue or profit from you --- that is what gives him the higher return, and you get no reward for it because your reward is fixed .
    \begin{itemize}
        \item Coming next: the modes of exploitation --- the concept of the length of the working day, absolute surplus value and relative surplus value .
    \end{itemize}
\end{itemize}

\paragraph{Is Durkheim's division of labour a source of solidarity a law?}
\begin{itemize}
    \item No, it is not a law . Durkheim was not a positivist sociologist when he discussed the division of labour . It is just his belief; it is not something that can be empirically tested, and he never said this was a law he was giving .
    \item One thing he did state law-like: with industrialisation the older form of solidarity weakens . But the claim that organic solidarity will then integrate society is not a law [as stated in lecture] .
    \item The reason: when Durkheim discusses the division of labour thesis he is always comparing society with an organism, and in an organism everything happens naturally . Since it was a natural, evolving outcome, he did not decide to give a law .
\end{itemize}

\paragraph{If firms underpay us and the government also takes a portion, are we doubly exploited?}
\begin{itemize}
    \item The student's formulation was that we have become new slaves and the government owns our labour . [Recorded as raised in class; the detailed answer was carried into the state-and-capitalism discussion of the later classes.] 
\end{itemize}

\paragraph{Is the value of wage in surplus value the same as variable capital?}
\begin{itemize}
    \item Variable capital is you . Constant capital is the machines from which the goods come --- but who created those machines? Labour power only . You created those machines and you work on them .
    \item Why are you called ``variable''? Because you are adding value to the commodities . The machine per se does not add any value . Which design of shoe is to be made --- do machines decide that? You do . You control the machines, you repair them, you make the whole chart, you decide sales and marketing . You are the main source of value to the capitalist, not the machines .
\end{itemize}

\paragraph{In this era of heavy industrialisation, is it right to assume constant capital is negligible?}
\begin{itemize}
    \item No . Constant capital is constant only --- in comparison with variable capital, which is the labour I am selling in exchange for income . Constant capital is something already invested in by the capitalist class .
    \item The machine, the assembly line, the transformer --- these are static . They are not going to give value to the commodity; the value is given by you . In fact those machines were created in the first place by your labour, so you are the ultimate source of profit, not the machines .
\end{itemize}

\paragraph{One more numerical}
\begin{itemize}
    \item You worked for the company and made a product whose worth in the market is \rupee2 lakh, but you were given \rupee40,000 . Surplus value = \rupee1.6 lakh, appropriated by the company, not by you .
    \item Tomorrow you will get a bonus, or after one year you will get a raise, and you will be happy . Why are you happy? This is false consciousness --- taken up in the next class .
\end{itemize}

\begin{QuickRevision}
\begin{itemize}
    \item Surplus value = $\text{value produced} - \text{value received}$; appropriated by the capitalist .
    \item The extra hours of the fixed working day are the source of the capitalist's return .
    \item Durkheim on division of labour: not a law --- society compared with an organism, so change is natural rather than law-governed .
    \item You are variable capital because you add value; machines are constant capital and were themselves made by labour .
    \item Constant capital is never negligible, but it is never value-creating either .
    \item Being pleased by a bonus while the surplus stays appropriated = false consciousness .
\end{itemize}
\end{QuickRevision}

\sectiondivider

\section{Surplus Value, the Law of Motion of Capitalism and Class Struggle}

\subsection{What the Test Was Actually For}
\begin{itemize}
    \item On the answers glanced through: the attempts were good . It does not matter how many questions you have left --- initially it will take time to attempt all the questions .
    \item Those who left four questions need to work on time management: the paper was not so lengthy that four questions should be left . Leaving four questions means something is wrong with your writing speed, or with your thinking process --- you are taking a lot of time to think .
    \item Those who left one or two questions are completely fine . Over time, those leaving four will come down to one or two, and those leaving one or two will complete the paper . This will happen naturally; we do not need to worry about it .
\end{itemize}

\paragraph{Why the test used new questions}
\begin{itemize}
    \item Since previous year questions had already been gone through, the main purpose of the test was to check how you manage everything when you confront new questions .
    \item In the examination you will get new questions only . We should not expect them to repeat questions --- at least for the last three years that is not happening in any optional .
    \item They are not going to see what you have memorised; what you have understood is what they are going to check .
\end{itemize}

\begin{QuickRevision}
\begin{itemize}
    \item Leaving four questions = a speed or thinking-process problem; leaving one or two is normal at this stage .
    \item Tests are set with new questions because the examination itself does not repeat questions .
    \item Understanding, not memorisation, is what is being examined .
\end{itemize}
\end{QuickRevision}

\subsection{Everything Has a Price Tag --- Use Value Turned into Exchange Value}
The whole of yesterday's micro-detail can be simplified into one movement . Take the same act twice, once at home and once in the market .

\begin{table}[h!]
\centering
\begin{tabularx}{\textwidth}{lYY}
\toprule
\rowcolor{AccentDark}
\thead{The same thing\dots} & \thead{\dots as use value} & \thead{\dots as exchange value} \\
\midrule
Your smile & Greeting guests at home --- ``hello uncle, hello auntie, how are you?'' The need being satisfied is respect for elders and guests. & Smiling in a shopping mall greeting customers --- ``hello sir, how may I help you?'' The customer is impressed by the service and buys; your smile has generated revenue for the company. \\
Noodles you cook & Making them for your brother or sister at home. & Making them for paying customers. \\
Your sense of humour & Recalling old stories and cracking jokes so your friends laugh. & Making an audience laugh. \\
Your car & You decide where to go and where to stop over. & You rent it out to a ride-hailing platform. \\
Your house & You sit in your room, study, talk to friends. & You rent it out to students or list it on a home-sharing platform. \\
\bottomrule
\end{tabularx}
\end{table}

\begin{itemize}
    \item What does it mean? It means only one thing: today everything has got a price tag . Your smile, your car, your house, your sense of humour, even your beauty, even your eyes . Even your body organs have value .
\end{itemize}

\begin{KeyConcept}
\begin{itemize}
    \item Marx: all that is solid melts into air, and all that is holy is profaned . This is what capitalism --- modern industrial capitalism --- is .
    \item This is not the world people once had, when they were hunting and gathering and simply consuming their primary needs .
    \item Capitalism is that spirit which can turn all use values into exchange values --- and this is how profit is generated .
\end{itemize}
\end{KeyConcept}

\begin{itemize}
    \item So, in simple words: in capitalism the use values are transformed into exchange values, and this generates profit .
    \item But that does not fully answer the question . The next question is: how does exchange value generate profit? 
\end{itemize}

\begin{QuickRevision}
\begin{itemize}
    \item The same activity has use value at home and exchange value in the market .
    \item Everything today carries a price tag --- smile, humour, car, house, beauty, even body organs .
    \item ``All that is solid melts into air, and all that is holy is profaned'' .
    \item Capitalism is the spirit that converts use values into exchange values; that conversion is the generation of profit .
\end{itemize}
\end{QuickRevision}

\subsection{What Counts as Labour, and What Counts as a Commodity}
\begin{itemize}
    \item When you made your friends laugh, was that not labour? What is the simple meaning of labour? Work .
    \item Making other people laugh is also work . Sometimes you do not charge for it --- that does not mean you are not doing work; it just is not paid work .
    \item When the same person made his friends laugh we did not call it work; now he makes an audience laugh and it is regarded as work . This is weird --- and that is the point .
    \item Marx: every commodity has both use value and exchange value .
\end{itemize}

\paragraph{Commodity does not mean physical object}
\begin{itemize}
    \item My communication skill is also a commodity . This is why companies ask for someone with good communication skill and good experience .
    \item Experience is also a commodity . Your educational background is also a commodity --- only because of these things do you get money .
    \item You exchange your past experience, your educational background, your communication skill with the other person, and in exchange you get money . You have exchanged a commodity --- and the name of that commodity is labour .
    \item Everything in the world is a commodity .
\end{itemize}

\begin{QuickRevision}
\begin{itemize}
    \item Labour simply means work --- including unpaid work such as entertaining friends .
    \item Every commodity has use value and exchange value .
    \item A commodity need not be a physical object: communication skill, experience and educational background are all commodities .
    \item What you actually exchange for a wage is the commodity called labour .
\end{itemize}
\end{QuickRevision}

\subsection{How is Exchange Value Determined?}
\begin{itemize}
    \item If I am a comedian who has to make people laugh, how is my exchange value going to be determined? What is the criterion for paying me? A very big question .
    \begin{itemize}
        \item Imagine a comedian whose video garners one billion views --- that means his worth is one billion views . But he is paid only for one million views . This generates profit .
        \item Noodles version: your skill can make a plate of noodles worth \rupee1,000, and the price tag on your noodles is \rupee900 or \rupee1,000 . But how much are you getting? \rupee100 per plate . Where has the \rupee900 gone? 
    \end{itemize}
\end{itemize}

\begin{KeyConcept}
Marx's idea is very simple: the worker in the production process creates value more than what is given to him or her . And that created value which is more than what is given to him --- is that not surplus? This is surplus value .
\end{KeyConcept}

\begin{QuickRevision}
\begin{itemize}
    \item Exchange value is set below the value the worker actually creates .
    \item The comedian paid for one million of a billion views; the noodle-maker paid \rupee100 on a \rupee1,000 plate .
    \item The gap is surplus value: the worker creates more value than what is returned to him .
\end{itemize}
\end{QuickRevision}

\subsection{Value of Commodity = C + V + S --- Worked Through a Smartphone}
\begin{itemize}
    \item Take the commodity as a smartphone . In common sense its value is, say, \rupee80,000 . How did \rupee80,000 become the price tag? Because the commodity carries many things which, when added, make the cost \rupee80,000 .
    \item Marx says those things are $C + V + S$ .
\end{itemize}

\paragraph{Definitions dictated in class}
\begin{itemize}
    \item Constant capital ($C$) is the amount of money invested in raw material, machines or tools .
    \item Variable capital ($V$) is the amount of money invested in hiring the labourer --- the amount of money invested on the wage of the labour, which we call labour power .
    \begin{itemize}
        \item If my worth is \rupee10 per month, that means the employer has invested \rupee10 on me . So variable capital is \rupee10 . My worth is \rupee10 . I have sold my labour for \rupee10 --- that is why labour power is variable capital .
    \end{itemize}
    \item Surplus ($S$) --- and remember, where did this surplus come from? It did not come from thin air . The surplus comes from the labourers only . Is the smartphone itself creating surplus, or is it the labourer who creates it? It is created by $V$ .
\end{itemize}

\paragraph{The chain of reasoning}
\begin{enumerate}
    \item My worth is \rupee10 . But what I am creating is worth more than \rupee10 . And what I am getting is only \rupee10 .
    \item Marx: the worker produces, in the work process, value worth $V + S$ .
    \item $V$ is variable capital --- the labour power, which is paid labour .
    \item $S$ is the surplus created by that labour power, worth more than the labour power itself . But $S$ is not given to the worker --- therefore $S$ is unpaid labour .
\end{enumerate}

\begin{CriticalPoint}
This is what exploitation is: you do not give him the unpaid labour; you have hidden the unpaid labour from him . You give him only what is required to meet his needs, and put the rest in your pockets .
\end{CriticalPoint}

\paragraph{Rate of exploitation}
\begin{itemize}
    \item Marx: the rate of exploitation = $\text{surplus value} \div \text{variable capital} = S / V$ .
    \item The more the $S$ --- the more surplus value is generated --- the higher the rate of exploitation .
    \item It also means: $\text{rate of exploitation} = \text{unpaid labour} \div \text{paid labour}$ .
\end{itemize}

\begin{QuickRevision}
\begin{itemize}
    \item $\text{Value of commodity} = C + V + S$ .
    \item $C = \text{money invested in raw material, machines and tools}$; $V = \text{money invested in hiring the labourer (the wage / labour power)}$; $S = \text{surplus, created by } V \text{ alone}$ .
    \item The worker produces $V + S$ but is paid only $V$ --- $S$ is unpaid labour .
    \item $\text{Rate of exploitation} = S / V = \text{unpaid labour} / \text{paid labour}$ .
\end{itemize}
\end{QuickRevision}

\subsection{Why Capitalism Gets the Detailed Treatment}
\begin{itemize}
    \item Notice that we did not go into such detail for feudalism, for ancient society or for primitive communism . Why? Because Marx himself never went into those details of those modes of production .
    \item Marx is the scholar of capitalism . He is always remembered as the scholar of capitalism --- not even of communism . Communism is a dream .
    \item Largely, whatever he has written is about capitalism --- thus \emph{Capital}, volumes 1, 2 and 3 (\emph{Das Kapital}) . He is mainly talking about capital .
\end{itemize}

\begin{QuickRevision}
\begin{itemize}
    \item Marx did not elaborate the earlier modes; he is the scholar of capitalism .
    \item Communism in his work is a dream, not an analysis; \emph{Capital}, volumes 1--3 is the body of the analysis .
\end{itemize}
\end{QuickRevision}

\subsection{The Two Modes of Exploitation}
\begin{itemize}
    \item We know the rate of exploitation is $S/V$ . The more you exploit, the more profit you gain; the more surplus value generated, the more profit . Now: how does exploitation actually take place? 
    \item There are two modes of exploitation: absolute surplus value and relative surplus value .
\end{itemize}

\paragraph{Absolute surplus value}
\begin{itemize}
    \item Take the length of the working day as 0 to 8 hours .
    \item Within that 0--8 span, the worker was getting only that amount required for him to meet his needs --- equivalent to 3 hours ($V$) . The remaining 3 to 8, that is 5 hours, was surplus generated for the capitalist class .
\end{itemize}

\begin{ExampleBox}
\begin{itemize}
    \item The services you offer the company are worth \rupee1 lakh per day . What are you getting per day? \rupee2,000 --- the bare minimum that should be given to you to meet your everyday needs . The remaining \rupee98,000 is retained .
    \item And that \rupee2,000 is made in about three hours . What are you doing for the remaining five hours? You are working for the capitalist class, not for yourself .
\end{itemize}
\end{ExampleBox}

\begin{itemize}
    \item If I am the industrialist and want to extract maximum profit, the first option is to increase the length of the working day . In 8 hours the workers produced \rupee1 crore; so double the hours to 16 and they can produce \rupee2 crore . Make it 10 hours .
\end{itemize}

\paragraph{Why this route is blocked}
\begin{itemize}
    \item This is not an easy job --- especially today, when you have labour laws, the Factory Act, a welfare state, and communist parties across states who will protest that making workers work 12 or 14 hours is inhuman and against human rights .
\end{itemize}

\begin{GovtPolicy}
\begin{itemize}
    \item During the lockdown, labour laws were suspended . Suspension of labour laws means the capitalist class was free to make workers work extended hours, and workers could not form unions or protest --- they had to work silently .
    \item Minimum wage was also withdrawn; maternity benefit was also withdrawn . This happened in Madhya Pradesh, Maharashtra and Uttar Pradesh, and a case also went to the Supreme Court .
    \item That was a mode of exploitation: minimum wage not guaranteed, and increased working hours definitely there .
\end{itemize}
\end{GovtPolicy}

\paragraph{The second route to absolute surplus value}
\begin{itemize}
    \item Surveillance --- tight monitoring of the workers . If they want a half-hour break, it should be strictly half an hour; managers are kept for exactly this, because every single minute counts .
    \item This is why companies are very strict and now have biometric mechanisms: come on time, leave on time and not before --- because it is important for them to generate surplus value .
    \item Something we do not realise: we feel we are working for ourselves . Marx says no --- you are working for yourself only about 20\%; for the remaining 80\% you are working for the company .
\end{itemize}

\paragraph{The lines to write}
\begin{itemize}
    \item Absolute surplus value is generated by increasing the length of the working day, and by tight monitoring or surveillance of the workers .
    \item However, the length of the working day cannot be increased beyond a point, due to labour laws and human rights .
\end{itemize}

\paragraph{Relative surplus value}
\begin{itemize}
    \item Since the first option is blocked, the last option left is relative surplus value --- the second way to exploit .
    \item Recall the formula $S/V$ . If you want the rate of exploitation in your favour, either increase $S$ or decrease $V$ . Extending the working day is what increases $S$ directly; that cannot be done . But decreasing $V$ can be done .
    \item So: 0 to 8 cannot be extended . Till 3 hours it was $V$, and the remaining 5 was surplus . Reduce the 3 to 2 --- and the surplus rises from 5 to 6 . That is relative surplus value .
\end{itemize}

\begin{table}[h!]
\centering
\begin{tabularx}{\textwidth}{lYY}
\toprule
\rowcolor{AccentDark}
\thead{} & \thead{Absolute surplus value} & \thead{Relative surplus value} \\
\midrule
Method & Increase the length of the working day; plus tight monitoring and surveillance of workers & Reduce the necessary labour time by increasing the productivity of labour and introducing machines \\
Effect on the 8-hour day & Stretch 8 hours to 10, 12, 16 & Keep 8 hours; push V down from 3 hours to 2, so S rises from 5 to 6 \\
Constraint & Blocked by labour laws, the Factory Act, the welfare state, unions and human rights & Not blocked --- but it is what eventually destroys capitalism \\
\bottomrule
\end{tabularx}
\end{table}

\paragraph{How relative surplus value is achieved in reality}
\begin{itemize}
    \item The task for the capitalist is to increase the productivity of labour . What he produced in 3 hours, if he can now produce in 2 hours, productivity has increased .
    \item And it does not stop: what you did in 2 hours you now do in 1; then in half an hour; then 15 minutes, 10, 5, 4, 3, 2, 1 --- and then you are unemployed .
    \item The increase in productivity is supplemented with the introduction of machines . Most workers are now gradually replaced by machines, precisely because the length of the working day cannot be increased .
    \begin{itemize}
        \item You will abuse the capitalist class --- they make us work all day, they do not even leave us alone at home . But machines cannot abuse the capitalist class .
    \end{itemize}
    \item Line to write: relative surplus value is achieved by increasing the productivity of the labour and the introduction of machines .
\end{itemize}

\paragraph{What machines then do}
\begin{itemize}
    \item Machines are highly productive; they will produce 24$\times$7 without any complaint .
    \begin{itemize}
        \item And you have now been appointed manager of the company . You are happy about the promotion --- but what does it mean? You now manage: that the machines are working properly, that staff comes on time, the repairs of the machines, the installation costs . You are doing monitoring now .
    \end{itemize}
    \item The fact is: machines have been introduced at the cost of workers .
    \item Although machines increase the amount of goods produced, since workers have been eliminated those goods will not be purchased . This leads to overproduction and under-consumption .
    \begin{itemize}
        \item In economy there are cycles of boom and bust --- sometimes overproduction and under-consumption, sometimes underproduction and over-consumption . Marx is explaining the reason behind that cycle --- what is happening in reality at the ground level .
    \end{itemize}
    \item However, this will sow the seeds of the destruction of capitalism .
\end{itemize}

\begin{DataFact}
\begin{itemize}
    \item From about 1850 to 2020 --- more than 150 years --- what we have seen is nothing but more and more mechanisation and rationalisation at workplaces .
    \item Old workers and unskilled workers are getting more and more unemployed .
    \item Marx's thesis: the increase in unemployment is inevitable --- it cannot be avoided .
    \item Surplus value is generated at the cost of the employment of millions: capital is accumulated and appropriated by a few people, and the maximum number are unemployed .
\end{itemize}
\end{DataFact}

\begin{QuickRevision}
\begin{itemize}
    \item Two modes of exploitation: absolute surplus value and relative surplus value .
    \item Absolute = lengthening the working day + surveillance; blocked by labour laws and human rights (except when those laws are suspended, as in the lockdown) .
    \item Relative = shrinking necessary labour time by raising productivity and introducing machines .
    \item $S/V$: absolute raises $S$ directly; relative lowers $V$ .
    \item Machines replace workers $\rightarrow$ overproduction and under-consumption $\rightarrow$ boom and bust; rising unemployment is inevitable .
\end{itemize}
\end{QuickRevision}

\subsection{The Industrial Reserve Army}
\begin{itemize}
    \item When machines are introduced and many people are replaced by them, the pool of these unemployed people is called the industrial reserve army (IRA) .
    \item Definition to write: the industrial reserve army consists of a pool of both the unemployed who want to work, and the underemployed .
    \begin{itemize}
        \item If you are underemployed, you are a member of the reserve army . If you are unemployed but want to work, you are a member of the reserve army .
        \item What is underemployment? Suppose you have a PhD and you are working in a menial job --- wiping floors in some office . That is underemployment . And for some positions people are over-employed .
    \end{itemize}
\end{itemize}

\paragraph{What the pool does to bargaining power}
\begin{itemize}
    \item The capitalist class now has many people who can work at whatever price the capitalist decides --- because they are unemployed . Readily available for exploitation .
    \item If you are the owner of a factory and a worker asks for a high wage, you can simply say: we have many people who will work for a hundred rupees; why should we give you more? So when there are many unemployed people, you lose bargaining power .
    \item The capitalist class is free to recruit any person from that pool --- and this pool is the creation of the capitalist class itself .
\end{itemize}

\begin{ExampleBox}
\begin{itemize}
    \item When the lockdown was imposed, workers said in interviews that they could do whatever it took to sustain their lives today, because they had nothing in their hands .
    \item This is what the reverse-migrant workers were saying while going back to their homes . They had become extremely vulnerable --- and that is the extent of vulnerability those people had fallen victim to .
    \item These people are, in the language of Marx, the industrial reserve army .
\end{itemize}
\end{ExampleBox}

\begin{QuickRevision}
\begin{itemize}
    \item Industrial reserve army = pool of the unemployed who want work + the underemployed .
    \item Underemployment: qualification far above the work performed .
    \item The pool destroys the bargaining power of those in work, and it is created by the capitalist class itself .
\end{itemize}
\end{QuickRevision}

\subsection{Dead Labour Dominates Living Labour}
\begin{itemize}
    \item As the cycle goes on --- necessary labour time falling from 2 to 1.8 to 1.6 to 1.4 to 1 to 0.5 --- you will lose your labour power . And when you go looking for a job, the capitalist can pull someone cheaper than you from the large pool of the underemployed . This is how capitalism is running .
\end{itemize}

\paragraph{The conclusion}
\begin{itemize}
    \item The worker produces values worth more than its own value at the workplace . The value of the commodity generated by you is worth more than your own self .
    \item And what you are producing is dead labour --- it is simply an expression of your labour, so it is dead .
    \item Dead labour is worth more than living labour . Dead labour dominates over living labour . Your worth is \rupee30,000 per month; what you are producing is worth \rupee1 lakh per day .
\end{itemize}

\begin{ExampleBox}
\begin{itemize}
    \item You are a painter . A company has hired you and gives you \rupee30,000 per month for making paintings, and those paintings are sold for \rupee11 crore .
    \item Your work is bigger than you . And this will happen only in capitalism .
\end{itemize}
\end{ExampleBox}

\begin{itemize}
    \item Since the commodity is worth more than the labour power, the dead capital dominates over the living capital, and --- the most important line --- the commodity assumes a reality external to the creator .
\end{itemize}

\begin{QuickRevision}
\begin{itemize}
    \item The worker produces more value than he himself is worth .
    \item What is produced is dead labour; dead labour dominates over living labour .
    \item The commodity assumes a reality external to its creator --- this happens only in capitalism .
\end{itemize}
\end{QuickRevision}

\subsection{Commodity Fetishism}
\begin{itemize}
    \item When you look at commodities like a premium coffee chain's coffee or a flagship phone, what is your first reaction? 
    \item The fact is that actors fail to realise that the worth of the commodity --- \rupee400 per glass of coffee, \rupee80,000 for the phone --- is actually the result of labour .
\end{itemize}

\paragraph{The Feuerbach parallel}
\begin{itemize}
    \item This is the same thing Feuerbach had told people: we have created God, but we have elevated God to such a level that the creation has become external to the creator --- and now we assume he is almighty, he is everything .
    \item And the same status we have now started granting to these objects . We fail to realise that these objects are the creation of their creators, and the objects assume a reality of their own .
    \item This is what is called the fetish of commodity . Line to write: this leads to fetishism of commodities .
\end{itemize}

\paragraph{What ``fetishism'' means}
\begin{itemize}
    \item Deep desire, revering, craving, lust --- ``I have a fetish for this, I have a fetish for that'' . It is almost equivalent to lust .
    \item And fetishism is, by the way, the name of a religion . It is about craving for something mystical and magical; those objects seem to create an aura, and you start believing they have talismanic features --- like pearls, or a highly precious stone that gives the impression that there is something magical inside it .
    \item Fetishism means to worship objects which seem highly magical to you . And the same worship we are doing of commodities .
\end{itemize}

\paragraph{The extension by the neo-Marxists}
\begin{itemize}
    \item We do this with actors and singers too . Adorno and Horkheimer call this the cultural industry; that is why they are neo-Marxists, and they were very much influenced by the writings of Marx .
    \item Reading Marx on commodity fetishism, they argued that we are not able to differentiate between the singers and performers to whom we have given so much value --- we start believing they have magical value, and they become superhero-like .
    \item They assume a creation just as God was to those people, and as the commodity is to the working class . The same thing happens when people see cultural personalities .
\end{itemize}

\paragraph{The line to write and the underlying argument}
\begin{KeyConcept}
Commodity fetishism is the result of actors failing to realise that the commodity is the creation of labour --- the creation of labour value only .
\end{KeyConcept}

\begin{itemize}
    \item Basically Marx is saying: the worth of the phone should be the worth of the amount given to the people who manufactured it . If the worth of the person who made it is \rupee20,000, then the amount of the phone should also be \rupee20,000 . What about the extra \rupee60,000? 
    \item Why do we give so much status to it? Because it differentiates people --- it helps people differentiate themselves from the rest . It becomes a status symbol, because investing \rupee1 lakh or \rupee2 lakh in a phone is a status .
    \item But why did it assume that status? Because it cost \rupee1 or \rupee2 lakh --- while the real cost is \rupee20,000 . The moment you break into this illusion, this smoke screen, you will not have a fetish for this object .
\end{itemize}

\begin{QuickRevision}
\begin{itemize}
    \item Commodity fetishism = failing to see that the commodity is the creation of labour .
    \item The Feuerbach parallel: we create God, elevate the creation above the creator, then worship it .
    \item Fetishism is originally the name of a religion --- worship of objects that appear magical or talismanic .
    \item Adorno and Horkheimer extend it to the cultural industry and to cultural personalities .
    \item The status symbol rests on the price, and the price rests on appropriated surplus --- break the illusion and the fetish dissolves .
\end{itemize}
\end{QuickRevision}

\subsection{The Law of Motion of Capitalism}
\begin{itemize}
    \item Kepler gave the law of planetary motion; Marx has given the law of motion of capitalism . A very important law .
    \begin{itemize}
        \item An aside noted in class: Durkheim and Marx are very much alike here --- both are macro scholars, both give importance to structures rather than to individuals, and both are saying that something becomes external to the individual . And when anything is external to you, it always assumes a coercive power over you .
    \end{itemize}
\end{itemize}

\paragraph{Capitalism rewards efficiency --- and you cannot beat a machine}
\begin{itemize}
    \item Capitalism rewards efficiency: the more efficient you are, the more you are rewarded . But you cannot be more efficient than a machine --- so the capitalist will always prefer machines over you .
    \item ``Machine'' here does not literally mean a machine . It means anything that is beyond human labour --- a supercomputer, artificial intelligence . You cannot be more efficient than these .
    \item After a point of time there will come a time when machines will start thinking like human beings . Today machines do not think like humans, and that is the only thing which differentiates a machine from us .
    \begin{itemize}
        \item When machines start thinking like humans, that is singularity . What will happen to humans then? What will be the means of subsistence? Will there still be an economy where you sell your labour and get paid every month? What kind of mode of production will there be? Very futuristic --- but some people say we will have singularity by 2050 .
    \end{itemize}
\end{itemize}

\paragraph{Why capitalism digs its own grave}
\begin{itemize}
    \item Capitalists are their own grave diggers . Workers do not have to do anything --- structurally it is happening .
\end{itemize}
\begin{enumerate}
    \item To extract maximum profit from workers, the working day cannot be lengthened, because those laws exist .
    \item So either the capitalist cuts the wages of labour, or he replaces the labourers with machines .
    \item In both cases you have developed the means of production in order to replace the workers .
    \item The moment the means of production get developed, the surplus value decreases . Why? Because the source of surplus value is labour . Machines are not the source of surplus value .
    \item Imagine you have many machines and become highly efficient with a great volume of goods . But who will purchase these goods? The people who could have bought them have been replaced by machines and have no purchasing power .
    \item Unless people purchase the goods, you cannot make profit . If you cannot make profit, the surplus value is decreasing .
\end{enumerate}

\paragraph{The lines to write}
\begin{itemize}
    \item The law of motion reveals the inherent contradictions within capitalism .
    \item With the development of means of production, the rate of surplus value will decrease --- because the source of surplus value is $V$, not $C$ . $S$ was the result of $V$; $S$ was not the result of $C$ .
    \item The law of motion describes the mechanism of the inherent contradiction .
\end{itemize}

\begin{QuickRevision}
\begin{itemize}
    \item Kepler: law of planetary motion . Marx: law of motion of capitalism .
    \item Capitalism rewards efficiency, and machines --- including supercomputers and AI --- are always more efficient; hence singularity as the extreme case .
    \item Replacing workers destroys purchasing power, hence profit, hence surplus value .
    \item The law of motion reveals --- and describes the mechanism of --- the inherent contradictions within capitalism .
    \item Marx and Durkheim converge: both macro, both structural, both on externality becoming coercive .
\end{itemize}
\end{QuickRevision}

\subsection{The Three Laws Within the Law of Motion}

\paragraph{First --- the law of capital accumulation}
\begin{enumerate}
    \item Each capitalist tries to expand by increasing output (surplus value) and reducing costs .
    \item They might hire more workers to increase the output .
    \item But the wages of these hired workers will eat up the company's profits .
    \item Therefore, in order to keep the wages low, the capitalist introduces labour-saving machinery .
    \item But this will create unemployment and a reserve army of labour . Trapped from both sides .
    \item And these unemployed people will not have purchasing power to buy the goods produced from the machines .
    \item This will result in overproduction and under-consumption --- which for the capitalist is complete loss .
    \item Therefore, in the long run, the rate of surplus value will fall, and cyclical recession or depression will creep in .
\end{enumerate}

\paragraph{Second --- the law of concentration of capital}
\begin{itemize}
    \item This is the direct outcome of the first law .
    \item What happens generally when companies are not making profits? They are eaten up by a bigger conglomerate or a bigger company --- and here come mergers and acquisitions (M\&A) .
    \begin{itemize}
        \item Bigger companies start merging with the companies making losses . What happens with public sector banks when they make losses? The bigger bank keeps them with itself . When small finance banks are in losses, the bigger banks buy them . Divestment is the corresponding term in public sector undertakings .
    \end{itemize}
    \item Line to write: the falling rate of profit will force concentration of capital, as some capitalists will try to buy out or merge with small businesses .
    \item This also shows that smaller companies will go out of business .
\end{itemize}

\paragraph{Downward mobility and proletarianisation}
\begin{itemize}
    \item When smaller companies go out of business, the people who worked in them also go out of jobs --- not all of them will be integrated into the big companies, which have their own demands and expectations of workers .
    \item These unemployed people will witness downward mobility in their life . They are sinking in --- they are going to join the ranks of the proletariat .
    \begin{itemize}
        \item Some had achieved upward mobility --- managers, who are not just proletariat . Some had opened their own shops: petty bourgeois --- the small shopkeepers, the small dairy owners .
        \item When the economy witnesses depression or recession, these small shops are either closed down or eaten up by the big corporates . This is what happened during the lockdown also --- many shops closed down, and in the language of Marx all those people have sunk in, joining the ranks of the proletariat . Those once making lakhs are now making thousands, earning their subsistence, trying to recover from the slowdown . There are many millions of people like this in India right now .
    \end{itemize}
\end{itemize}

\begin{KeyConcept}
Raymond Aron (R-A-Y-M-O-N-D A-R-O-N) calls this process proletarianisation .
\end{KeyConcept}

\paragraph{Third --- the law of increasing misery}
\begin{enumerate}
    \item One company aims to achieve higher efficiency and introduces machines; workers go out of jobs .
    \item With increased machines, surplus value inevitably decreases .
    \item Then under-consumption and overproduction happen; eventually recession and depression come .
    \item This company is eaten up by a bigger company .
    \item All those who became unemployed or underemployed join the ranks of the proletariat --- the reserve army . That is proletarianisation .
    \item The more people join the reserve army, the clearer the division of society into two --- polarisation .
    \item Society will increasingly become polarised between a shrinking capitalist class and a massive proletariat that suffers worsening misery .
    \item Finally a point will come when this cannot continue, and revolution shall occur .
\end{enumerate}

\begin{CriticalPoint}
\begin{itemize}
    \item In all three laws you will see the elements of the laws of dialectics .
    \item In the third law --- increasing misery --- the element of the first law is clear: increasing misery accumulates quantitatively until a point comes at which revolution occurs . Exactly like water being heated and heated until it turns to gas .
\end{itemize}
\end{CriticalPoint}

\begin{QuickRevision}
\begin{itemize}
    \item Law of capital accumulation: expand output, cut costs $\rightarrow$ hire $\rightarrow$ wages eat profit $\rightarrow$ labour-saving machinery $\rightarrow$ unemployment + reserve army $\rightarrow$ no purchasing power $\rightarrow$ overproduction and under-consumption $\rightarrow$ falling rate of surplus value $\rightarrow$ cyclical recession/depression .
    \item Law of concentration of capital: falling profit forces buy-outs and mergers; small firms die; their workers and the petty bourgeois sink into the proletariat --- Raymond Aron's proletarianisation .
    \item Law of increasing misery: polarisation between a shrinking capitalist class and a massive, worsening proletariat, until revolution occurs .
    \item The first law of dialectics --- quantity into quality --- is visible in the third law .
\end{itemize}
\end{QuickRevision}

\subsection{Class Struggle --- Class in Itself and Class for Itself}\label{sec:class-in-for-itself}
\begin{itemize}
    \item Class struggle is the struggle between two classes: the bourgeoisie and the proletariat .
    \item Marx says class has two dimensions: objective and subjective .
    \begin{itemize}
        \item Subjective is class in itself . The objective dimension of class is class for itself . [Note: in the previous class the same pair was presented with the labels the other way round --- see section 2.13 . Both formulations are recorded as delivered.] 
    \end{itemize}
\end{itemize}

\paragraph{The subjective dimension --- class in itself}
\begin{itemize}
    \item Marx says a class consists of people having similar relation to the means of production . That is the subjective dimension of class .
    \item All of you are proletariat . We do not own the means of production . We only have our labour to sell to the corporate people .
    \begin{itemize}
        \item It does not matter whether you are a worker in a large software firm or a worker at a street food stall in a busy market --- all of you are similar; you are a worker, proletariat, at the end of the day .
    \end{itemize}
\end{itemize}

\paragraph{The objective dimension --- class for itself}
\begin{itemize}
    \item The objective dimension of class is when this class in itself realises that whether we are in the software firm or at the street stall, both of us are exploited --- and our common enemy is the owning class .
    \begin{itemize}
        \item When these people stop posting pictures showing off that they work at a prestigious firm, and realise that they and the street-stall worker are the same, with no difference between them --- then they become class for itself .
    \end{itemize}
\end{itemize}

\begin{KeyConcept}
Class for itself = a class which is aware or conscious of the exploitation; being conscious of it they share solidarity; out of this solidarity they have also identified their common enemy; and therefore they bring revolution against the common enemy .
\end{KeyConcept}

\begin{itemize}
    \item Till now we do not have class for itself . We have only class in itself . We have groups working in companies; we do not have groups across companies identifying all these owners as their common enemy . We have only dispersed people working in different companies .
    \item When we have the conscious set of people who recognise that whether they work in this company or that, all of them are exploited, and that their common enemy is one --- the ruling class, whatever the individual owners' names may be --- then class becomes class for itself .
    \item And this solidarity is the result of exploitation . This solidarity is not the result of interdependence --- that is the difference from Durkheim . There is a difference even in solidarity .
\end{itemize}

\begin{QuickRevision}
\begin{itemize}
    \item Class struggle = bourgeoisie vs proletariat .
    \item Class in itself = people sharing a similar relation to the means of production, without awareness of it .
    \item Class for itself = conscious of exploitation + solidarity + common enemy identified + revolution against that enemy .
    \item We have only class in itself today --- dispersed, company-bound, not cross-company .
    \item Marxian solidarity comes from shared exploitation, not from Durkheimian interdependence .
\end{itemize}
\end{QuickRevision}

\subsection{False Consciousness and the Chains That Divide Workers}
\begin{itemize}
    \item The revolutionary class is class for itself, because they have identified the common enemy and are aware they are exploited . They do not become happy on being promoted .
    \begin{itemize}
        \item You will see people updating a status: ``I have just got a promotion --- from junior sales manager to executive sales manager.''  In the language of Marx, that person is the victim of false consciousness .
    \end{itemize}
    \item Unless that person attains true consciousness, he will never be a member of the revolutionary class .
    \item Line to write: Marx believed the members of class in itself are the victims of false consciousness .
    \item Meaning of false consciousness: false consciousness develops when the degree of exploitation is not realised by the worker .
    \begin{itemize}
        \item And even if he realises the exploitation --- if your friend in a big firm says ``this is everywhere, you cannot escape exploitation'' --- then to Marx that person is still a member of class in itself, because he is speaking in a helpless fashion .
        \item The teacher added the real-world qualification the students themselves would raise: these are bookish things, and in reality we have families, people we look up to, and our own next generation to provide for .
    \end{itemize}
\end{itemize}

\paragraph{Why it has become so difficult --- the chains}
\begin{itemize}
    \item In the \emph{Communist Manifesto} Marx writes: ``Workers of the world, unite. You have nothing to lose but chains.'' 
    \item Which chains? Rousseau, the Enlightenment scholar, once said: men are born free, but they are constrained by chains . Marx picked up that chain .
\end{itemize}

\paragraph{Chain 1 --- the nation}
\begin{itemize}
    \item Nation is one chain . Chinese workers are not going to unite with Indian workers, because to them we are Indians and to us they are Chinese .
    \item Simply, nationalism will never let class consciousness develop --- and that is why Marx will never like the nation . And remember, nation is in the superstructure; being part of the superstructure it will never let class consciousness get developed .
\end{itemize}

\paragraph{Chains 2 to 5 --- race, caste, gender, sexuality}
\begin{itemize}
    \item Workers are also divided by race: Black worker, white worker .
    \item By caste: Dalit worker, upper-caste worker .
    \item By gender: male worker, female worker . And by sexuality .
    \item Gender, race, sexuality, nation --- all these are superstructures which will never let the workers unite . These are the ideologies which are ruling over you and will never let class consciousness develop .
    \item Marx says that with these superstructural forms of consciousness, the capitalist class has always justified inequality and exploitation . These things have actually helped the capitalist class exploit the people --- the capitalist class knows that hostile nations will never unite, so let us exploit .
\end{itemize}

\begin{QuickRevision}
\begin{itemize}
    \item False consciousness = the degree of exploitation is not realised; helpless acceptance of exploitation is also class in itself .
    \item Celebrating a promotion is, in Marx's language, false consciousness .
    \item ``Workers of the world, unite. You have nothing to lose but chains'' --- Rousseau supplies the image of the chain .
    \item The chains: nation, race, caste, gender, sexuality --- all superstructural, all preventing class consciousness .
    \item These forms of consciousness are what the capitalist class uses to justify inequality and exploitation .
\end{itemize}
\end{QuickRevision}

\subsection{The Three Prerequisites --- Homogenisation, Pauperisation, Depression}
\begin{itemize}
    \item For a class in itself to become a class for itself, there are certain prerequisites --- certain conditions which have to be met: 
    \begin{enumerate}
        \item Homogenisation 
        \item Pauperisation --- from ``pauper'', meaning a poor person . (Not Karl Popper.) 
        \item Depression 
    \end{enumerate}
\end{itemize}

\paragraph{Homogenisation}
\begin{itemize}
    \item ``Homo'' means same; ``-isation'' means the process of creating a homogeneous reality . You, me, everyone --- irrespective of our skills and education --- will have to become similar, and we will become similar at a certain point of time .
    \item When technologies and machines are introduced in the name of efficiency, all of us will join one rank . You will be thrown out of your company and I will be thrown out of mine; we become a pool of homogeneous people who have been replaced by machines, and whose differences have been obliterated by the introduction of machines .
    \item It does not matter whether you are high-skilled, low-skilled, semi-skilled or unskilled . The highly skilled person will also be replaced by another machine which is also highly skilled .
    \begin{itemize}
        \item At one point of time it was the low-skilled worker on low wages that capitalism required, and that is how capitalism actually came into existence . Today, with the ever-increasing urge for profit, they have hired many high-skilled workers --- but those high-skilled workers do not feel secure in their jobs . ``There is no job security; next day we may be fired.''  And this is why people prepare for the civil services .
    \end{itemize}
    \item So homogenisation means the obliteration of differences between high-skilled, low-skilled, semi-skilled and unskilled --- all these differences will end with the introduction of labour-saving machines .
    \item And then solidarity will be achieved between the high-skilled, low-skilled, semi-skilled and unskilled . A process is going on --- a bond is being made from class in itself towards class for itself .
\end{itemize}

\paragraph{Pauperisation}
\begin{itemize}
    \item Pauperisation means the process of increasing poverty .
    \item Obviously, when all these people are homogenised by the introduction of labour-saving machines, all these victims will become poor . Purchasing power will fall down . If they do not have purchasing power, how will they consume? That means they have become poor --- and this is the process of becoming poor . It is a process .
\end{itemize}

\paragraph{Depression}
\begin{itemize}
    \item Depression here refers to the pressing down of the intermediate strata . The managers, executives and senior executives --- that upper stratum gets depressed, pushed downward, and presses into the proletariat .
    \item When they join the ranks of the proletariat they realise that at the end of the day all of them are exploited only --- none of them are getting what they are producing . This is what Marx says we have to realise .
\end{itemize}

\begin{ComparisonBox}
\begin{itemize}
    \item Recession vs depression: an economy is called to be in recession when it goes through the cycle of lower consumption for at least two consecutive turns .
    \item When the cycle continues for about four, it is called depression . Depression is the extended version of recession, where finally consumption is not happening at all . First recession comes, then depression .
\end{itemize}
\end{ComparisonBox}

\begin{itemize}
    \item When these three conditions are met, class in itself will become class for itself .
\end{itemize}

\begin{QuickRevision}
\begin{itemize}
    \item Three prerequisites: homogenisation, pauperisation, depression .
    \item Homogenisation = obliteration of skill differences by labour-saving machines, producing solidarity .
    \item Pauperisation = the process of increasing poverty; purchasing power collapses .
    \item Depression = the intermediate strata (managers, executives) pressed downward into the proletariat .
    \item Recession = two consecutive cycles of falling consumption; depression = the extended version, roughly four .
\end{itemize}
\end{QuickRevision}

\subsection{Why the Revolution Has Not Happened}
\begin{itemize}
    \item The biggest question: why has class in itself not yet become class for itself? Why has the revolution not yet occurred? 
\end{itemize}

\begin{KeyConcept}
Line to write: the only reason why class in itself has not yet become class for itself is the prevalence of false consciousness perpetuated by the superstructure .
\end{KeyConcept}

\begin{itemize}
    \item Illustration from the lockdown: on social media people started writing that Chinese goods should be dumped and nobody would buy them . This is nothing but false consciousness --- in the name of the nation you got divided .
    \item What is happening is that the bourgeoisie and proletariat of one country are revolting against the bourgeoisie and proletariat of another . What should be happening is that the proletariats of both countries should be criticising the bourgeoisie of both countries . But the opposite is happening --- because the nation has divided us .
    \item In the name of the nation, the difference between bourgeoisie and proletariat disappears entirely . Nation is one ideology which has fooled you --- not letting you realise that you are exploited .
\end{itemize}

\begin{QuickRevision}
\begin{itemize}
    \item The single reason revolution has not occurred: false consciousness perpetuated by the superstructure .
    \item Boycott campaigns organised along national lines are the textbook case --- the class line is erased by the national line .
\end{itemize}
\end{QuickRevision}

\subsection{Alienation Introduced --- Species Being and the Circuit of Commodity}
\begin{itemize}
    \item Alienation is a by-product of everything already covered .
\end{itemize}

\paragraph{Species being at its extreme potential}
\begin{itemize}
    \item In primitive communism the workers had full control over what they wanted to produce . Homo Faber was fully creative, fully expressive; species being --- species essence --- was at its extreme potential .
    \item Meaning: the worker was free to express himself the way he wanted to . He was free to make a fishing boat, free to make a bow and arrow, free to cut and etch stones, free to hunt anything (there was no wildlife protection Act either), free to make paintings on the cave walls and to carve figures on the caves .
\end{itemize}

\paragraph{The decline across modes of production}
\begin{table}[h!]
\centering
\begin{tabularx}{\textwidth}{lY}
\toprule
\rowcolor{AccentDark}
\thead{Mode of production} & \thead{What happens to expression} \\
\midrule
Primitive communism & Full control over what to produce; species being at its extreme potential; free to create, hunt, paint, carve \\
Ancient / slave society & The moment you become a slave you are not free; you do what the master tells you. You are not even free to have a child; what you do and for how many hours is decided by the master. The first society in which a human being became the private property of another human being --- and private property cannot be expressive or creative \\
Feudal society & Serfs were not free to produce what they wanted; they produced what they were asked to produce --- indigo, because the colonial planters wanted indigo. How much they got for the produce was also decided by the lords. No power over what or how much is produced \\
Capitalist society & The proletariat is the worst --- worse than serfs, worse than slaves --- because the relation between humans has turned into a relation between things \\
\bottomrule
\end{tabularx}
\end{table}

\begin{itemize}
    \item The same logic operates in offices: if you want to be creative you are not given enough freedom or room to express it; you are simply asked to do whatever minimum is required from you, and there are different people for the rest .
\end{itemize}

\paragraph{The circuit of commodity}
\begin{itemize}
    \item Primitive communism: $C_1 - C_2$ . I am the producer of tomato, you are the producer of garlic . I am free to produce tomato, you are free to produce garlic . I want garlic, you want tomato --- I gave you tomato, you gave me garlic . Direct exchange . (In the earlier era: I gave you deer, you gave me lion.) 
    \item Ancient and feudal, once commercial purposes appear: $C_1 - M - C_2$ . I produced tomato, I sold tomato, I got money; with this money I can now purchase garlic . There is no direct relation any more --- the relation is through money .
    \item Capitalism: $M_1 - C - M_2$ . You produce tomato, you sell tomato, you get money --- but with this money you are not going to the shop to get garlic . You are going to buy one more acre of land; on this land you will sow more tomatoes so that you have more money; so that you invest more in land; so that in future you buy more and more acres .
    \item It is money in a cycle . You are not consuming; you are only finding ways to earn more and more money .
\end{itemize}

\begin{KeyConcept}
Money, which at one point of time was just a means to the end of consumption, has now in itself become the end . Money has become the end today .
\end{KeyConcept}

\paragraph{The fisherman illustration}
\begin{itemize}
    \item Subsistence (a coastal belt such as Odisha): you are a fisherman, you fish, you consume it . Finished .
    \item Commercial: you fish, you sell the fish, you get money, and with this money you meet some other needs .
    \item Capitalism: you fish, you sell the fish, you get money, and with this money you buy a bigger net to catch more fish, so that you get more money when you sell them, so that you get a still bigger net --- and this goes on and on, until one day you become the leader of the fishermen's association .
    \item You are producing more than what is required --- this is why Gandhi told us to limit our needs .
    \item Marx: labour creates worth more than its value, and this value is pocketed by a class . And when you buy a larger number of nets, ships and boats, you hire more people to manage them --- and there begins capitalist exploitation .
    \item So the relation between humans has turned into a relation between things .
\end{itemize}

\begin{QuickRevision}
\begin{itemize}
    \item Species being was at its extreme potential in primitive communism; expression declines across slavery, feudalism and capitalism .
    \item Circuit of commodity: $C_1-C_2$ (primitive) $\rightarrow C_1-M-C_2$ (ancient and feudal) $\rightarrow M_1-C-M_2$ (capitalism) .
    \item In capitalism money ceases to be a means and becomes the end itself .
    \item Fisherman: subsistence $\rightarrow$ commercial $\rightarrow$ capitalist accumulation of nets; production exceeds requirement .
    \item The relation between humans has turned into a relation between things .
\end{itemize}
\end{QuickRevision}

\subsection{Doubts Taken in Class}

\paragraph{Is depression the same as the formation of the petty bourgeois and lumpen proletariat?}
\begin{itemize}
    \item When there is depression in the economy there is no consumption; with no consumption there is no profit; and when profit is not being gained, you and other people join the same ranks .
    \begin{itemize}
        \item In the lockdown many people became unemployed --- the data shows many people lost their jobs . In fact, during the lockdown the Indian economy witnessed informalisation . The modern idea of development is to move towards formalisation of the economy, but what we actually witnessed was informalisation: more and more people became underemployed and vulnerable, available as forced labour, readily available to do anything for any amount of money .
    \end{itemize}
\end{itemize}

\paragraph{Lumpen proletariat}
\begin{itemize}
    \item Not so important, but note the term . The lumpen proletariat are even below the proletariat, and Marx is not very happy with them .
    \item Who are they? Those who have not studied, who are dependent on alcohol, who lie about with no idea of what is going on in the world --- drug users, alcoholics, peddlers, beggars . Marx does not expect anything from these people .
\end{itemize}

\paragraph{The worker has dissatisfaction and disinterest at work --- how does Marx cater to that?}
\begin{itemize}
    \item Marx will not cater to those people; he is telling us what is going to happen .
    \item Marx clearly says boredom is a very natural process in work . Why do you get bored? When you are doing the same thing every day --- a monotonous, repetitive task .
    \item When you perform a monotonous and repetitive task you become, in the language of Karl Marx, alienated .
    \begin{itemize}
        \item Imagine that worker who had full creativity and could do whatever he wanted . Now you will do only what the capitalist tells you to do, because you are dependent on him for your survival . And that is how the death of expression and the death of creativity creeps in, and you become like a cog in the machine .
    \end{itemize}
\end{itemize}

\paragraph{Difference between unemployment and underemployment according to Marx}
\begin{itemize}
    \item By this logic all are underemployed, as surplus value will always be more than variable capital . Everybody is underemployed .
    \item To an economist, underemployment means something narrower: your degree is B.Tech but your work is that of a peon or a clerk . But underemployment in the Marxian sense is something everybody is facing .
\end{itemize}

\paragraph{If capitalists are already making losses and some are becoming petty bourgeois, why will workers revolt against them?}
\begin{itemize}
    \item There is long literature on this; taken up in the next class --- along with who exactly are the workers who will bring the revolution .
    \item What we do know is that capitalism is not the end of history, and that capitalism has evolved --- it is not the same as the one described by Karl Marx . We also have socialism; but do we have communism anywhere? 
    \begin{itemize}
        \item When Marx talks about communism he is talking about a communist society in a stateless society . What we have today is the communist state --- we have states being ruled by communist parties, and that is not what the idea of communism is to Marx . So when people criticise Chinese or Russian communism they believe they are criticising the Marxian idea of communism, but they are not --- they are only criticising a nation, because communism in the language of Marx has not yet arrived .
    \end{itemize}
\end{itemize}

\paragraph{The point at which Marxism could be proven wrong}
\begin{itemize}
    \item In the language of Marx: if I am earning a lot of money and I am happy and partying all night, I am living in false consciousness --- not able to realise that I am still underpaid and exploited .
    \item Second case: if I am getting bored of my job, of the everyday routine task, then again Marx is right --- because in that case I am alienated .
    \item So what is that one point at which Marxism can be proven wrong? If I am happy, Marx says false consciousness; if I am sad, Marx says alienation . You are right in both cases --- so when will you be wrong? 
\end{itemize}

\begin{CriticalPoint}
This is what Karl Popper is saying: if you cannot be proven wrong ever, then Marxism cannot be science . It is pseudoscience .
\end{CriticalPoint}

\paragraph{Is ours an investment-oriented society?}
\begin{itemize}
    \item Yes --- purely investment-oriented . When you become a millionaire, will you consume millions? No --- you will invest in the stock market, in fixed deposits, recurring deposits, insurance policies .
    \item It is the chase for money that we are doing every day . Consumption is not the end today; money has become the end .
\end{itemize}

\paragraph{Why did Marx not discuss the work process in the administrative arena, since government employees are also proletariat?}
\begin{itemize}
    \item At that point of time there was no concept of the state providing work for the workforce . It was all the industrial capitalist class .
    \begin{itemize}
        \item Remember Adam Smith on self-interested individuals, laissez-faire, the free market economy, the free hand --- let the state not intervene .
    \end{itemize}
    \item The state providing you employment is not a very old concept --- it is a very new concept, which came basically after the World War . Reference given: the Beveridge Report / Beveridge Commission, established after the World War .
\end{itemize}

\begin{QuickRevision}
\begin{itemize}
    \item Depression $\rightarrow$ no consumption $\rightarrow$ no profit $\rightarrow$ ranks merge; the lockdown pushed the economy towards informalisation, not formalisation .
    \item Lumpen proletariat sit below the proletariat and Marx expects nothing of them .
    \item Monotonous, repetitive work produces alienation --- the cog in the machine; boredom is natural to such work .
    \item In Marxian terms everyone is underemployed; the economist's definition is narrower .
    \item The Popper objection: happy = false consciousness, unhappy = alienation; unfalsifiable, therefore pseudoscience .
    \item Money, not consumption, is the end in an investment-oriented economy .
    \item The state as employer is a post-war concept --- Beveridge; in Marx's time it was laissez-faire .
\end{itemize}
\end{QuickRevision}

\sectiondivider

\section{Alienation, Socialism and the First Criticisms}

\subsection{The Essence of Marx, and What Alienation Means}
\begin{itemize}
    \item The essence of the writings of Marx is that workers do not realise that they are exploited . We do not realise that we are exploited, and we are living in false consciousness .
    \item And as long as we are living in false consciousness, the new society will never get developed .
    \begin{itemize}
        \item The analogies used: just as smokers do not realise that the cigarette is harming them; just as people wearing very high heels do not realise the harm; just as certain heads of state did not acknowledge that climate change was happening, or that the pandemic was real --- similarly workers do not realise that exploitation is there .
    \end{itemize}
\end{itemize}

\paragraph{The meaning of alienation}
\begin{itemize}
    \item Alienation simply means you are estranged --- estrangement .
    \begin{itemize}
        \item Suppose you have made something and you cannot buy it . You made the phone; you cannot buy it . You made the car; you cannot buy it . Is it not a very strange thing? 
    \end{itemize}
\end{itemize}

\begin{KeyConcept}
Alienation is the social breakdown of the natural interconnectedness between the worker and his produce or product, and between the worker and the production relations .
\end{KeyConcept}

\begin{QuickRevision}
\begin{itemize}
    \item Essence of Marx: workers do not realise they are exploited; false consciousness blocks the new society .
    \item Alienation = estrangement --- you make the thing and cannot possess it .
    \item Definition: social breakdown of the natural interconnectedness between worker and product, and between worker and production relations .
\end{itemize}
\end{QuickRevision}

\subsection{Alienation is Not Exclusive to the Proletariat}
\begin{itemize}
    \item Generally what we believe is that proletariats are the victims of alienation . Marx said even the bourgeoisie are alienated . Just remember this .
    \item Because the bourgeoisie are also individuals . The production process and the relations of production --- everything is external to the actor, and not just to the proletarian actor but to the bourgeois actor as well .
    \begin{itemize}
        \item The whole economy is over and above the biggest industrialist . The industrialist is not over the economy; the economy as a system is over and above the biggest business houses . Therefore these are also alienated .
    \end{itemize}
\end{itemize}

\begin{CriticalPoint}
Marx believed alienation is not exclusive to proletariats, because the bourgeoisie are equally alienated . However, the alienation of proletariats is of an extreme type and of a different degree --- quantitatively higher alienation, qualitatively different alienation .
\end{CriticalPoint}

\paragraph{Two illustrations of proletarian alienation}
\begin{itemize}
    \item The simple example: you are a worker in a noodles factory . You make the noodles, but you cannot buy the noodles --- because you are not given enough income or salary to buy it . You are alienated from your own produce .
    \item The most basic example: all of you want to become civil service officers . But you know very well that you are highly talented and highly creative in your respective fields --- maybe singing, maybe dancing, maybe acting, maybe the guitar .
    \begin{itemize}
        \item But you are not trying to polish those skills . Rather you want to become an officer . In that sense you are alienated from your desires --- because playing the guitar, singing, dancing or acting are not given so much status in society . So your wish to become an officer is not your wish .
        \item You are alienated from your creative self, your inner self, your potential self --- from the species being, the species essence .
    \end{itemize}
\end{itemize}

\paragraph{Why the job itself was never a choice}
\begin{itemize}
    \item Recall: before you become a philosopher, before you start thinking and speaking on any given subject, you have to eat something first . The very first, primary need is survival .
    \item For survival you would inevitably enter into social relations of production . So the basic question --- why do you want to do a job? To survive . Going for a job is therefore not your choice; it was inevitable .
    \item But why did you not go for the job that you like? You are going for the job that is given status in society, not for the thing you feel creative with . So you are alienated from your desires and from your creative potential .
\end{itemize}

\begin{QuickRevision}
\begin{itemize}
    \item Bourgeoisie are alienated too, because the economy as a system is external to them as well .
    \item But proletarian alienation is quantitatively higher and qualitatively different --- of an extreme type .
    \item You make the product and cannot buy it; you have talent and abandon it for a higher-status occupation .
    \item Taking a job at all is inevitable (survival); which job you take is where alienation from desire shows .
\end{itemize}
\end{QuickRevision}

\subsection{The Four Types of Alienation}
Alienation is of four types .
\begin{enumerate}
    \item Alienation from species being 
    \item Alienation from the product 
    \item Alienation from the production process 
    \item Alienation from production relations (sociability) 
\end{enumerate}

\paragraph{1 \& 2 --- Species being and the product}
\begin{itemize}
    \item The first two have already been covered above . Alienation from the product: what you have made does not belong to you . It belongs to the person who asked you to make that thing .
\end{itemize}

\paragraph{3 --- Alienation from the production process}
\begin{itemize}
    \item First of all, you did not make the whole product by yourself . In a car you made only the rubber of the tube that goes into the tyre at the very back --- not even the whole tyre .
    \item Imagine there was a point of time when you could make the whole product all by yourself and you used to take pride in your craftsmanship .
    \item Craftsmanship does not exist today . It is all machine work that you see today .
\end{itemize}

\begin{ExampleBox}
\begin{itemize}
    \item When you go to the Taj Mahal or the Red Fort, nobody today can make that Taj Mahal or that Red Fort at today's rates --- even though machines have arrived .
    \item And look at our expression at these monuments: ``how could those people make this kind of thing in that time, when we cannot make it today?'' That question is wrong . They could make those things because they did not have machines .
    \item When machines start making things, things look uniform --- a machine works to create standardisation . A machine can create ``Taj Mahals'' --- the ten-rupee replicas sold in the shops outside --- but it cannot create the Taj Mahal . The Taj Mahal can only be created by craftsmen .
\end{itemize}
\end{ExampleBox}

\begin{itemize}
    \item The whole production process has been broken down . Some workers work on the tyre; some on the window of the car; some on the nuts and screws that fix the wheel to the platform; and some others design the platform itself .
    \item And the people who make the tyre do not really know what final shape their produce is going to take .
    \begin{itemize}
        \item Software engineer: they code and code and code all their lives --- but they do not even know what they are actually coding for . They are operating on a minuscule part of a larger operation they are not even conscious of . They only know that they have to learn the big languages and write that code .
        \item The same with people working on mathematical operations --- solving a theorem without knowing what its application in the real world is . They simply solve the equations, because they know solving them will get marks .
    \end{itemize}
    \item Those making the glass of the car do not even know what the final car is going to look like . Every day they come to the assembly line and make these glasses in the same shape with the same dimensions .
    \item The creative potential of that man has been finished . He could make all the things by himself; now, in the name of specialisation, he has actually become alienated from the larger process .
\end{itemize}

\paragraph{4 --- Alienation from production relations, or sociability}
\begin{itemize}
    \item Think of the people who work whole day and night . They are so involved in their work that they do not even know who their neighbours are .
    \item They do not get enough time to go out with their family; they do not get enough time to spend with their friends .
    \item And moreover, they do not even get time to chat with their fellow employees in the factory or office where they work . They cannot even have a casual conversation .
    \begin{itemize}
        \item They register themselves with the biometric, go to the computer, keep working there, eat there, sit there, and then leave . This is the life .
    \end{itemize}
    \item If you look at the life of those people from the top, it seems a repeated, monotonous, stereotyped pattern that they have become part of --- a cog in the machine . The machine keeps going, and looked at from above, they keep going in exactly the same process . No change at all .
\end{itemize}

\begin{QuickRevision}
\begin{itemize}
    \item Four types: species being, product, production process, production relations (sociability) .
    \item Product: what you make belongs to whoever ordered it .
    \item Production process: the process is broken into fragments; the maker of a part cannot see the whole --- the death of craftsmanship in the name of specialisation .
    \item Machines standardise: they can make Taj Mahals, not the Taj Mahal .
    \item Production relations: no neighbours, no family time, no friends, not even casual conversation at work --- a cog in the machine .
\end{itemize}
\end{QuickRevision}

\subsection{How the Bourgeoisie are Alienated and in False Consciousness}
\begin{itemize}
    \item Marx says it is not only that the bourgeoisie are alienated --- even the bourgeoisie are the victims of false consciousness . False consciousness is on both sides; alienation is on both sides . Generally people believe only the working class lives in false consciousness --- no, everyone is in false consciousness .
\end{itemize}

\paragraph{The false consciousness of the bourgeoisie}
\begin{itemize}
    \item The bourgeoisie think that they are earning profit through constant capital --- through machines .
    \begin{itemize}
        \item Otherwise why would they introduce machines? The only reason for the introduction of machines is that they think machines can give them profit .
    \end{itemize}
    \item But when machines are introduced, in the long run they realise that it is a loss --- the rate of surplus value is declining . It takes them a very long time to realise it; and then Marx comes and explains it .
\end{itemize}

\paragraph{The alienation of the bourgeoisie}
\begin{ExampleBox}
\begin{itemize}
    \item Take the standard flavour of instant noodles --- the classic yellow packet . We have been eating that same packet forever .
    \item It is not that the manufacturer did not try to introduce different flavours . They tried . But the market did not accept them --- the consumers did not accept those flavours, even with advertising campaigns behind them .
    \item So even the manufacturers have become alienated from their own produce . They are not able to introduce any other flavour, though they think every day about how to move beyond that yellow packet .
\end{itemize}
\end{ExampleBox}

\begin{itemize}
    \item The market has assumed a reality of its own . They have created the product, but they are not in control of this product now --- people want that particular version and nothing else .
    \item That is why the bourgeoisie have also become alienated from the product . The same is visible in a car model that has run essentially unchanged for years, and in phone models numbered one after another .
    \item But this is not a very great concern to Marx . His main concern was the alienation of proletariats .
\end{itemize}

\begin{KeyConcept}
What Marx actually wanted to tell us by discussing bourgeois alienation is that the market is over and above both these players . The market has assumed a reality which none of these players can control .
\begin{itemize}
    \item An aside from class: people call every brand of instant noodles by the name of the leading brand --- brand names becoming generic is itself a mark of this .
\end{itemize}
\end{KeyConcept}

\begin{QuickRevision}
\begin{itemize}
    \item Even the bourgeoisie live in false consciousness --- they believe machines (constant capital) generate their profit .
    \item In the long run they discover the rate of surplus value is falling instead .
    \item Bourgeois alienation: the producer can no longer control the product; the market dictates it .
    \item The real point: the market is over and above both classes and beyond the control of either .
\end{itemize}
\end{QuickRevision}

\subsection{Which Type of Alienation is Commodity Fetishism?}
\begin{itemize}
    \item Answer: alienation from the product .
    \item You have made a part of the phone --- not the full phone . No person today makes a full product . Some make the screen; some make the battery, in an altogether different department; some make other components .
\end{itemize}

\paragraph{The global value chain}
\begin{itemize}
    \item Today we are in the era of the global value chain --- the whole production process has become globalised . The company is American, but the components and the production process are in China . The assembly line spans from America to China .
    \item And it is not just the raw-material components: marketing, operations and sales all go into it . It is not so simple today .
    \begin{itemize}
        \item Think of the contribution, in percentage terms, of each person in the whole production process of one phone . Even the person who greets you when you enter the showroom is part of the larger production process --- because the brand runs on an image, an image of being very customer-friendly, and the phones themselves are user-friendly . So the very smile of the seller is also part of the production process .
    \end{itemize}
\end{itemize}

\paragraph{What happens at the point of purchase}
\begin{itemize}
    \item When you as a consumer go to the showroom and see the newly launched model, you start looking at it as something magical --- as extraordinary compared with the rest .
    \item In the whole process of attributing magical quality to the commodity, you do not realise that it is simply a product of the labour process --- the product of a labourer who did not get anything like the amount you are going to shell out from your pocket to buy it .
    \item So the very idea of the status symbol is not real . But this is something you would not realise; you would simply buy it . This is the fetish of the commodity .
\end{itemize}

\paragraph{The religious parallel and Freud}
\begin{itemize}
    \item Just as with God: you have created God, but you do not realise that it is your own creation --- and then you ask that creation for favours, a minute after having made it .
    \item This idea comes straight from Feuerbach's writing on religion .
    \item Sigmund Freud --- a very important scholar whose importance will be felt in later topics --- says religion is all about wish fulfilment: you wish that someone were there, the grandfather, the grandest of all fathers, God . Similarly, you simply wish to have that phone .
    \begin{itemize}
        \item You may already have a better phone, but you would still want it . There is no rationality in it .
    \end{itemize}
\end{itemize}

\begin{CriticalPoint}
\begin{itemize}
    \item If we were rational --- in the language of the economists who modelled human nature as Homo economicus, a conventional model not applicable today: highly utilitarian, rational, interest-maximising, selfish --- then why is he not saving his money? 
    \item Why is he spending ruthlessly on products that do not have any value per se? It is all about the fetish that people have about these commodities .
\end{itemize}
\end{CriticalPoint}

\begin{QuickRevision}
\begin{itemize}
    \item Commodity fetishism corresponds to alienation from the product .
    \item Nobody makes a whole product: the global value chain fragments it across countries and functions, including the salesperson's smile .
    \item At the point of purchase the commodity appears magical and the labour behind it disappears .
    \item Feuerbach on religion and Freud on wish fulfilment are the parallels; Homo economicus cannot explain the behaviour .
\end{itemize}
\end{QuickRevision}

\subsection{The Contradiction Inside Marx --- Orthodox and Hegelian Marxists}
\begin{itemize}
    \item Till capitalism everything is fine . Marx said inherent contradictions within capitalism will sow the seeds of its destruction, and then we will have socialism . In that sense, workers do not need to do anything --- capitalism will collapse automatically, just as every previous mode of production collapsed .
    \item But there is one more statement in the \emph{Communist Manifesto}: ``Workers of the world, unite. You have nothing to lose but chains.'' 
    \item These two things are very contradictory . On the one hand he is pleading with the workers to realise that they are exploited, that they are in false consciousness, and to revolt . On the other hand he says capitalism is destined to collapse automatically due to inherent contradictions --- the law of motion, capital accumulation, polarisation, homogenisation, pauperisation, depression, all inevitably forming two clear-cut classes .
    \begin{itemize}
        \item Maybe people are not smart enough to see that these two are not opposites . But since Marx is dead, we cannot ask him which side he is on --- so people became Marxists all by themselves . Some go with one philosophy, some with the other .
    \end{itemize}
\end{itemize}

\begin{table}[h!]
\centering
\begin{tabularx}{\textwidth}{lYY}
\toprule
\rowcolor{AccentDark}
\thead{Strand} & \thead{Which statement they follow} & \thead{Position on revolution} \\
\midrule
Orthodox Marxists & Inherent contradictions will inevitably collapse capitalism & Revolution is inevitable \\
Hegelian Marxists & ``Workers of the world, unite --- you have nothing to lose but chains'' & Revolution is not inevitable --- the working class will have to overthrow the capitalists; they will have to do something \\
\bottomrule
\end{tabularx}
\end{table}

\paragraph{Why ``Hegelian''?}
\begin{itemize}
    \item Marx was against Hegel --- so what is a Hegelian Marxist? Hegel was all about consciousness --- the world consciousness, which is evolving; and whatever changes we see in society are nothing but the evolution of consciousness, the evolution of ideas .
    \item Marx says we are living under false consciousness, so we need true consciousness --- true consciousness being the realisation that these people are exploiting us .
    \item So when Marx said ``workers of the world unite'', the interpreters of Marx saw Hegel and Marx coming together: Hegel was all about the evolution of consciousness, and Marx wants the workers to develop the very consciousness that has not yet let them realise that they are exploited .
    \begin{itemize}
        \item Doubt raised: Hegel's consciousness is about the world spirit --- how will man develop it? Answer given: ``it will come'' --- when philosophers talk about anything there is no concept or logic; it is random . That is exactly why philosophy has a problem getting the status of science .
    \end{itemize}
\end{itemize}

\begin{QuickRevision}
\begin{itemize}
    \item Contradiction: automatic collapse (inevitability) vs the call to unite (agency) .
    \item Orthodox Marxists = revolution is inevitable .
    \item Hegelian Marxists = revolution is not inevitable; the working class must act .
    \item The Hegelian label comes from Hegel's primacy of evolving consciousness, which Marx's call for true consciousness echoes .
\end{itemize}
\end{QuickRevision}

\subsection{Liberalism and the Three Reactions to It}

\paragraph{Liberalism}
\begin{itemize}
    \item Liberalism is basically Adam Smith's idea: the market should be liberated from any intervention of the state . Laissez-faire . Free market economy .
    \item The principle of liberalism: the more wealth is accumulated, the more development happens . The more money you make, the more developed you become .
    \item The problem: when you are making a lot of wealth, your concern is only to make wealth . And when your only concern is to accumulate wealth, you are not going to take care of the workers who are working for you --- not their physical or mental health, not the working conditions they live in, not whether they can survive on the income you give them . They have no pension, no insurance, and you are not concerned about these things .
    \item The victims of liberalism were the working class .
\end{itemize}

\begin{DataFact}
\begin{itemize}
    \item 1750 --- the industrial revolution era begins .
    \item 1750 to 1800 --- the dominant period of liberalism .
    \item After 1800 --- great movements begin, starting in Britain .
\end{itemize}
\end{DataFact}

\paragraph{Reaction 1 --- the Luddites (1800 to the 1810s)}
\begin{itemize}
    \item The Luddites are those people who are against technological development .
    \item Why would anyone be against technological development? Because technology was developing at the cost of their employment .
    \item These were the people who attacked factories and attacked the machines, because with the introduction of machines they were the ones expelled from the factories . So they formed a band and charged against the machines that had cost them their jobs .
    \item When they charged in full-fledged fashion against those machines, the industrialists called the state administration and the bureaucracy --- the police and administrative officers .
    \begin{itemize}
        \item And what would those officers do? Asked who is wrong, they conclude the workers are wrong . So you are serving the interest of the capitalist class --- because you do not know the whole story, the whole narrative, or the philosophical idea that led to these developments . You just know one action --- the surface act . You do not know the deep inner working of the law of motion of capitalism . So you arrest the workers; and the workers abuse you too, saying you are also involved with these people .
    \end{itemize}
    \item The Luddite movement was a grassroot movement --- it started out as a local movement, but eventually spread all over England .
\end{itemize}

\paragraph{Reaction 2 --- Chartism}
\begin{itemize}
    \item Chartists were those who said we need to introduce the working class into the parliamentary system .
    \item Before Chartism, workers were not represented in parliament --- it was all elites and the bourgeoisie . The industrialist class was also the parliamentarian class, so it was inevitable for the police and the officers to serve them .
    \item Why were workers not represented? Because they were not owners of private property . In those times, if you wanted to become a member of parliament in the British house, you needed to have private property --- that was one condition .
    \begin{itemize}
        \item Today there is no such condition: you need to be a citizen and to meet the age requirement . But at that time the private-property condition automatically served the interest of the capitalist class, because they were the ones who had private property, while the working class was working merely for subsistence .
    \end{itemize}
    \item Chartism was also a grassroot movement . A People's Charter was implemented --- dated in the lecture as 1834 [as stated in lecture] --- through which the working class could now find themselves represented in the British house .
\end{itemize}

\paragraph{Reaction 3 --- socialism}
\begin{itemize}
    \item Socialism was the third --- and the important --- reaction to liberalism .
    \item Note on usage: when the word liberalism is used here, it automatically implies capitalism . Liberalism and capitalism are synonyms in this discussion .
    \begin{itemize}
        \item Why could Marx say that with the introduction of machines workers would be replaced? Because he could see it . Although he was born in 1818 and died in 1883, he had seen in the historical records the movements --- such as the Luddites --- of those who were the victims of technological advancement .
    \end{itemize}
\end{itemize}

\paragraph{What socialism means}
\begin{itemize}
    \item Socialism is not a unified ideology --- socialism is of multiple types .
    \item Liberalism served the self-interested individuals of Adam Smith's language: pursue self-interest, no matter what; there is an invisible hand which will automatically serve all the people, because if you want to become a billionaire you will obviously give jobs to many people, so your self-interest serves society as a whole . But liberalism is all about self-interest at the end of the day --- it is highly individualist .
    \item On the other hand, socialism is about collectivism --- against individualism .
    \item Collectivism is all about collectively owning the resources --- that is, public ownership .
    \begin{itemize}
        \item Who is the owner of a public sector undertaking? You and I . Because in India the public is sovereign --- whereas in Britain Parliament is sovereign . We are sovereign; we give money to run it; so we are the owners .
    \end{itemize}
    \item But it is very utopian . How is the public going to own? How does it happen in reality? Is there any method? 
\end{itemize}

\paragraph{The word ``utopia''}
\begin{itemize}
    \item Thomas More wrote a book called \emph{Utopia}, and coined the word, specifically to describe the ideal society --- the ideal society we should have, where we do not have wealthy individuals exploiting workers .
    \item And all this is happening before Marx --- Marx has not yet arrived; these discussions come earlier .
    \item Utopia is a dream world . This early strand is therefore called utopian socialism .
\end{itemize}

\begin{QuickRevision}
\begin{itemize}
    \item Liberalism = Adam Smith, laissez-faire, free market; principle: more accumulation $\rightarrow$ more development; victims: the working class .
    \item 1750 industrial revolution; 1750--1800 the dominant period of liberalism; after 1800, reactions begin in Britain .
    \item Luddites (1800--1810s): against technological development because it cost them employment; attacked machines and factories; a grassroot movement that spread across England .
    \item Chartism: demanded working-class representation in parliament, blocked by the private-property qualification; grassroot; the People's Charter [1834 as stated in lecture] .
    \item Socialism: the third reaction --- collectivism and public ownership against liberal individualism .
    \item Thomas More coined ``utopia'' in his book describing the ideal society, before Marx .
\end{itemize}
\end{QuickRevision}

\subsection{The Utopian Socialists}

\paragraph{Robert Owen (Owenism)}
\begin{itemize}
    \item Owen said the biggest problem is the condition of workers, and we can improve it if we give them good education . Education is the basic necessity: if people are educated they will take care of themselves .
    \item So our duty should be to give them good health and good education --- the money they will earn on their own .
    \item He also focused on cooperatives .
    \begin{itemize}
        \item What is a cooperative, and how is it different from a private limited company? In a cooperative, whatever revenue the company makes is distributed to everyone . The dairy sector has many such cooperatives .
    \end{itemize}
    \item He also talked about trade unions --- form trade unions . He was basically concerned about the working class only, and that is why he gave the ideas of cooperatives, trade unions, good health and good education .
\end{itemize}

\paragraph{The New Lanark experiment}
\begin{itemize}
    \item The problem: he did not have enough money or capital to give everyone health and education, or to run campaigns for the cooperative movement . So Owen said: I will show this to the world by doing one experiment in my countryside .
    \item The name of that place was New Lanark . He established a huge textile factory there . Today New Lanark is a World Heritage Centre .
    \begin{itemize}
        \item Comparison offered in class: much like a company town where the school, the market and everything else belongs to one industrial house .
    \end{itemize}
    \item He then demonstrated New Lanark to the capitalists in Britain, in London, and pleaded with them: you should also do something like this for the workers . If you do not take care of your workers, how will you earn profit? These are the people whose hard work is giving you profit .
    \begin{itemize}
        \item He was requesting them with folded hands --- and that is precisely why the capitalists were not impressed .
    \end{itemize}
\end{itemize}

\paragraph{Charles Fourier and Claude Henri de Saint-Simon}
\begin{itemize}
    \item After Robert Owen came Charles Fourier . Owen and Fourier were British citizens [as stated in lecture] . Fourier's idea was not very different from Owen's .
    \item And then came Claude Henri de Saint-Simon, the Frenchman --- also a utopian socialist . Saint-Simon was more about cooperatives and centralised planning .
\end{itemize}

\paragraph{Why ``utopian''}
\begin{itemize}
    \item Utopian in the sense that they believed the capitalist would help the workers, would establish cooperatives for them, would give them good health and good education . But that does not happen .
\end{itemize}

\begin{CriticalPoint}
\begin{itemize}
    \item Present-day link: corporate social responsibility (CSR) . If you look into CSR reports you will see how much is spent on welfare --- and what is actually done with that money .
    \item The idea of CSR is actually the gift of the writings of these utopian socialists .
\end{itemize}
\end{CriticalPoint}

\begin{QuickRevision}
\begin{itemize}
    \item Robert Owen: good health, good education, cooperatives, trade unions; the New Lanark textile factory experiment, today a World Heritage Centre; pleaded with capitalists and was ignored .
    \item Charles Fourier: broadly the same idea .
    \item Saint-Simon (French): cooperatives and centralised planning .
    \item Utopian because it assumes capitalists will voluntarily help --- CSR is the modern descendant of this idea .
\end{itemize}
\end{QuickRevision}

\subsection{Scientific Socialism}
\begin{itemize}
    \item After utopian socialism came scientific socialism --- what we today call Marxism . Marxism is scientific socialism, not utopian socialism .
    \item Marx believes we do not have to plead with the capitalist class . Peaceful methods are not going to help anyone . In fact the more you request, the more it reveals the exploitation you have been the victim of .
    \item You should be authoritative --- because it is only because of your labour that those people have gained profit . The very rooms in which they sit, the staff room, the ball room, everything, is made by your labour . So go with full authority and full confidence; do not be nervous at all .
    \item The proletariats have to bring a revolution against the bourgeoisie, because they have already suffered a great deal of exploitation .
    \begin{itemize}
        \item Many people died simply by working in the factories --- bad health, unsafe and insecure working conditions, no pension, no income, no insurance, and child labour on top of that .
        \item And it is not happening to one generation only . Generation by generation, the working class produces more working class . Since they are not educated they cannot educate their children; these uneducated people act as cheap labour to the capitalist, who is very happy because cheap labour does not have to be paid much . They will not even encourage them to study more --- because with education they would start talking about rights .
    \end{itemize}
\end{itemize}

\paragraph{The definition}
\begin{KeyConcept}
Scientific socialism is the result of the inherent contradictions within capitalism . These inherent contradictions will automatically reveal two clearly polarised classes --- bourgeoisie and proletariat . Either the proletariats will violently overthrow these people, or, if they do not, capitalism will collapse automatically . In both ways capitalism is doomed to collapse .
\end{KeyConcept}

\paragraph{Why is it called ``scientific''?}
\begin{itemize}
    \item Because it is a natural outcome . Inherent contradiction --- just like any other mode of production destined to collapse . It is happening naturally: law of dialectics, law of motion .
    \item Therefore socialism is a natural outcome, and that is why it is scientific . It is not the outcome of the working class per se; the inherent contradictions will automatically polarise the world into two, and then one will overthrow the other .
\end{itemize}

\paragraph{The principles of scientific socialism}
\begin{enumerate}
    \item The means of production would be owned and managed by the working class .
    \item The working class will produce and distribute the goods and services according to the work --- according to the working labour . If you have done work for so many hours, whatever should have been given to you in the first place, and was instead pocketed by the capitalist, will now be given to you . There will be no surplus value appropriation by the capitalist .
    \item Once the distribution process is over, the state shall wither away .
    \begin{itemize}
        \item On point 3: capital has been accumulated for a very long time, so this accumulated surplus capital has to be distributed first . Imagine you are sitting on behalf of the workers, and each worker comes and says ``I was working in this factory for ten years'' . You calculate how much surplus value should have been received in those ten years against what was actually received, and the difference is paid out . Once that whole process is over, the state withers away .
    \end{itemize}
\end{enumerate}

\paragraph{The socialist state}
\begin{itemize}
    \item In scientific socialism the state is managed by the working class, and the objective of this socialist state is the distribution of goods and services .
\end{itemize}

\begin{GovtPolicy}
In fact India also had this objective: in 1950, when the Planning Commission was established, its objective was precisely to distribute goods and services .
\end{GovtPolicy}

\begin{itemize}
    \item Socialism is nothing but the dictatorship of the proletariat --- the working class manages the affairs of the economy . They decide what is to be produced, how much is to be produced, and what shall be given to those people involved in the production process .
    \item Once everything has been given according to their labour time, the state shall wither away . Once the state has withered away --- because we have achieved equality --- we will have communism: stateless, classless, moneyless, profitless .
\end{itemize}

\begin{QuickRevision}
\begin{itemize}
    \item Scientific socialism = Marxism; no pleading, no peaceful petitioning --- the proletariat must act with authority .
    \item It is ``scientific'' because socialism is the natural outcome of inherent contradictions, not of working-class will per se .
    \item Three principles: means of production owned and managed by the working class; distribution according to labour performed with no appropriation of surplus; the state withers away once distribution is complete .
    \item Socialism = dictatorship of the proletariat; the socialist state's only objective is distribution (compare the 1950 Planning Commission) .
    \item Communism: stateless, classless, moneyless, profitless .
\end{itemize}
\end{QuickRevision}

\subsection{Democratic Socialism}
\begin{itemize}
    \item Socialism is not of one type . There is Owen, Fourier and Saint-Simon; there is Marx's version; and there is another version --- democratic socialism .
    \item Democratic socialism is not Marxist socialism or scientific socialism . Why? 
\end{itemize}

\begin{table}[h!]
\centering
\begin{tabularx}{\textwidth}{lYY}
\toprule
\rowcolor{AccentDark}
\thead{} & \thead{Marxist / scientific socialism} & \thead{Democratic socialism} \\
\midrule
Role of the state & The state is simply an instrument that will achieve equality & The state is not the instrument --- the state is the end \\
Fate of the state & Once equality is achieved the state withers away & The state is there to stay forever \\
What is equally important & Only the economic transformation & Both democracy and socialism are equally important \\
\bottomrule
\end{tabularx}
\end{table}

\begin{itemize}
    \item Democracy is a political concept; socialism is an economic concept . In democratic socialism, socialism is managed and run by democracy --- the politics .
    \item And it is a democracy chosen by the people, the working class . The working class chooses workers . In Marxian terms, we are the proletariats giving votes to proletariats to rule over us --- proletariats ruling the proletariats --- and these proletariats are running the economy . And that is why it is a socialist economy .
    \begin{itemize}
        \item The teacher noted that describing us all as proletariats is not highly accurate --- that will become clear when the criticism is taught . But let us use the Marxist terminology .
    \end{itemize}
    \item And since we are socialist, nobody can be poor . Everybody will be given education, insurance and pension --- and that is why so many welfare schemes are doled out every year: employment guarantee schemes, health missions covering primary, secondary and tertiary care, and free education from the age of 6 to 14 [as stated in lecture] .
    \begin{itemize}
        \item Scandinavian countries follow this model . But there is a huge difference between them and us: the population issue does not exist there --- that is why they are always ranked first on happiness . The scale of the problem here is different; even getting out of a city can take two hours .
        \item A doubt raised --- is the process becoming a hoax? The answer given: you are privileged people living in 2020 and can easily call it a hoax; think about the times in which Marx was writing .
    \end{itemize}
\end{itemize}

\begin{QuickRevision}
\begin{itemize}
    \item Democratic socialism $\neq$ Marxist socialism: the state is the end, not an instrument, and never withers away .
    \item Democracy is political, socialism is economic; the working class votes for the working class .
    \item Welfare schemes --- employment guarantee, health missions, free schooling --- are its expression .
    \item The Scandinavian countries are the model case; population scale is the difference noted .
\end{itemize}
\end{QuickRevision}

\subsection{Communism as Full Expression of the Self}
\begin{KeyConcept}
When Marx says we have attained communism, that means we have attained a civilisation where you can hunt in the morning, fish in the afternoon, graze cattle in the evening and criticise after dinner --- without ever becoming either a fisherman, or a hunter, or a grazer, or a critic .
\end{KeyConcept}

\begin{itemize}
    \item In that sense you would express your full creativity .
    \item Today you are either a teacher or an engineer, either a doctor or a professor, either a civil servant or a teacher . You cannot be both .
    \begin{itemize}
        \item Very basic question: can any person be both a doctor and a civil service officer? If you are a civil servant you can never practise medicine --- never . The civil service code forbids pursuing anything for personal gain; charity is permitted, but nothing for personal gain --- the office-of-profit logic .
    \end{itemize}
    \item Marx believed that today we are constrained in occupations . We cannot express our potential .
    \item But tomorrow, when we have communism, we will be fishermen, hunters, gatherers, herders and critics --- everything at one time, in one whole day, without ever being called or recognised as any one of them .
\end{itemize}

\paragraph{Why we cannot pursue two passions equally --- the working day again}
\begin{itemize}
    \item There are many people who are not able to pursue both interests and passions equally, because they cannot give equal time to both . And this is what the concept of the working day is all about .
    \item If you are an officer you would have to work from 10 a.m. to 5 p.m. The length of the working day will not let you do anything else . It is something external to you, and it restricts you from pursuing any other interest .
    \item And this length of the working day is a capitalist concept that we are following .
\end{itemize}

\begin{CriticalPoint}
In fact the very idea of Sunday --- of the weekend --- is the result of working-class movements, labour movements and Ambedkarite movements . These people have given us Sundays . Otherwise every day was the same working day .
\end{CriticalPoint}

\begin{QuickRevision}
\begin{itemize}
    \item Communism = hunt in the morning, fish in the afternoon, herd in the evening, criticise after dinner --- without becoming any one of these .
    \item Today occupations constrain us; you cannot be a doctor and an officer at once .
    \item The length of the working day is external, coercive and capitalist in origin .
    \item The weekend itself is a product of labour and Ambedkarite movements .
\end{itemize}
\end{QuickRevision}

\subsection{Doubts on Socialism and Communism}

\paragraph{We develop efficient growth through competition, which is not in communism}
\begin{itemize}
    \item How do you know it is not in communism? Nobody has seen communism . If nobody has seen it, how do you criticise it? 
    \item ``I am telling you what Marx is telling me to tell you. I am not praising Marx, nor am I criticising you. You cannot say Marx is wrong, because you have not seen that thing yet.'' 
\end{itemize}

\paragraph{Does Marx not talk about incentivising merit, talent and hard work?}
\begin{itemize}
    \item All of that is utopia . The question of ``hard-working people'' is taken up properly when stratification is taught .
\end{itemize}

\paragraph{Tribal societies have communism, but we see power relations there}
\begin{itemize}
    \item Yes --- power relations are definitely there . Power you cannot avoid .
\end{itemize}

\begin{CriticalPoint}
And this is what the critique of Marxism is also about: society may be classless, but society can never be powerless . And the moment you have power, class comes again . Society can be classless, but can never be statusless --- and class and status are two different things .
\end{CriticalPoint}

\begin{itemize}
    \item Many critics of Marx come from elite theory --- Vilfredo Pareto; and the idea of lions and foxes from Machiavelli's \emph{The Prince} .
    \item Elites are born to rule over us; society can never be equal; elites are destined to be born . Some will be lions and some will be foxes, and the ideal ruler will be the one who combines the best elements of both .
    \begin{itemize}
        \item The teacher applied this typology to contemporary national leaders by way of illustration, while noting that such applications cannot be empirically tested and remain a matter of opinion .
    \end{itemize}
\end{itemize}

\begin{QuickRevision}
\begin{itemize}
    \item You cannot criticise communism empirically --- it has never existed .
    \item Merit and incentives are utopian in this frame; taken up under stratification .
    \item Key critique: a society can be classless but never powerless or statusless; class and status are different .
    \item Pareto's elite theory and Machiavelli's lions and foxes: the ideal ruler combines both; such labels remain untestable opinion .
\end{itemize}
\end{QuickRevision}

\subsection{Criticism of Alienation --- Framing, and C. Wright Mills}

\paragraph{What ``classical scholar'' means}
\begin{itemize}
    \item Durkheim, Marx and Weber are classical scholars . Meaning: they can never be criticised .
    \begin{itemize}
        \item That does not imply criticising them is immoral . It implies that even if you criticise them, you will find yourself supporting their ideas . At the end of the day you remain in the trap and cannot escape it .
    \end{itemize}
    \item So a critic of Marx is a Marxist only . He may be a non-Marxist; he may be an anti-Marxist --- but he will always use Marx's own concepts (capitalism, socialism, communism) while criticising him . You are using Marxist concepts to criticise Marxism, and therefore you cannot escape this vicious cycle .
\end{itemize}

\paragraph{C. Wright Mills}
\begin{itemize}
    \item C. Wright Mills --- you know him from sociological imagination . He is the first person in America to introduce Marxism in universities; Marxism was never taught in US universities until the 1930s .
\end{itemize}

\begin{DataFact}
\begin{itemize}
    \item Sociology is a very young discipline --- in fact the youngest of all the disciplines, hardly 100 years long; a centenary was celebrated only recently .
    \item Philosophy is as old as human beings --- the mother of all sciences, older even than mathematics . Political science is also a very old subject .
    \item Sociology is young because it came as a response to modernity --- and modernity is itself a very young concept .
\end{itemize}
\end{DataFact}

\paragraph{The argument}
\begin{itemize}
    \item Mills says: when Marx was analysing capitalism, the capitalism of those times was industrial manufacturing capitalism --- what we today call the secondary sector .
    \item The tertiary or service sector did not exist in as developed a form as we have today . In those times it was a full-fledged manufacturing sector, and that is why the workers were blue-collar workers .
    \begin{itemize}
        \item ``Blue collar'' is a symbolic term --- it does not mean they literally wore blue shirts .
    \end{itemize}
    \item Line to write: Marx's analysis of capitalism, and of alienation, was limited to the manual working class --- that is, blue-collar workers .
    \item C. Wright Mills has discussed alienation among non-manual workers --- the workers of the service sector, the white-collar workers .
\end{itemize}

\paragraph{Prostituting the personality}
\begin{itemize}
    \item Mills says these white-collar workers are ``beautiful idiots'' . Why? 
\end{itemize}

\begin{KeyConcept}
Because in the sales room, in the board room, in the ball room and in the conference room, these people are prostituting their personalities --- selling the pieces of their personalities in pursuit of self-interest or personal gain .
\end{KeyConcept}

\begin{itemize}
    \item What does it mean? They will sell you their smile . They will sell you their glamour, their services, their expression --- to seek personal gain, whether a promotion or the reward for giving good service .
    \begin{itemize}
        \item Examples given: the cab driver who talks pleasantly about everything through the whole ride --- the economy, politics, everyone --- and at the end asks for a five-star rating . The delivery executive who asks for a rating at the door . A salesperson in a showroom who talks to you very sweetly; an air hostess who does the same . You would not see a software engineer or a mechanic talking to a customer like this --- only the service-class workers .
    \end{itemize}
    \item So in the language of C. Wright Mills these are the people who have become alienated from their own real selves --- because their smile is an expression of the reward they seek from you, in the form of a higher rank, a promotion, or good feedback from the customer . They are serving the interest of larger companies by selling pieces of their personalities .
    \item This is known as white-collar alienation .
\end{itemize}

\begin{CriticalPoint}
\begin{itemize}
    \item But was it a critique of alienation? No --- it is not a criticism of Marx .
    \item It is simply a refinement of Marxism in the modern era --- the application of the same concept to today's time . Application, not criticism .
\end{itemize}
\end{CriticalPoint}

\begin{QuickRevision}
\begin{itemize}
    \item Classical scholars cannot really be criticised --- critics end up using their concepts and remain inside their frame .
    \item C. Wright Mills: author of the sociological imagination; first to introduce Marxism in US universities (Marxism untaught there until the 1930s) .
    \item Marx analysed the blue-collar, manual, secondary-sector worker; Mills extends alienation to the white-collar service worker .
    \item White-collar workers prostitute their personalities in the sales room, board room, ball room and conference room --- selling smile, glamour and expression for personal gain .
    \item This is white-collar alienation --- an application and refinement of Marx, not a criticism .
\end{itemize}
\end{QuickRevision}

\subsection{Andr\'e Gorz and Herbert Marcuse --- Leisure Shaped by Capitalism}
\begin{itemize}
    \item During the era of Marx the working day was 12 to 16 hours .
    \begin{itemize}
        \item In fact the labour laws of the time said only that you could not make a child work for more than 12 hours --- so more than 12 hours was against child rights, and anything less was fine . Twelve hours was the minimum length of the working day .
    \end{itemize}
    \item When you are working 12 to 16 hours you would not have time for leisure activities --- no time for your own entertainment .
    \item Today the working days have become shorter --- 8 hours, 9 hours, at maximum 10 hours . Beyond that it becomes a court case, because we have labour laws (except where someone is operating in the informal economy --- but these scholars are writing about Europe) .
\end{itemize}

\paragraph{The argument}
\begin{itemize}
    \item At the surface it seems we now have a great deal of time for leisure . But just as the working day was shaped by capitalism, leisure has also been shaped by capitalism .
    \item After work, what are you going to do? Watch a streaming platform . Play games on a console . Watch videos online . Listen to music . These leisure activities have also been shaped by capitalism .
    \item And when you are watching these things, you are again working for these people --- because when you watch, they earn, and you do not get anything . So the vicious cycle is complete .
    \item How you are going to spend your time has also been shaped by capitalism . So even when you are out of work, out of the factory, you are still in the sphere of capitalism --- you cannot escape it .
    \item Instead of finding leisure in talking to your friends, your parents or your family members, you find leisure in these commodities . And this is how you have become one-dimensional .
    \begin{itemize}
        \item Ask anyone what their source of entertainment is, and the prompt reply is a streaming platform or a video site --- and that is what their entertainment is .
    \end{itemize}
\end{itemize}

\paragraph{Herbert Marcuse}
\begin{itemize}
    \item Herbert Marcuse says leisure is shaped by capitalism, and we have become passive robots chasing false needs . These are false needs .
    \item This is very much like commodity fetishism --- because both of these scholars are neo-Marxists .
\end{itemize}

\paragraph{Cultural industries and cultural goods}
\begin{itemize}
    \item These streaming platforms and video sites are cultural industries, and what comes out of them are cultural goods .
    \item Neo-Marxists say the cultural industry produces cultural goods, and no cultural good is very different from another --- all of them are the same, because it is a repetitive flavour .
    \item Capitalists have identified what kind of taste you have, so they repeat similar things: they introduce different characters, but with similar lyrics, similar scenes and a similar ending scene as well .
    \begin{itemize}
        \item The teacher admitted candidly that even he cannot tell the difference between two of the major streaming platforms --- a very basic question that still confuses him .
        \item Two doubts deferred: whether this amounts to surveillance (taken up later), and whether alienation is the same as McDonaldization (no --- taken up later) .
    \end{itemize}
\end{itemize}

\subsection{Welfare Capitalism, Keynes, Beveridge and the Welfare State}

\paragraph{The liberals' response to the responders}
\begin{itemize}
    \item The people responding to liberalism were Owen, Fourier, More, Saint-Simon and Marx . The liberals then responded to them:
    \begin{itemize}
        \item ``Your objectives are very nice . We really like your objectives . You want proper safe working conditions for the workers, you want to give them good education and good health . Ultimately it will benefit us only . We agree.'' 
    \end{itemize}
    \item So what they proposed was: let us have welfare capitalism .
    \item Welfare capitalism means the capitalist giving you education, good health, schooling, good residence --- good quarters, working-class quarters --- everything provided by the capitalist .
    \begin{itemize}
        \item Like the housing given by a large industrial house: a house within the complex, with water and electricity and everything else .
    \end{itemize}
    \item The purpose: so that they do not revolt, and work silently . ``We have given you everything, so now there is no need to bring a revolution --- just work.'' 
\end{itemize}

\paragraph{Why it failed --- the Great Depression}
\begin{itemize}
    \item But this welfare capitalist model was not sustainable . It failed --- the biggest failure was in the 1930s, the period of the Great Depression .
    \begin{itemize}
        \item At that point of time welfare capitalism was not purely welfare capitalism; it was largely capitalism with a little bit of welfare . And it collapsed in the 1930s .
    \end{itemize}
\end{itemize}

\paragraph{What led to the Great Depression}
\begin{enumerate}
    \item After the World War, America had become highly prosperous --- producing cars and much else .
    \item The cultural industry got established: radio, FM, music, Hollywood, cinema, and popular music forms .
    \item People had become consumers of these cultural goods and were ready to buy a car .
    \item But after a point, when France and others also recovered and started making these goods, the world was flooded with goods while demand did not keep pace with the production .
    \item Ultimately, with overproduction and under-consumption, the Great Depression came in .
\end{enumerate}

\paragraph{Why the revolution still did not happen --- John Maynard Keynes}
\begin{itemize}
    \item When the Great Depression came in 1930, people got highly excited --- the time of Marxism has come . Depression is the last thing in the laws; after this revolutionary consciousness would develop and these people would overthrow the capitalists .
    \item But it did not happen . Why? Only because of one man: John Maynard Keynes .
\end{itemize}

\begin{KeyConcept}
Keynesian economics overthrew Marxist economics . Keynesian economics was developed to fight this demand-deficit economy .
\end{KeyConcept}

\begin{itemize}
    \item And it is because of Keynes that we today study monetary policy and fiscal policy, and the idea of stimulus --- fiscal stimulus and monetary stimulus .
    \begin{itemize}
        \item Cut the repo rate; cut taxes under fiscal policy; inject money into the economy; cut the CRR . All of these come from Keynesian economics .
    \end{itemize}
    \item Keynesian economics simply tells us that when there is a demand deficit, we should work on increasing the demand --- and demand can be increased if people are given that much amount of money, in whatever form it may be .
    \item These ideas were much liked by the man who became President of America after the Great Depression --- Franklin Roosevelt . He transformed the American economy into a Keynesian economy and actually established Keynesian economics .
    \item And from there the drive of Marxism started weakening . Somehow it was believed that whenever the economic or business cycle goes through boom and bust, Keynesian economics would always fix the anomaly of demand deficit or supply deficit . So working-class consciousness would not develop, and the overthrow would never happen . These were the things that actually disrupted the idea of Marxism .
\end{itemize}

\paragraph{From welfare capitalism to the welfare state}
\begin{itemize}
    \item After Franklin Roosevelt came the idea of the welfare state . From welfare capitalism we moved towards the welfare state .
    \item By the 1960s the whole of Europe and the USA had absorbed this idea --- the welfare state had become the new normal .
    \item But this welfare state also got over: the demise of the welfare state in the 1970s . Why it collapsed is taken up later .
\end{itemize}

\paragraph{Keynes and Beveridge}
\begin{itemize}
    \item Keynes in America and Beveridge in Britain --- these are two people who were working towards a new system where economies remain stable and do not fluctuate the way they did during the World War or the Great Depression .
    \item To fight the Great Depression we had Keynes . To fight the deficits after the World War we had the Beveridge Report . These two people actually established the idea of the welfare state .
    \begin{itemize}
        \item A doubt clarified in passing: a software engineer is white collar; a hardware engineer is blue collar .
    \end{itemize}
\end{itemize}

\begin{QuickRevision}
\begin{itemize}
    \item Welfare capitalism = the capitalist supplies education, health, housing, so that workers do not revolt .
    \item It collapsed in the Great Depression of the 1930s --- post-war American prosperity, cultural industry, then global overproduction against insufficient demand .
    \item Keynes prevented the expected revolution: demand-deficit economics, monetary and fiscal policy, stimulus; adopted by Franklin Roosevelt .
    \item Keynes (America) and Beveridge (Britain) established the welfare state; it became the new normal by the 1960s and died in the 1970s .
\end{itemize}
\end{QuickRevision}

\subsection{A Doubt on Fascism}
\begin{itemize}
    \item First, a logical point: you cannot ask how an orange is different from a car . One is a fruit and one is a vehicle; there can be no comparison between them .
    \item Fascism is a political ideology . Capitalism is an economic ideology . So comparing fascism and capitalism is comparing two very different things .
\end{itemize}

\paragraph{What fascism is}
\begin{itemize}
    \item Fascism is an ultra-nationalist ideology --- in the sense that the whole of nationalism is subordinated to the interest of one leader; the nation as a concept serves one leader .
    \begin{itemize}
        \item Under Nazism, German nationalism was not over and above the leader --- the leader was over and above German nationalism . Therefore to be German became equal to being a supporter of that leader's version of Germany: the Nazi Germany that despised Jews, that wanted to establish its supremacy across Europe, to become prosperous across Europe and to expand beyond its confines .
    \end{itemize}
    \item Fascism and Nazism are different . Fascism was in Italy, under the leadership of Mussolini --- Italian nationalism subordinated to the interest of Mussolini, very much similar to Hitler's Nazism .
\end{itemize}

\paragraph{Gramsci and the Prison Notebooks}
\begin{itemize}
    \item Under fascism, whoever protested against Mussolini was jailed . One of those jailed was Antonio Gramsci, a neo-Marxist .
    \item Gramsci wrote \emph{The Prison Notebooks}, in which he explained how these totalitarian, fascist regimes are run .
    \item Gramsci's reading was also very instrumental in the development of critical theory --- Adorno and Horkheimer .
\end{itemize}

\begin{CriticalPoint}
\begin{itemize}
    \item A caution given on ideal types: these are concepts . The use of these ideal types is very wrong sometimes --- it becomes a journalistic use when people casually label a leader a fascist . A full treatment was deferred to the politics section .
    \item Context: after the World War people expected the rise of socialism, but what came instead was the rise of totalitarian leaders and, in America, the cultural industry . Dictatorship, fascism and Nazism are the ideologies which prevented the establishment of the new utopian world of socialism --- which is why people start equating fascism with capitalism .
    \begin{itemize}
        \item Also clarified: neo-Marxism is not only about criticising --- it is much more than that .
    \end{itemize}
\end{itemize}
\end{CriticalPoint}

\begin{QuickRevision}
\begin{itemize}
    \item Fascism is a political ideology, capitalism an economic one --- they are not comparable categories .
    \item Fascism = ultra-nationalism in which the nation is subordinated to one leader; Nazism (Germany) and fascism (Mussolini's Italy) are distinct .
    \item Gramsci, jailed under Mussolini, wrote \emph{The Prison Notebooks} on how totalitarian regimes are run; this fed into critical theory .
    \item Ideal types are analytical concepts, not journalistic labels .
\end{itemize}
\end{QuickRevision}

\sectiondivider

\section{The Criticism of Marx}

\subsection{Opening Doubts}

\paragraph{How exactly are the bourgeoisie alienated?}
\begin{itemize}
    \item If I am the owner of a company and I am not able to produce something out of my own wish --- if whatever I produce comes out of subjugation to market forces --- that is alienation .
    \item I have become subjugated to the demand-and-supply game; I have become subordinated to market supply-and-demand forces . I cannot do anything on my own .
    \begin{itemize}
        \item Illustration from the teacher's own trade: I may know I can teach the subject in ten days, but nobody will accept a ten-day programme as credible, so I have to run a three-and-a-half to four-month programme . If I drop the fee very low it will look cheap; if I raise it very high nobody will pay . So the price and the length are both decided by the market, not by me . I am alienated .
    \end{itemize}
\end{itemize}

\paragraph{Why is Marxism known as a romantic philosophy?}
\begin{itemize}
    \item Anything is romantic which seems rosy --- when you romanticise the idea of being an ambassador, or of being a district collector, that is romanticising .
    \item Many people romanticise Marx, and others therefore call Marxism a romantic philosophy --- in the sense that you can talk about it, but it is not going to happen . It will remain in the book, and that is why you can romance about it .
\end{itemize}

\paragraph{Apart from communism, is there any solution to alienation?}
\begin{itemize}
    \item Yes --- Marx gave a solution for alienation also .
    \begin{enumerate}
        \item Decrease the working hours . At that point of time the working day was 12 to 16 hours, so he suggested 8 to 10 --- drop it .
        \item The end of capitalism --- communism .
    \end{enumerate}
    \item But Andr\'e Gorz and Herbert Marcuse say that even though working hours have decreased, the leisure we have gained is something we have become trapped in . It is also an illusion; it is not leisure; it is also false needs that we are chasing . Leisure is shaped by capitalism, and we have become alienated from leisure also . So the circle is complete .
\end{itemize}

\paragraph{Was Marx vindictive because of his own suffering, or did he really want equality?}
\begin{itemize}
    \item Look at the idea itself . Socialism will come . What is the purpose of socialism? To distribute the goods and resources, to minimise all inequality, to end all the oppression of the working class that was happening in capitalism .
    \item Once that is done, the work of the state is over . Who is going to distribute the goods and resources? The working class only --- they will decide their own fate . The capitalists have had their turn .
    \item Once distribution is over, the state shall wither away, and then we have communist society .
\end{itemize}

\begin{QuickRevision}
\begin{itemize}
    \item Bourgeois alienation: the owner cannot produce out of his own wish either --- the market dictates product, price and format .
    \item Marxism as ``romantic philosophy'': attractive to talk about, not going to happen .
    \item Marx's own solutions to alienation: shorter working hours, and communism . Gorz and Marcuse close that exit --- leisure too is shaped by capitalism .
    \item The purpose of socialism is distribution and the ending of oppression; then the state's work is over and it withers away .
\end{itemize}
\end{QuickRevision}

\subsection{Criticism of Alienation --- Max Weber}

\begin{CriticalPoint}
A general point about how criticism works among these scholars: you would not find them mentioning the name of the person explicitly in their books . You will not see Weber writing Karl Marx's name while criticising him .
\end{CriticalPoint}

\begin{itemize}
    \item Max Weber says alienation is something that you cannot escape . Our societies are going to be more and more rationalised --- we are going towards a process called rationalisation .
\end{itemize}

\paragraph{What rationalisation means (in brief)}
\begin{itemize}
    \item Today we talk about numbers, rules and laws . Everything is quantifiable today . Everything is bound by rule today . Gone are the days of tradition, religious beliefs, customs and norms; everything today is either rule-bound or quantifiable .
\end{itemize}

\begin{table}[h!]
\centering
\begin{tabularx}{\textwidth}{lYY}
\toprule
\rowcolor{AccentDark}
\thead{} & \thead{Two hundred years back} & \thead{Today} \\
\midrule
Exercise & After half an hour of jogging you used to feel fresh & After half an hour of jogging you check how many calories you have burnt \\
Greeting & You would greet people --- ``hello, how are you?'' & You would rate them --- 3 out of 10, 4 out of 10 \\
Education & Qualitative --- theological and religious education & Quantifiable --- what percentage did you get in class 11 and 12; a numbered curriculum from class 1 to 12 \\
\bottomrule
\end{tabularx}
\end{table}

\begin{itemize}
    \item Now, when everything is rule-bound, number-bound and quantifiable, there is no space for human emotions . Rationality displaces emotion .
\end{itemize}

\paragraph{Bureaucrats as the highest manifestation of rationality}
\begin{itemize}
    \item Max Weber says that bureaucrats --- the administrative and police service officers --- are the highest manifestation of rationality .
\end{itemize}

\begin{ExampleBox}
\begin{itemize}
    \item Imagine a poor person who lost everything in the lockdown and wants a licence to open a restaurant . He goes to the district collector and pleads: all my wealth is gone, I have nothing in my hands, my family will die hungry, please understand, I really need this licence .
    \item The collector says: wait for one or two minutes, I have an urgent meeting . The man waits --- and the collector returns after an hour . Then the collector does not even look at him; he is on a call with another officer, quoting rule six .
    \item Finally: ``many people have applied for a licence; you can apply for a licence.''  There is no emotion .
    \item But the collector is also helpless . He is a mouthpiece, a spokesperson of the civil service code, and cannot do anything beyond that code . If the person is dying, the officer will say: write an application stating that you are dying .
    \item Everything has a process, a rule, a procedure; everything has to be in written form; the paper travels a long chain of higher authorities; then the rule book is opened to see which rule fits the licence .
\end{itemize}
\end{ExampleBox}

\paragraph{The line}
\begin{KeyConcept}
Max Weber: bureaucrats are ``specialists without spirit, sensualists without heart.'' 
\begin{itemize}
    \item That means the very idea of being impersonal and rational has alienated those people from emotions, empathy and sympathy . All those values are gone; only rule, rationality, process and procedure remain .
\end{itemize}
\end{KeyConcept}

\begin{itemize}
    \item Remember: Weber is not talking about the alienation of the working class . He is talking about the alienation of humanity . All people are going to act like this; the larger trend is the rationalisation of societies, and bureaucrats are the manifestation of that rationality .
    \item The larger trend of rationalisation will eliminate everything magical, spiritual, emotional and humane in humanity .
    \item Impersonality is one of the features of bureaucracy --- you cannot afford to get emotional with the client, you have to remain strictly rule-bound, you change your personality when you enter service . And this is nothing but alienation --- alienation from human emotions, empathy and sympathy .
\end{itemize}

\begin{QuickRevision}
\begin{itemize}
    \item Weber: alienation cannot be escaped; societies move towards rationalisation .
    \item Rationalisation = everything quantifiable and rule-bound; no space for emotion .
    \item Bureaucrats are the highest manifestation of rationality --- ``specialists without spirit, sensualists without heart'' .
    \item Weber generalises alienation from the working class to humanity as such .
    \item Impersonality, a defining feature of bureaucracy, is itself alienation from empathy and sympathy .
\end{itemize}
\end{QuickRevision}

\subsection{Robert Blauner --- Alienation is Relative, Not Universal}
\begin{itemize}
    \item Blauner's starting objection: all the Marxist scholars say we are alienated, alienation will definitely be there, we are helpless, nothing is going to happen . Blauner says: at least ask those people once .
    \begin{itemize}
        \item Are they really alienated? Or are you going to sit on your chair and comment on them relentlessly? First tell them the definition of alienation, and then ask them: are you feeling this? 
    \end{itemize}
\end{itemize}

\begin{KeyConcept}
Robert Blauner's statement: alienation is not a universal phenomenon . Rather, it is a relative phenomenon .
\end{KeyConcept}

\begin{itemize}
    \item The reasoning: all industries are not the same . We have the textile industry, the print industry, the chemical industry, the automobile industry --- all these industries are different from one another .
    \item How can you say that in all industries alienation is there? In some industries maybe alienation is not there; in some, some people may be alienated and some may not .
    \item What does this show? That Blauner is not a structuralist . He wants to ask the people themselves --- as a bureaucrat, as a worker in any factory, are you really alienated? He wants to check it himself; he wants to conduct a survey .
\end{itemize}

\paragraph{The survey}
\begin{itemize}
    \item Blauner conducted a behavioural attitude survey in four industries: craft (print), chemical, textile and automobile .
\end{itemize}

\paragraph{The problem of measurement --- operationalisation}
\begin{itemize}
    \item How do you measure alienation? There is no scale . Alienation is not a quantifiable thing .
    \item When you decide to measure something that does not seem measurable, social scientists use one word: they operationalise the concept .
    \item Operationalisation of the concept means you establish certain criteria of alienation . If those criteria are met, then the person is alienated . This is the way we measure alienation .
\end{itemize}

\paragraph{The four dimensions Blauner used}
\begin{enumerate}
    \item Powerlessness --- ``nobody listens to me; I simply obey other people . Although they have told me I am an executive manager, nobody listens to me.'' 
    \item Meaninglessness --- ``I keep doing it only because the money comes at the end of the month . Otherwise I do not find any meaning in the work that I am doing.'' 
    \item Self-isolation (self-estrangement) .
    \item Normlessness --- there is no system, no proper rule or governance in place . You believed you were the one to be promoted to the managerial level, but someone very close to the boss got the promotion instead .
\end{enumerate}

\paragraph{The result}
\begin{itemize}
    \item After collecting the data from all four industries, he plotted degree of alienation against the different industries --- print, textile, automobile, chemical --- and formed an inverted U-curve .
    \item Conclusion: in the automobile sector alienation is extreme; in the print or craft sector alienation is least .
    \begin{itemize}
        \item Doubt raised: if the names of the industries change, there will be no inverted U-shape --- for instance if you put a toy industry or restaurants in place of print . Answer: obviously; but he analysed only these four industries .
    \end{itemize}
\end{itemize}

\begin{QuickRevision}
\begin{itemize}
    \item Blauner: alienation is relative, not universal --- it varies by industry, so ask the workers .
    \item Not a structuralist; he surveys behaviour and attitude in four industries: craft/print, chemical, textile, automobile .
    \item Alienation is unmeasurable directly, so it is operationalised through four dimensions: powerlessness, meaninglessness, self-isolation, normlessness .
    \item Result: an inverted U-curve --- extreme alienation in automobile, least in print/craft .
\end{itemize}
\end{QuickRevision}

\subsection{Goldthorpe and Lockwood --- Ask What Meaning They Attach}
\begin{KeyConcept}
Line to write: instead of categorising the working class as victims of alienation, the job of a social scientist should be to see what meaning they attach to their work .
\end{KeyConcept}

\begin{itemize}
    \item You cannot simply say all of them are alienated . Ask them: how exactly does this job matter to you? Is it only a source of income to you, or is it something more than that? 
    \begin{itemize}
        \item Think about people working in a large, respected industrial house --- they feel proud of being an employee there, because the whole ecosystem of that house is very different from that of other companies . They feel they will do anything for the company, so there is no scope for them to feel alienated .
        \item Think about people in a big software firm, sitting on bean bags, working on their laptops, enjoying it --- they feel even better there than at home, and that is why they are considered highly productive also . Mutual needs are being exchanged .
    \end{itemize}
    \item Of course it depends which particular section of worker you are talking about . But if it depends, that simply means alienation is not universal .
    \item So instead of simply calling them victims of alienation, we should ask them . This is what Goldthorpe and Lockwood are saying .
\end{itemize}

\begin{QuickRevision}
\begin{itemize}
    \item Goldthorpe and Lockwood: study the meaning workers attach to their work rather than assigning them the label of alienation .
    \item Every company has its own ecosystem; pride and mutual exchange can displace alienation .
    \item If alienation depends on the case, it is by definition not universal .
\end{itemize}
\end{QuickRevision}

\subsection{Simone de Beauvoir and Arlie Hochschild}

\paragraph{Simone de Beauvoir}
\begin{itemize}
    \item Simone de Beauvoir simply argues that it is the women who are alienated --- it is not the working class .
    \item Alienated from what? From the species being --- from their true potential .
    \item Women have a huge potential to transform the world .
\end{itemize}

\begin{DataFact}
The figure cited in class: according to the World Economic Forum, the World Bank and similar organisations, if the workforce gets feminised --- if women are brought into the workforce --- the GDP of the nation will grow by 50\% of what it is today in the next ten years [figure as stated in lecture] .
\end{DataFact}

\begin{itemize}
    \item And what are they doing? They are at home taking care of the father, the husband, the child . The true creativity and the potential that they have --- they have been alienated from it .
    \begin{itemize}
        \item Women can fly fighter jets . Now imagine a woman who has spent all her life cooking rice and dal --- that means she has been alienated from her true potential, from the fact that she could do anything .
    \end{itemize}
    \item And it is only because of patriarchy that women have become alienated . Patriarchy gained its power and privilege by subjugating women to its interests . How does man gain power? Only by suppressing women . And with the suppression of women comes her alienation .
\end{itemize}

\paragraph{Arlie Hochschild --- emotional labour}
\begin{itemize}
    \item Arlie Hochschild has given one concept: emotional labour .
    \item Emotional labour is more laborious than physical labour . It is mentally tortuous .
    \begin{itemize}
        \item To keep smiling every day is not an easy job . A physical task is still easy; but to communicate 24$\times$7 --- ``hello sir'' --- and to talk in a sweet fashion, is very difficult . Think of the customer-care call where the customer speaks in one language and the agent must answer sweetly in another throughout .
    \end{itemize}
    \item Hochschild says emotional labour is the source of alienation in the service sector --- very much similar to C. Wright Mills's white-collar alienation, but she is specifically talking about women .
    \begin{itemize}
        \item Note running through all of these: none of them are actually critics . They are applications and extensions .
    \end{itemize}
\end{itemize}

\begin{QuickRevision}
\begin{itemize}
    \item De Beauvoir: it is women who are alienated --- from species being and true potential --- and patriarchy is the cause .
    \item Feminisation of the workforce as the counterfactual (GDP figure as stated in lecture) .
    \item Hochschild: emotional labour --- harder and more mentally tortuous than physical labour --- is the source of service-sector alienation, especially for women .
    \item Parallel with Mills's white-collar alienation, but gendered .
\end{itemize}
\end{QuickRevision}

\subsection{Critique of Class in Itself / Class for Itself --- Vladimir Lenin}

\paragraph{Recap of the problem}
\begin{itemize}
    \item Class in itself: people having similar relation to the means of production . The teacher of sociology and the worker making street food at a large sweet-shop chain are the same --- because both are not owners of the means of production; both are only selling their labour .
    \begin{itemize}
        \item And yet we still look down on that worker --- ``how much would he earn in a day?''  Marx would abuse me for that: who are you to say anything to that person? You belong to that class only .
    \end{itemize}
    \item When would I get this realisation? Only when exploitation attains a level at which the revelation of the inherent contradiction becomes inevitable and the polarisation becomes inevitable --- until there are clearly only two kinds of people on earth: the owners, and those who are not owners and work for them .
    \item When exploitation takes us to that level, this exploited class will form a union and overthrow these people --- class for itself . But it does not happen like this . The world is not so easy .
\end{itemize}

\paragraph{The Russian problem}
\begin{itemize}
    \item Vladimir Lenin wanted to mobilise the working class --- the peasants --- in Russia .
    \item In Russia there was no working class . The industrial revolution came later . The major class that was fighting was the peasant class --- it was landlord and serf .
    \item Lenin believed that these peasants cannot bring any revolution by themselves . They are protesting and agitating, but everything is in a very individual fashion . They need to be organised together .
    \item Unless they are organised, unless they are mobilised, the revolution would not happen --- their agitation and protest will be loosely taken, or will seem very individualistic .
    \begin{itemize}
        \item Analogy: if you want to protest against a decision of a university, you alone cannot do anything unless it is a mass, mobilised protest, covered by the media and given a lot of footage over ten days . Until that happens, no change comes .
    \end{itemize}
\end{itemize}

\paragraph{The All-Russia newspaper and the Vanguard Party}
\begin{enumerate}
    \item First, these peasants have to be imparted the consciousness that they are exploited .
    \item For that, Lenin talked about an All-Russia newspaper . In that newspaper different editors would write, every day, about how the working class is exploited --- and they would write in the language legible to those peasants .
    \item The peasants read the newspaper every day and then discuss it with their friends .
    \item Gradually a Vanguard Party --- a political party --- will be established: the party of the peasants, or the party of the workers .
    \item This Vanguard Party will bring forward the revolution, working to bring consciousness among these people that they are exploited .
    \begin{itemize}
        \item The comparison drawn in class: this is what the Communist Party of India was doing initially --- going into the affected districts and telling people about the oppression of the state, about land distribution and land reform .
    \end{itemize}
    \item These are the people who impart consciousness among the people, because they believe that people are not going to become conscious all by themselves . An external party has to be there to mobilise them around a similar conscious feeling .
\end{enumerate}

\begin{KeyConcept}
Lenin: the Vanguard Party will help in imparting consciousness among the oppressed . That means for a class in itself to become a class for itself, you need not just homogenisation, depression and pauperisation --- you also need one political party . It is not that easy .
\end{KeyConcept}

\paragraph{Why this is not a critique either}
\begin{itemize}
    \item Lenin was also a Marxist --- and that is why the ideology became Marxist--Leninist .
\end{itemize}

\begin{table}[h!]
\centering
\begin{tabularx}{\textwidth}{lYY}
\toprule
\rowcolor{AccentDark}
\thead{} & \thead{Marxism} & \thead{Marxism--Leninism} \\
\midrule
Born from & The factories & The agrarian land \\
Class concerned & The industrial working class of the Industrial Revolution & The peasants of an agrarian society \\
Consciousness & Will develop by itself --- it is inevitable; no external agency needed & Must be imparted by an external political party --- nothing happens unless we act upon it \\
Political party & Marx never believed in establishing one; he held democracy to be a bourgeois ideology & The Vanguard Party is the essential addition \\
\bottomrule
\end{tabularx}
\end{table}

\begin{itemize}
    \item The act is familiar, the discourse is familiar, the oppression is familiar, the concerns are familiar, the classes are almost the same --- only the characters have changed .
\end{itemize}

\begin{QuickRevision}
\begin{itemize}
    \item Class in itself only becomes class for itself when exploitation and polarisation become undeniable --- but Marx's automatic route does not work .
    \item Russia had no working class; the fighting class was the peasantry, and it was unorganised .
    \item Lenin: All-Russia newspaper written in the peasants' own language $\rightarrow$ discussion $\rightarrow$ Vanguard Party $\rightarrow$ revolution .
    \item A fourth prerequisite is added to Marx's three: a political party .
    \item Marxism came out of the factory, Marxism--Leninism out of the land; Marx rejected the party, holding democracy to be a bourgeois ideology .
\end{itemize}
\end{QuickRevision}

\subsection{Nancy Hartsock --- Women as Commodities}

\paragraph{The circuit of commodity, recalled}
\begin{itemize}
    \item $C_1 - C_2$, then $C_1 - M - C_2$, then $M_1 - C - M_2$ . Initially we exchange commodities; then money gets introduced and the commodity is sold for money, which helps in buying another commodity; and now money has become the end . The commodity is not the end . Investment, reinvestment .
\end{itemize}

\paragraph{The argument}
\begin{itemize}
    \item Nancy Hartsock, a feminist, wrote an article titled ``Women and/as Commodities'' .
    \item She says it is the women who are commodities . Just as a commodity has both use value and exchange value, women also have both use value and exchange value .
\end{itemize}

\begin{table}[h!]
\centering
\begin{tabularx}{\textwidth}{lYY}
\toprule
\rowcolor{AccentDark}
\thead{} & \thead{Use value of women} & \thead{Exchange value of women} \\
\midrule
What it is & Services that are useful but not traded for money --- the caring services: caring for the child, looking after the elderly, waiting for the husband, cooking food for the family & The woman acquires exchange value on the occasion of marriage --- she is exchanged for two things: money and marriage \\
Charge & She is not charging anything for these services & The exchange market: bride price, dowry, trafficking, the trophy wife \\
\bottomrule
\end{tabularx}
\end{table}

\begin{CriticalPoint}
The very concept of family has come into existence only because of the exchange of women . If women were not exchanged, the family would not have been possible .
\end{CriticalPoint}

\paragraph{The manifestations she lists}
\begin{itemize}
    \item Bride price in some regions --- that means it works like a commodity .
    \item Dowry --- a manifestation of the fact that the woman is a commodity; that is why she is traded for money .
    \item Trafficking of women --- one of the biggest international crimes; women trafficked across borders through networks .
    \item The trophy wife --- an eighty-year-old man about to die marries a twenty-five-year-old . Why? He has gained status . That woman is a commodity to him, because by keeping that commodity in his house he symbolises status . The existence of the human is reduced to the level of a commodity .
\end{itemize}

\begin{KeyConcept}
A note on how sociologists speak . When sociologists speak like this, they are not criticising us . They simply reveal to us something that is hidden --- something not visible to common-sensical observation . You knew all those things, but the way it is told was not known to you . This is the art of sociology: defamiliarisation . You become defamiliar with something that was familiar to you, and you start looking for the strange in seemingly familiar scenes .
\end{KeyConcept}

\paragraph{The circuit applied to the gender wage gap}
\begin{itemize}
    \item Hartsock says $C_1 - M - C_2$ is a highly patriarchal notion, and it is applicable to the gender wage gap .
    \item For women, what actually happens is $C_1 - M(\text{minus}) - C_2$ . A man and a woman are both working in the same organisation; for the same work the man gets more money and the woman gets a lower wage .
    \item The man's labour power exchanged for money lets him buy a house . Her labour power, sold for less, does not --- so she will have to remain dependent on him .
    \item The whole gender wage gap is a constructed notion through which these patriarchal marriage structures are maintained, so that women are always dependent on the male .
\end{itemize}

\begin{QuickRevision}
\begin{itemize}
    \item Hartsock, ``Women and/as Commodities'': women too have use value (unpaid caring services) and exchange value (marriage, money) .
    \item The family exists because women are exchanged .
    \item Manifestations: bride price, dowry, trafficking, the trophy wife --- the human reduced to a commodity .
    \item Defamiliarisation is the art of sociology: making the familiar strange .
    \item $C_1-M(\text{minus})-C_2$: the gender wage gap is constructed to keep women dependent within patriarchal marriage .
\end{itemize}
\end{QuickRevision}

\subsection{Nancy Hartsock --- Women as the Real Source of Surplus Value}

\paragraph{First, why machines cannot be the source of surplus}
\begin{itemize}
    \item Some students were confused about why machines cannot give surplus, since so many things come out of them and money keeps flowing .
    \item Look at it practically . A computer lying idle on a table is dead labour . The moment you sit at the chair and operate it, pages come out of the printer, and those pages are sold for a thousand rupees .
    \item What is the source of surplus value? You --- not the computer . The computer will not produce a page by itself . And the valuation of what comes out depends on the work you did, not on the fact that it came from a computer .
    \item And by the way, who made the computer? Labour only . So the source of the fixed capital is also variable capital . Ultimately it is variable capital alone which is the source of all the surpluses .
\end{itemize}

\paragraph{The argument}
\begin{KeyConcept}
Nancy Hartsock says the real source of surplus value is women .
\end{KeyConcept}

\begin{enumerate}
    \item Before you sit at the computer and expend the labour power that generates the surplus, you must have eaten something; someone must have taken a great deal of care of you --- breakfast, lunch, care, love, affection .
    \item Hence the saying that behind every successful man there is a woman . That man is highly productive in the office only because he was taken care of very well back at home . Unless you are peaceful at home your mindset will not be right --- how can you be a productive labourer in the office? 
    \item So the real source of surplus value is that woman's labour power --- this labour power has generated the other labour power ($LP$ has generated $LP'$) .
    \item It is the women who give workers to the society . If workers are productive, their productivity is the result of the reproduction process that all workers have gone through .
    \item So reproductive labour is not rewarded --- only productive labour is rewarded . And ``reproductive'' does not mean only biological reproduction; it means care as well .
\end{enumerate}

\begin{itemize}
    \item And that is why we have the debate on wage for housework . For three or four months it became a dominant discourse and the newspapers were flooded with articles on whether it should be given; some argued that it definitely should be .
\end{itemize}

\begin{QuickRevision}
\begin{itemize}
    \item Machines cannot be the source of surplus: the idle computer is dead labour, and labour made the computer in the first place .
    \item Hartsock: the real source of surplus value is women's labour --- the care and reproduction that makes the productive worker possible .
    \item Reproductive labour (biological and care work) is unrewarded; only productive labour is rewarded .
    \item Hence the wage-for-housework debate .
\end{itemize}
\end{QuickRevision}

\subsection{Critique of Scientific Socialism --- Does Socialism Exist?}

\paragraph{The sequence did not hold}
\begin{itemize}
    \item Marx said the history of society has been the history of class struggle: primitive era $\rightarrow$ ancient society $\rightarrow$ feudal society $\rightarrow$ capitalist society $\rightarrow$ socialism $\rightarrow$ finally communism .
    \item Now look at Russia, India and China . They had feudalism, and they jumped directly to socialism --- with no capitalism in between .
    \item The Chinese revolution and the Russian revolution were both agrarian revolutions --- peasant-class revolutions, not working-class revolutions .
    \item What does it show? That history is not so unilinear or predictable . In empirical observation, the sequence Marx laid down was not followed in either of these societies .
\end{itemize}

\paragraph{Do we really have socialism today? --- the case of China}
\begin{itemize}
    \item China is a socialist society with a capitalist economy and a communist party . All three at once .
    \item China calls itself a socialist society . Being a socialist society, your duty should be to be concerned about the welfare of the oppressed and marginalised sections --- equality . As a socialist society China should be producing socially useful goods --- good health, education, public goods .
    \begin{itemize}
        \item But, as stated in class: inequality in China is even bigger than that of India; the rate of poverty is increasing; dropouts from school are increasing; people are not getting proper education; malnutrition and health issues are worse . Only during the pandemic did China gain --- and that too relatively, because everybody else was almost in depression . [claims as stated in lecture] 
    \end{itemize}
    \item That means it is not socialist in that sense .
\end{itemize}

\paragraph{When are you actually categorised as socialist?}
\begin{itemize}
    \item It is not enough that you own the means of production --- the strategic sectors: rail, power, nuclear, petrochemicals, atomic energy, defence .
    \item You should also distribute all the benefits . Just owning the means of production is not going to help; it does not give you the status of being socialist . The next duty is to distribute the resources --- which is not being done .
    \item What is the revenue being spent on instead? Warfare, defence goods, and the expansion of territory --- the China--Pakistan Economic Corridor, a ``look west'' / pivot-to-Eurasia policy (just as the United States once had a pivot to Asia), and One Belt One Road . They want to expand; they are not concerned about the poor people; they are only concerned about making wealth .
    \item Dumping goods all over the world at the cheapest prices: if you dump goods at the cheapest price, what wage will your worker get? A worker who should have got a hundred will get fifty --- because the main idea is to establish hegemony worldwide through ``made in'' products . So you would not be giving enough wage to the workers . That is not what a socialist economy is about .
\end{itemize}

\begin{KeyConcept}
Lynette Ong: China is a capitalist country with a socialist facade . It is like Adam Smith wearing the suit of Karl Marx .
\end{KeyConcept}

\paragraph{Power is also a resource}
\begin{itemize}
    \item The Communist Party has usurped power . Power is also a resource --- and it is not distributed . No one outside the party can be in power, and even within it the chance is small .
    \item So power is not distributed, wealth is not distributed, and poverty and inequality are rising . In no way should it be called socialist that way .
    \item The only reason for calling China socialist would be that the state is the owner of the strategic sectors . But if China is socialist because of that, China is also non-socialist, because the wealth generated from those enterprises is not distributed across the population --- how else would you explain rising poverty? 
    \item So China is socialist and non-socialist both . It depends on what is important to you --- only the ownership, or also what happens after the ownership .
\end{itemize}

\begin{DataFact}
\begin{itemize}
    \item China became a communist nation in 1949 .
    \item In 2004 China recognised private property in its constitution . If private property is also recognised, then what is socialist? 
\end{itemize}
\end{DataFact}

\paragraph{The alternatives are no different}
\begin{itemize}
    \item J. K. Galbraith: under capitalism, man exploits man; and under communism it is just the opposite --- which, as the teacher noted, is not really opposite at all, because the opposite of ``man exploits man'' is still ``man exploits man'' .
    \item The fact is that the alternatives beyond capitalism are no different from capitalism . You call yourself socialist or communist, but your practices are highly capitalist .
    \item Therefore Francis Fukuyama, in his book \emph{The End of History}, believed that capitalism is the end of history --- because the alternatives beyond capitalism are no different from capitalism . (Randall Collins was named alongside in the same connection.) 
    \item Everywhere you can see $M - C - M_2$ . Money is generated and not distributed; money is generated and then appropriated . Earlier we had appropriation of money by the capitalist class; today we have appropriation of money by the state .
    \item And in those times you could escape one capitalist by joining some other company . Today you cannot escape the state --- and therefore the state is the super bourgeois . You can move from one company to another, but you cannot leave the country .
\end{itemize}

\begin{QuickRevision}
\begin{itemize}
    \item Russia, India and China jumped from feudalism to socialism without capitalism, and through agrarian, not working-class, revolutions --- history is not unilinear or predictable .
    \item Socialism requires distribution, not merely state ownership of strategic sectors .
    \item Lynette Ong: a capitalist country with a socialist facade --- Adam Smith wearing the suit of Karl Marx .
    \item Power is an undistributed resource; private property recognised in 2004 .
    \item Galbraith's epigram; Fukuyama's \emph{End of History}: the alternatives are no different from capitalism .
    \item Appropriation has moved from the capitalist to the state --- the state as super bourgeois, because you cannot exit it .
\end{itemize}
\end{QuickRevision}

\subsection{Critique of the Mode of Production --- Anthony Giddens}
\begin{itemize}
    \item According to Marx, history has been the history of class struggle, and class struggle is the motor of history; societies evolve only because of the struggle between classes, and cannot evolve if classes do not struggle .
    \item Anthony Giddens, in \emph{A Contemporary Critique of Historical Materialism} (described in class as an article rather than a book per se), says there is no absolute universal formula that with class struggle history will evolve .
    \item And it is not the only struggle in the world . There are many other struggles which Marx has ignored .
\end{itemize}

\begin{ExampleBox}
\begin{itemize}
    \item Look at the partition of India and Pakistan . Was it class struggle? Was it the result of bourgeoisie versus proletariat, or lords versus serfs? No . It was about religious communities, not about a Hindu bourgeoisie and a Muslim proletariat .
    \item The World Wars were not bourgeoisie versus proletariat --- these were nation-states at war .
\end{itemize}
\end{ExampleBox}

\begin{itemize}
    \item The fact is that class struggle is only one of the struggles . There are many other forms of struggle Marx did not take into consideration .
\end{itemize}

\paragraph{Fault line conflict}
\begin{itemize}
    \item One of the other forms is fault line conflict --- conflict on the lines of ethnicity and religion .
    \begin{itemize}
        \item The Israel--Palestine conflict, Jews versus Arabs, is not a struggle of bourgeoisie versus proletariat .
    \end{itemize}
    \item Ethnicity is a cultural identity . Son-of-the-soil movements are the clear case: in the north-east, people from elsewhere are considered outsiders . With limited resources, if the local people are not served first and outsiders come and take away jobs and resources, it automatically leads to conflict between locals and outsiders --- whoever that outsider may be .
    \begin{itemize}
        \item In that conflict the seven states unite . It is not the mobilisation of proletariats --- it is the mobilisation of a region against outsiders . Two classes clearly emerge: insiders and outsiders, purely on ethnic lines .
    \end{itemize}
    \item The concept of fault line conflict is given by Samuel Huntington in his book \emph{Clash of Civilizations} .
    \item The India--China war, the India--Pakistan wars, Doklam, surgical strikes --- these are not class struggles . They are ethnic struggles or nation-state struggles . Not even religious struggles: it is not Buddhism fighting Hinduism, or Confucianism and Taoism fighting Sikhism or Islam .
    \begin{itemize}
        \item And ethnicity cuts across the nation: if I trace my blood ancestry to Lahore, then a connection with that side exists for me . On the lines of nation they are divided; on the lines of ethnicity they are united .
    \end{itemize}
\end{itemize}

\paragraph{How the world we live in was actually produced}
\begin{itemize}
    \item The world history we are part of today is not the result of class struggle . We adopted a socialist pattern of society only because of the World War .
    \item After the World War the world was divided into two blocs --- capitalist and communist --- and Nehru had to choose one flavour, and chose the one that seemed better and closer .
    \item So what we have is Fabian socialism --- which is nothing but democratic socialism . Both are one and the same; Fabian is a British term, there is no difference between them .
    \begin{itemize}
        \item How democratic socialism differs from Marxian socialism you already know: democracy is there; the state will be there forever, not only until inequalities are removed; and we achieved socialism not through violent struggle by overthrowing the bourgeoisie .
    \end{itemize}
\end{itemize}

\begin{QuickRevision}
\begin{itemize}
    \item Giddens, \emph{A Contemporary Critique of Historical Materialism}: no universal formula that history evolves through class struggle .
    \item Class struggle is only one struggle among many; partition and the World Wars are not class conflicts .
    \item Fault line conflict --- conflict along ethnic and religious lines --- is Samuel Huntington's concept from \emph{Clash of Civilizations} .
    \item Ethnicity is a cultural identity: son-of-the-soil movements produce insiders and outsiders, not bourgeoisie and proletariat .
    \item India's socialist pattern came from the post-war bloc choice, not from class struggle: Fabian = democratic socialism .
\end{itemize}
\end{QuickRevision}

\subsection{Ernesto Laclau and Chantal Mouffe}
\begin{itemize}
    \item Ernesto Laclau and Chantal Mouffe say that it is not important to predict what will happen in the future --- as Marx predicts that socialism is inevitable and will come with the demise of capitalism .
    \item Instead of predicting the future and presenting yourself as a prophet, we should rather be interested in studying why some revolutions are happening only in some places .
    \begin{itemize}
        \item Why are feminist revolutions happening in some places and not in others?  Why did the civil rights movement happen in some places and not others?  Why do environmental movements happen in some places and not others? 
    \end{itemize}
    \item This is what sociologists should study: why movements are localised and not universal . Marx says workers of the world unite --- but why did workers unite only in a few pockets of the world and not the world over? What prevents them from uniting? 
\end{itemize}

\paragraph{Their three points}

\paragraph{1. For every revolution to happen, there has to be a discourse --- and it has to be a democratic discourse}
\begin{itemize}
    \item Very important line: for every revolution to happen there has to be a discourse, and the discourse has to be a democratic discourse .
    \begin{itemize}
        \item Before the 1700s, the relation between man and woman, and the subordination of women, was seen as normal . It was very normal for a woman to obey the man; there was no question about it at all . If in the 1700s you had talked about feminism, that would have been a problem .
        \item But now you can easily speak as a feminist on social media or anywhere you want --- because we have the discourse available . The discourse of gender equality exists; there is a widely accepted notion that women and men are equal, and that equality of opportunity should be given irrespective of sex . And if a democratic discourse is not available, revolutions will not happen .
    \end{itemize}
    \item And when the democratic discourse is there, there is no need to bring forward a violent revolution . People will write editorials; students will make notes of them and reproduce them; some will clear the civil services writing essays on them; and then an officer will bring that change in her district . It happens peacefully, through the democratic process .
\end{itemize}

\paragraph{2. Revolutions are not limited to class struggle}
\begin{itemize}
    \item We have movements related to environmental rights, human rights, women's rights, civil rights, and racial and ethnic movements . We have LGBT and sexual-identity movements .
    \item We have movements on the lines of every identity you can think of --- sexual identity, women's identity, ecological and eco-feminist and environmental movements . So we have multiple kinds of movements and multiple kinds of struggles . It is not just class struggle and capitalism that everybody is concerned with --- capitalism is just one thing .
\end{itemize}

\paragraph{3. Do empirical study instead of prediction}
\begin{itemize}
    \item Instead of making predictions about the future of history or the end of history, the job of a sociologist should be to conduct empirical study on why revolutions happened in a few pockets and why they did not happen in others .
    \begin{itemize}
        \item For example: why did feminism develop in Europe and not in India? And how is feminism in India different from feminism in Europe? This is what sociologists should study --- instead of making predictions .
    \end{itemize}
\end{itemize}

\begin{QuickRevision}
\begin{itemize}
    \item Laclau and Mouffe: study why revolutions are localised rather than predicting an inevitable future .
    \item Condition for revolution: an available democratic discourse --- feminism before the 1700s was unsayable; today it is ordinary .
    \item Where democratic discourse exists, change comes peacefully rather than through violent revolution .
    \item Revolutions are not limited to class: environmental, human rights, women's, civil rights, racial, ethnic, LGBT movements .
    \item The sociologist's job is empirical study of where movements do and do not occur .
\end{itemize}
\end{QuickRevision}

\subsection{J\"urgen Habermas --- Legitimation Crisis}
\begin{itemize}
    \item J\"urgen Habermas is still alive; he is a second-generation neo-Marxist belonging to critical theory .
    \item He says: I am surprised that capitalism is still alive --- despite the World Wars, the rise and fall of fascism and Nazism, the crash of Wall Street in 1929, the Great Depression of 1930, the subprime crisis of 2008, the Occupy Wall Street movement, and the lockdown . Why does this wall not fall? 
    \begin{itemize}
        \item The teacher's own aside: it is not a very great answer --- but the biggest answers come only from the biggest questions .
    \end{itemize}
\end{itemize}

\paragraph{His answer}
\begin{KeyConcept}
The problem is that today the state will never let capitalism fall . What we have today, in the name of capitalism, is state-regulated late capitalism .
\end{KeyConcept}

\begin{itemize}
    \item Whenever the economy goes through a bust, the state will help and come to the rescue --- the state will help the private player . This is what happened in the United States, where money was doled out even for people employed privately, and in India as well .
\end{itemize}

\paragraph{How the discontent gets redirected}
\begin{itemize}
    \item When the state is regulating late-capitalist societies and the economy goes through a depression period, people will not criticise the owners of these private companies . People will criticise the government .
    \begin{itemize}
        \item If you lose your job after a lockdown, whom will you abuse --- the owner of your company, or the government?  The owner was sitting at home and may not even know that you were removed --- and yet you abuse the government on social media .
    \end{itemize}
    \item So the discontent generated at the private place will be channelled towards the public sphere .
    \item And when you start abusing the state, the state gets confused . What should we do? Should we be concerned about the welfare rights of these people, or about the declining profits of those people --- because those people are running us, and these people are voting us? The state is trapped .
    \item And this is known as legitimation crisis, in the language of Habermas .
\end{itemize}

\begin{QuickRevision}
\begin{itemize}
    \item Habermas: second-generation critical theory; puzzled that capitalism survived every crisis of the last century .
    \item Answer: state-regulated late capitalism --- the state rescues capital at every bust .
    \item Consequence: discontent generated in the private sphere is channelled into the public sphere; people blame the state, not the employer .
    \item The state, caught between welfare rights and declining profits, faces a legitimation crisis .
\end{itemize}
\end{QuickRevision}

\subsection{Erik Olin Wright --- Contradictory Class Location}
\begin{itemize}
    \item Erik Olin Wright has given the concept of contradictory class location .
    \item He says today we are living in finance capitalism, run by insurance conglomerates, stock markets, banks and mega banks --- these are the agencies which are running capitalism . Money comes from these banks and insurance companies, and we are the shareholders of the big corporations .
\end{itemize}

\paragraph{The contradiction}
\begin{itemize}
    \item You are a proletarian in the company . The moment you go back home, you are a shareholder of that company --- and as a shareholder you are an owner, holding a certain percentage of ownership .
    \item So you are owner and worker both . Are you bourgeoisie or proletariat? You cannot tell . At today's rate these are invalid concepts .
    \item The manager --- if there was ever such a category between bourgeois and proletariat --- is bourgeoisie in comparison to the proletariat, but proletariat in comparison to the bourgeoisie .
    \begin{itemize}
        \item The manager acts like a bourgeois towards the worker: he gives commands and instructions, cuts his salary, writes a bad review . And the same manager acts like a worker in front of the bourgeoisie .
    \end{itemize}
    \item The petty bourgeoisie --- the shop owners, owners of small means of production --- are bourgeoisie when they employ workers in their own shop, and the same is true of everyone acting as a proletarian in another company .
    \item So you would not be able to point at anyone and say ``this is bourgeoisie'' . There is no clear-cut line of distinction today in the era of finance capitalism .
\end{itemize}

\begin{QuickRevision}
\begin{itemize}
    \item Erik Olin Wright: contradictory class location, under finance capitalism run by banks, insurers and stock markets .
    \item The same person is worker by day and shareholder-owner at home .
    \item The manager is bourgeois downwards and proletarian upwards; the petty bourgeois is both at once .
    \item The bourgeois/proletariat line no longer holds cleanly .
\end{itemize}
\end{QuickRevision}

\subsection{Laclau and Mouffe Again --- The Complexity of Identities}
\begin{itemize}
    \item Today we have a highly complex, diverse, fluid, globalised world . There are no clear-cut categories available, and therefore everybody has different concerns --- as a member of a different nation, a different class, a different religion .
    \item And since all of them have different concerns, the struggles would also be different . Therefore it is not only economic class .
    \item There is not just an economic class: we have an environmental class, a sexual class, a religious class, an ethnic class --- different classes . Not classes in the terms of Marx, but classes all the same --- groups of people .
    \item When you have different concerns, how would you mobilise? How would you become a class for itself?  And when this class is also divided across identity --- upper caste, Black, white --- it is not very easy for them to get united under one umbrella for a single cause .
\end{itemize}

\begin{QuickRevision}
\begin{itemize}
    \item A fluid, globalised world offers no clear-cut categories; concerns differ by nation, class and religion .
    \item Environmental, sexual, religious and ethnic ``classes'' exist alongside the economic one .
    \item Cross-cutting identities make mobilisation under one umbrella extremely difficult --- hence no class for itself .
\end{itemize}
\end{QuickRevision}

\subsection{Pierre Bourdieu --- Four Forms of Capital}
\begin{itemize}
    \item Pierre Bourdieu says there are four forms of capital: economic capital, social capital, cultural capital and symbolic capital .
    \item When Marx talks about capital --- volumes 1, 2, 3 --- he is talking only about one form of capital: economic capital . And he believes that only because of this economic capital will we have struggle .
    \item But Bourdieu says we have multiple forms of capital in the world --- and when we have multiple forms of capital, you cannot say that we will have only economic class struggle .
\end{itemize}

\begin{table}[h!]
\centering
\begin{tabularx}{\textwidth}{lYY}
\toprule
\rowcolor{AccentDark}
\thead{Form of capital} & \thead{What it is} & \thead{How it converts} \\
\midrule
Economic & Ownership of the means of production; wealth & The endpoint the other three are converted into \\
Social & Your contacts --- you know politicians and bureaucrats, you sit with them in a club every weekend & With contacts it is not a difficult task to become a millionaire --- social capital can be exchanged for economic capital \\
Cultural & Your education, qualification, the way you speak, language, accent, the way you carry yourself, the clothes you put on & A certain demeanour and good qualification make it easier to establish contacts --- cultural capital converts into social, then into economic \\
Symbolic & Power --- aura, charisma. Some people you cannot speak in front of when you first confront them & An outcome of all the capitals: once you have accumulated the others, you gain symbolic capital \\
\bottomrule
\end{tabularx}
\end{table}

\begin{itemize}
    \item This is known as convertibility of capital --- you have converted your contacts into wealth .
\end{itemize}

\begin{ExampleBox}
\begin{itemize}
    \item If you are the son or daughter of a senior civil service officer, there are very high chances that you will clear the civil services . Not because of any eligibility criterion --- but because you have cultural capital .
    \item You have been trained; you have been seeing those things since childhood; the programmes watched at home, the clubs visited, the way of speaking and writing are all different . These are accumulated advantages gained purely by virtue of being born into that family --- and they can be exchanged for social capital and then for economic capital .
\end{itemize}
\end{ExampleBox}

\begin{itemize}
    \item On symbolic capital: a very elite, sophisticated professor speaking elegantly --- whatever she says seems right and everyone notes it down . There are many millionaires and billionaires, but only a few of them have symbolic capital; it is highly exceptional .
    \item Caution: these are forms of capital, not forms of capitalism . There is no such thing as ``symbolic capitalism'' .
    \item Bourdieu's argument: humans have multiple forms of capital, and each form of capital can be exchanged for any other form . Ultimately anybody can become an owner of the means of production that way . Economic class is not the single group or identity in the world --- class is not only economic .
\end{itemize}

\begin{QuickRevision}
\begin{itemize}
    \item Four forms of capital: economic, social, cultural, symbolic .
    \item Social = contacts; cultural = education, accent, demeanour, taste; symbolic = aura and power, the outcome of all the others .
    \item Convertibility of capital: each form converts into another, ending in economic capital .
    \item The officer's child clears the exam through inherited cultural capital, not eligibility .
    \item Because class is not only economic, class struggle cannot be only economic .
\end{itemize}
\end{QuickRevision}

\subsection{Why Capitalism Sustained Itself --- Gramsci and Digital Capitalism}
If the question is asked --- why did the revolution never occur, and why has capitalism sustained itself --- these are the scholars you can write about .
\begin{itemize}
    \item Laclau and Mouffe --- localised movements, multiple identities, no single class subject .
    \item J\"urgen Habermas --- state-regulated late capitalism and the legitimation crisis .
    \item Adorno and Horkheimer of critical theory --- the cultural industry producing cultural goods which have nullified our consciousness . We have become consumerists; we consume these goods, we like these goods, these are false needs, and we do not realise it . We chase them and are now in a vicious trap . When you are in the habit of watching these things or buying these commodities, you would never think about bringing a revolution --- in fact, if someone discusses revolution with you, you may even start considering him an idiot .
\end{itemize}

\paragraph{Antonio Gramsci and cultural hegemony}
\begin{itemize}
    \item Antonio Gramsci's writing had a very great influence on the development of critical theory . Adorno and Horkheimer are part of critical theory, and Habermas is the extension of that, although his own account of why the revolution did not come is quite different .
    \item Gramsci asks: what is the meaning of success today?  For all of us, the meaning of success is to have wealth, to speak good English, to have a good community, to have great social capital and good contacts .
    \item And how did you come to know that this is what success is? Because this is what you have been studying, watching, listening to and reading --- every single time .
    \item That means education, media, the state and laws --- all these institutions and their dominant ideas --- have made those dominant ideas of material success into common sense to you .
    \begin{itemize}
        \item You see the biggest industrialists photographed with heads of government, and administrative officers photographed with them, so you start thinking that to be successful you must become either an industrialist or an officer . The state advertises this: these are the people who are celebrated; education is what you do to get in there; English is what you learn to be a global, cosmopolitan citizen --- and it is also, today, the meaning of success .
    \end{itemize}
\end{itemize}

\begin{KeyConcept}
Gramsci: when the ideas or the ideologies of the dominant sections of society become common sense to the members, they have become victims of cultural hegemony .
\begin{itemize}
    \item Line to write: cultural hegemony is when the dominant notions and ideas --- of leisure, of entertainment, of success, of good education, of status --- become common sense to the members of society .
\end{itemize}
\end{KeyConcept}

\begin{itemize}
    \item Example: ask someone their hobby and they say they play marbles --- a two-rupee game . And you laugh at that person, and tell him to study, or to become like us and watch the video sites and streaming platforms instead . Basically you are telling him to accept the dominant notion of entertainment, to become like us --- the victims of cultural hegemony --- and never to dare to think beyond it .
    \item When cultural hegemony is established, there is no scope for revolution . Everybody would be involved in speaking good English, making a lot of wealth, having good status and lounging away their time with the great ones . There would be no time and no discourse of revolution . Nobody would find that they are exploited, because they have been brainwashed through media, education and the state as institutions .
\end{itemize}

\paragraph{Ursula Huws --- i-capitalism / digital capitalism}
\begin{itemize}
    \item Capitalism has matured a lot --- we have multiple forms of capitalism now . One of them is i-capitalism, that is, digital capitalism, given by the scholar Ursula Huws .
    \item When we are sitting on the internet updating our status, or watching anything, what is happening is that we are working for those companies . We are giving them our data; we are telling them what our likes are .
    \begin{itemize}
        \item You tell the video platform the kind of videos you like; therefore from the next day it starts giving recommendations of that same type . So you become more and more of a consumer, and you are very happy with the platform .
    \end{itemize}
    \item You think that you are having a time pass, that you are having leisure, that you are having fun . But Ursula Huws is saying that is also work that you are doing --- and you are not able to realise it . You are in false consciousness . You do not think that it is work .
    \item Remember: proletariats give labour power, and labour power gives surplus . Today social media users give their so-called intellectual and leisure time, and that generates surplus . So at the end of the day surplus is accumulated in both ways --- either by working, or by having leisure and intellectual time .
\end{itemize}

\begin{CriticalPoint}
\begin{itemize}
    \item These may be proletariat, but these are cyber-tariats --- who do not even realise that they are working . They think they are having fun .
    \item This is false consciousness at its peak . Earlier people at least knew that they were working; today they do not even know that they are working .
\end{itemize}
\end{CriticalPoint}

\begin{QuickRevision}
\begin{itemize}
    \item Why capitalism survived: Laclau and Mouffe, Habermas, Adorno and Horkheimer's cultural industry, Gramsci, Huws .
    \item Gramsci: cultural hegemony = the dominant sections' ideas of success, leisure and status becoming common sense, transmitted by education, media, state and law .
    \item With hegemony established, there is neither time nor discourse for revolution .
    \item Ursula Huws: i-capitalism / digital capitalism --- leisure online is unpaid work, generating surplus through data .
    \item The cyber-tariat does not even know it is working --- false consciousness at its peak .
\end{itemize}
\end{QuickRevision}

\subsection{Final Doubts}

\paragraph{Wage for housework}
\begin{itemize}
    \item A separate debate --- it would take about 45 minutes to conclude, and is covered under work and economy . The prior question is: what do we call work?  How do we measure and categorise something as work? 
    \begin{itemize}
        \item One economist argues that even filling in application forms for examinations is work --- because when you do the same thing in an office you get paid for it . When you register as an applicant you are telling the state that you want to sit this exam and want to be a member of the bureaucracy, and you are giving your time and your labour . Some people give ten years to the exam --- very precious years, in fact more precious than the years they will later spend working .
    \end{itemize}
\end{itemize}

\paragraph{What does ``dominant discourse'' mean?}
\begin{itemize}
    \item Discourse means the way you speak . The dominant discourse today is equality --- and therefore LGBT persons can speak out, voice their concerns and protest against the state for not letting them have the tag of marriage .
    \item Marriage as an institution has been hijacked, usurped and appropriated by heterosexuals --- as if marriage is the private property of heterosexuals and homosexuals cannot have it .
    \begin{itemize}
        \item You will see Marxist ideologies and discourse everywhere . The way that scholar speaks is very much similar to what Marx would have said about the economic classes and about the oppression of the proletariats .
    \end{itemize}
\end{itemize}

\paragraph{False consciousness --- one line}
\begin{itemize}
    \item False consciousness simply means that we do not realise that we are exploited .
\end{itemize}

\paragraph{How can we convert symbolic capital into economic capital?}
\begin{itemize}
    \item These are forms of capital, not forms of capitalism --- there is no ``symbolic capitalism'' .
    \item Symbolic capital is an outcome of all the capitals . Once you have accumulated cultural capital and economic capital, then you gain symbolic capital . There are many millionaires and billionaires, but only a few billionaires have symbolic capital which the others do not . It cannot be shown objectively --- but symbolic capital is possessed by very few people; it is highly exceptional .
\end{itemize}

\paragraph{Why does the Chinese party keep the communist tag and not the socialist tag?}
\begin{itemize}
    \item China and Russia are both run by communist parties, but they call their states socialist states .
    \item They are not going to call their state a communist state, because ``communist state'' is an oxymoron . The two words cannot be put together --- there cannot be a dark day or a bright night . It should be communist society .
    \item Since they know it is not a communist society, they use the term socialist state . The communist party rules, and capitalism is the form of economy they have been following for a long time --- it is not new to their economy .
\end{itemize}

\paragraph{Why can we not compare Marxian socialism with the socialism of Russia and China?}
\begin{itemize}
    \item In Marxian socialism the state is a temporary body: the state will distribute everything, and once everything is done it will wither away .
    \item In Russia and China the state has become the end . There is no alternative beyond the state socialism they follow .
    \item Hence the term state socialism --- state and socialism always staying together . The state will never wither away .
\end{itemize}

\paragraph{Did Marx know he was a sociologist?}
\begin{itemize}
    \item Marx did not know that he was a sociologist, because at that point of time there was no department of sociology .
    \item And he was never writing like a philosopher, an economist, a political scientist or a sociologist --- because he is everything . He is like a Renaissance man . Look him up and you will see all the labels: political philosopher, political scientist, lawyer, economist, sociologist .
    \begin{itemize}
        \item The same is true of Ambedkar --- you cannot say he was a sociologist, or a historian; he was all of it .
    \end{itemize}
\end{itemize}

\paragraph{Restating Lenin's position}
\begin{itemize}
    \item Nobody can criticise Marx --- that has to be understood first . Lenin is simply saying one thing: you cannot mobilise the classes against the oppressor --- whoever that oppressor may be, a landlord in feudal society or the bourgeoisie in capitalism --- unless you have a political party .
    \begin{enumerate}
        \item The political party will impart consciousness among those people who are exploited .
        \item The political party will tell the people how they are exploited .
        \item The political party will publish journals, newspapers and editorials to educate them on how they are exploited .
        \item And when they are finally convinced that they are exploited, these convinced people will be mobilised against the oppressors . That is what Lenin did --- he formed the Vanguard Party .
    \end{enumerate}
    \item Marx never said to form a political party . Marx said class in itself will automatically become class for itself with homogenisation, monopolisation and depression [as stated in lecture; earlier delivered as pauperisation] . Was there any fourth condition of a political party? No .
    \item Lenin says all of that will not be enough --- ultimately you will have to build a political party to impart consciousness among these people . That is why the philosophy is called Marxist--Leninist: ``Leninist'' means the agenda of the political party has now been added, which Marx never proposed .
    \item In fact Marx did not believe in democratic instruments like a political party . He believed democracy is a bourgeois ideology --- false consciousness, telling us equality, equality, so that revolution never comes .
\end{itemize}

\begin{QuickRevision}
\begin{itemize}
    \item Wage for housework turns on the prior question of what counts as work at all .
    \item Dominant discourse = equality; hence the claim that marriage has been appropriated as the private property of heterosexuals .
    \item False consciousness = not realising that we are exploited .
    \item Symbolic capital is the outcome of all other capitals and is possessed by very few .
    \item ``Communist state'' is an oxymoron; hence socialist state + communist party + capitalist economy .
    \item Marxian socialism: the state is temporary . Russia and China: state socialism, where the state is the end .
    \item Marx was a Renaissance man, not a sociologist by self-description .
    \item Lenin adds the political party as the vehicle of consciousness; Marx rejected the party, holding democracy to be a bourgeois ideology .
\end{itemize}
\end{QuickRevision}

\sectiondivider

\section{Answer-Writing Toolkit}
Every item below was stated by the teacher across the five classes as guidance for writing answers . Nothing here has been added from outside the lectures .

\subsection{The Standard of Understanding an Answer Must Reflect}
\begin{itemize}
    \item Can you explain the concept in simple language in two minutes to someone with no background --- a parent, or a rickshaw puller? If yes, your understanding is at the required level .
    \item Explain the essence and the meaning of the term --- not the features .
    \item Once understanding is secure, revision is only for keywords and facts, and takes very little time .
    \item Whatever you struggle to articulate, scribble down --- that becomes your own mind map, and you refer to it before handouts or class .
\end{itemize}

\subsection{Examples are Compulsory}
\begin{itemize}
    \item Whenever you study anything theoretical, always look out for examples .
    \item Unless you have examples in your head to justify what you have written in theory, the answer will not reflect the level of understanding expected from a candidate .
    \item The examiner can only know you have understanding when you write examples . Without them you are not showing your true potential .
\end{itemize}

\subsection{Managing Time in the Paper}
\begin{itemize}
    \item Time is lost in the thinking process --- typically 20--30 minutes across the paper, in fragments .
    \item The cause is either lack of revision or, more often, lack of understanding of what you are revising .
    \item Leaving one or two questions is acceptable at the start; leaving four points to a problem with writing speed or with the thinking process .
    \item Questions will always be new --- the examination does not repeat questions . Memorisation is not what is being tested .
\end{itemize}

\subsection{Handling the Previous-Year Question Flagged in Class}
\begin{PYQLink}
\begin{itemize}
    \item PYQ: ``Is non-positivist methodology scientific?'' --- deferred to Unit 4 in class, but unanswerable unless you know what ``scientific'' means .
    \item If you know non-positivism but assume positivism = science, you will wrongly conclude that non-positivism is not science .
    \item The question is very basic --- and basic questions are the toughest .
\end{itemize}
\end{PYQLink}

\paragraph{The method demonstrated on ``higher unemployment is the source of higher crime rate''}
\begin{enumerate}
    \item Identify the type of statement: objective, cause--effect, testable $\rightarrow$ a positivist proposition . A phenomenologist would never make it: for phenomenology there are multiple causes behind multiple effects and objectivity is an illusion .
    \item Note that Marxism would also reject it, because to Marx the owning, ruling class is the criminal class --- so defining crime as working-class crime is not Marxist .
    \item Establish that scientific = testable = falsifiable, not ``true'' .
    \item Think about possibilities beyond the given statement . You do not have to quote data or reports . If A $\rightarrow$ B is falsifiable, what else could lead to B --- C, D, E? 
\end{enumerate}

\subsection{Lines the Teacher Dictated to be Written Down Verbatim}
\begin{itemize}
    \item The theory of historical materialism reflects how societies develop class characters, and how the relations of production between these classes are justified by the forms of consciousness --- that is, the superstructure .
    \item All modes of production developed as a result of inherent contradictions in the previous modes of production, which were the result of the contradiction between relations of production and means of production .
    \item Class is a group of people having similar relations to the means of production .
    \item The value of commodity = $C + V + S$ . Rate of exploitation = $S / V = \text{unpaid labour} / \text{paid labour}$ .
    \item Constant capital does not add any value to the commodity; it is the variable capital that adds all the value .
    \item Absolute surplus value is generated by increasing the length of the working day, and by tight monitoring or surveillance . However, the length of the working day cannot be increased beyond a point, due to labour laws and human rights .
    \item Relative surplus value is achieved by increasing the productivity of the labour and the introduction of machines .
    \item The industrial reserve army consists of a pool of both the unemployed who want to work, and the underemployed .
    \item Commodity fetishism is the result of actors failing to realise that the commodity is the creation of labour .
    \item The law of motion reveals the inherent contradictions within capitalism, and describes the mechanism of the inherent contradiction .
    \item The members of class in itself are the victims of false consciousness . False consciousness develops when the degree of exploitation is not realised by the worker .
    \item The only reason why class in itself has not yet become class for itself is the prevalence of false consciousness perpetuated by the superstructure .
    \item Alienation is the social breakdown of the natural interconnectedness between the worker and his produce, and between the worker and the production relations .
    \item Marx's analysis of capitalism and of alienation was limited to the manual working class --- blue-collar workers . C. Wright Mills discussed alienation among non-manual, white-collar workers .
    \item Instead of categorising the working class as victims of alienation, the job of a social scientist should be to see what meaning they attach to their work (Goldthorpe and Lockwood) .
    \item Cultural hegemony is when the dominant notions and ideas --- of leisure, entertainment, success, education and status --- become common sense to the members of society (Gramsci) .
\end{itemize}

\subsection{Terminology Cautions}
\begin{itemize}
    \item Use infrastructure, not ``economic base'', when pairing with superstructure --- ``economic base'' is general language, right but technically not sound .
    \item Do not confuse means of production with mode of production . Write MOP in capitals for Mode of Production and means of production in small letters .
    \item Forces of production = means of production --- the difference is only one of translation; use them interchangeably .
    \item Do not say ``pure capitalism'' --- the term is pure communism .
    \item Never list the Asiatic mode of production with the four main modes --- thin literature, and Marx never wrote extensively on it .
    \item Bourgeoisie is a noun; bourgeois is the adjective .
    \item Four forms of capital, not forms of capitalism --- there is no ``symbolic capitalism'' (Bourdieu) .
    \item ``Communist state'' is an oxymoron --- the correct pairing is communist society, or socialist state ruled by a communist party .
    \item Pauper, not Popper --- pauperisation comes from pauper, meaning poor .
\end{itemize}

\subsection{Strategy Notes}
\begin{itemize}
    \item The syllabus sits at BA honours level --- between BA and MA . Some recent questions come from MA-level areas .
    \item Do not chase books for the one MA-level question out of twenty; you may lose the larger picture of the nineteen you are supposed to do . There is choice in the paper .
    \item No scholar should be your favourite . You cannot afford to be in love with one scholar .
    \item There are hundreds of critics of every thinker; limits have to be set on how many scholars you carry .
    \item Note-making method is your own choice --- device or notebook . Devices allow editing and additions later; there is no single recommended method .
\end{itemize}

\subsection{Ready-Made Answer Frames from These Classes}
\begin{itemize}
    \item ``Why has the revolution not occurred / why has capitalism sustained itself?'' $\rightarrow$ false consciousness perpetuated by the superstructure; Laclau and Mouffe on localised movements and multiple identities; Habermas on state-regulated late capitalism and the legitimation crisis; Adorno and Horkheimer on the cultural industry; Gramsci on cultural hegemony; Ursula Huws on digital capitalism and the cyber-tariat .
    \item ``Is Marxism a science?'' $\rightarrow$ the Popper problem stated in class: happy = false consciousness, unhappy = alienation; if it cannot be proven wrong, it is not science .
    \item ``Criticism of alienation'' $\rightarrow$ Weber (rationalisation and impersonality); Blauner (relative, operationalised in four dimensions, inverted U-curve); Goldthorpe and Lockwood (meanings attached to work); C. Wright Mills (white-collar); Gorz and Marcuse (leisure); de Beauvoir and Hochschild (women, emotional labour) . Note that most of these are applications, not criticisms .
    \item ``Criticism of class in itself / class for itself'' $\rightarrow$ Lenin (Vanguard Party as a fourth prerequisite); Erik Olin Wright (contradictory class location); Bourdieu (four forms of capital); Laclau and Mouffe (fragmented identities) .
    \item ``Criticism of the mode of production and scientific socialism'' $\rightarrow$ Giddens (no universal formula, other struggles); Huntington (fault line conflict); the empirical jump from feudalism to socialism in Russia and China; Lynette Ong on China; Galbraith; Fukuyama .
    \item History vs sociology question $\rightarrow$ Marxism was a response to the Industrial Revolution; sociology is rootless without history, but history is fruitless without sociology . Origins of patriarchy, private property and class inequality are all historical sociology .
\end{itemize}